%%%%%%%%%%%%%%%%%%%%%%%%%%%%%%%%%%%%%%%%%%%%%%%%%%%%%%%%%%%%%%%%%%%%%%%%%%%%%
%%  main.tex — پایان‌نامه کارشناسی ملیکا مقدم‌کیا
%%  دانشگاه آزاد اسلامی واحد تهران مرکزی — زمستان ۱۴۰۴
%%  موتور: XeLaTeX  |  بسته اصلی: xepersian
%%%%%%%%%%%%%%%%%%%%%%%%%%%%%%%%%%%%%%%%%%%%%%%%%%%%%%%%%%%%%%%%%%%%%%%%%%%%%

%% ─── کلاس سند ──────────────────────────────────────────────────────────────
\documentclass[a4paper,14pt]{extreport}

%% ─── بسته‌های پایه ──────────────────────────────────────────────────────────
\usepackage{fontspec}          % مدیریت فونت با XeLaTeX
\usepackage{amsmath,amssymb}   % ریاضیات
\usepackage{graphicx}          % تصاویر
\usepackage{booktabs}          % جداول حرفه‌ای
\usepackage{longtable}         % جداول چند صفحه‌ای
\usepackage{multirow}          % ادغام سطرها در جدول
\usepackage{array}             % ستون‌بندی پیشرفته جدول
\usepackage{tabularx}          % جداول با عرض ثابت
\usepackage{geometry}          % تنظیم حاشیه‌ها
\usepackage{setspace}          % فاصله خطوط
\usepackage{hyperref}          % لینک‌ها و مراجع کلیک‌پذیر
\usepackage{caption}           % کنترل کپشن
\usepackage{subcaption}        % زیر کپشن
\usepackage{float}             % کنترل موقعیت شکل/جدول
\usepackage{fancyhdr}          % سرصفحه و پاصفحه
\usepackage{titlesec}          % قالب‌بندی عناوین
\usepackage{tocloft}           % قالب‌بندی فهرست‌ها
\usepackage{enumitem}          % کنترل لیست‌ها
\usepackage{xcolor}            % رنگ‌بندی
\usepackage{microtype}         % بهینه‌سازی حروف‌چینی
\usepackage{pdfpages}          % درج PDF خارجی
\usepackage{afterpage}         % دستورات بعد از صفحه
\usepackage{emptypage}         % صفحات خالی بدون سرصفحه

%% ─── تنظیم حاشیه‌ها (A4، استاندارد IAU) ───────────────────────────────────
\geometry{
  a4paper,
  top=3cm,
  bottom=2.5cm,
  left=2.5cm,     % حاشیه راست در چاپ دوطرفه (صحافی)
  right=3cm,      % حاشیه چپ
  headheight=1.5cm,
  headsep=0.5cm
}

%% ─── فاصله خطوط ─────────────────────────────────────────────────────────────
\linespread{1.5}

%% ─── تنظیمات hyperref ──────────────────────────────────────────────────────
\hypersetup{
  colorlinks=true,
  linkcolor=black,
  citecolor=black,
  urlcolor=blue,
  bookmarksnumbered=true,
  bookmarksopen=true,
  pdftitle={طراحی معماری امنیتی ترکیبی پساکوانتومی و توزیع کلید کوانتومی در ارتباطات ماهواره‌های مدار پایین زمین},
  pdfauthor={ملیکا مقدم‌کیا}
}

%% ─── xepersian باید آخرین بسته بارگذاری شود ────────────────────────────────
\usepackage{xepersian}

%% ─── تنظیم فونت‌ها ──────────────────────────────────────────────────────────
\settextfont{B Lotus_0.ttf}
\setboldtextfont{B Lotus Bold_0.ttf}
\setlatintextfont{times.ttf}
\setdigitfont{B Lotus_0.ttf}

%% ─── تنظیم عدد‌نویسی فارسی در شماره صفحات و شکل‌ها ────────────────────────
\renewcommand{\thefigure}{\thechapter-\arabic{figure}}
\renewcommand{\thetable}{\thechapter-\arabic{table}}
\renewcommand{\theequation}{\thechapter-\arabic{equation}}

%% ─── قالب‌بندی فصل‌ها و بخش‌ها ──────────────────────────────────────────────
\titleformat{\chapter}[display]
  {\bfseries\Large\centering}
  {فصل \thechapter}
  {0.5em}
  {\bfseries\Large}
\titlespacing*{\chapter}{0pt}{20pt}{30pt}

\titleformat{\section}
  {\bfseries\large\raggedleft}
  {\thesection}
  {0.7em}
  {}
\titlespacing*{\section}{0pt}{12pt}{6pt}

\titleformat{\subsection}
  {\bfseries\normalsize\raggedleft}
  {\thesubsection}
  {0.7em}
  {}
\titlespacing*{\subsection}{0pt}{10pt}{5pt}

%% ─── کپشن فارسی ─────────────────────────────────────────────────────────────
\captionsetup[figure]{
  name={شکل},
  labelsep=period,
  position=below,
  font=small,
  labelfont=bf
}
\captionsetup[table]{
  name={جدول},
  labelsep=period,
  position=above,
  font=small,
  labelfont=bf
}

%% ─── سرصفحه و پاصفحه ───────────────────────────────────────────────────────
\pagestyle{fancy}
\fancyhf{}
\fancyhead[R]{\small\rightmark}
\fancyhead[L]{\small\leftmark}
\fancyfoot[C]{\thepage}
\renewcommand{\headrulewidth}{0.4pt}

%% ─── تعریف شماره‌گذاری ابجد برای صفحات مقدماتی ────────────────────────────
\newcommand{\abjadpage}[1]{%
  \setcounter{page}{#1}%
}

%% ─── دستور \lr{} برای متن لاتین در محیط فارسی ──────────────────────────────
%% (xepersian این دستور را به‌طور خودکار فراهم می‌کند)

%%%%%%%%%%%%%%%%%%%%%%%%%%%%%%%%%%%%%%%%%%%%%%%%%%%%%%%%%%%%%%%%%%%%%%%%%%%%%
%%  آغاز سند
%%%%%%%%%%%%%%%%%%%%%%%%%%%%%%%%%%%%%%%%%%%%%%%%%%%%%%%%%%%%%%%%%%%%%%%%%%%%%
\begin{document}

%%%%%%%%%%%%%%%%%%%%%%%%%%%%%%%%%%%%%%%%%%%%%%%%%%%%%%%%%%%%%%%%%%%%%%%%%%%%%
%%  صفحات مقدماتی — شماره‌گذاری ابجد (الف، ب، ج، ...)
%%%%%%%%%%%%%%%%%%%%%%%%%%%%%%%%%%%%%%%%%%%%%%%%%%%%%%%%%%%%%%%%%%%%%%%%%%%%%
\pagenumbering{harfi}   % شماره‌گذاری فارسی ابجدی توسط xepersian
\pagestyle{empty}       % بدون سرصفحه در صفحات مقدماتی

%% ─────────────────────────────────────────────────────────────────────────
%% ۱. طرح روی جلد فارسی
%% ─────────────────────────────────────────────────────────────────────────

%% نشان دانشگاه — در صورت وجود فایل logo.png آن را فعال کنید
%\includegraphics[width=3cm]{logo.png}\\[1cm]

%% ─────────────────────────────────────────────────────────────────────────
%% صفحه الف — بسم الله الرحمن الرحیم
%% ─────────────────────────────────────────────────────────────────────────
\newpage
\thispagestyle{empty}
\vspace*{\fill}
\begin{center}
{\Huge\bfseries بِسْمِ اللَّهِ الرَّحْمَنِ الرَّحِیمِ}
\end{center}
\vspace*{\fill}
\newpage

%% ─────────────────────────────────────────────────────────────────────────
%% صفحه ب — منشور اخلاق پژوهش
%% ────────────────────────────────
%% ─────────────────────────────────────────────────────────────────────────
%% صفحه ج — تعهدنامه اصالت پایان‌نامه
%% ────────────────────────────────

%% ─────────────────────────────────────────────────────────────────────────
%% صفحه ه — سپاسگزاری
%% ─────────────────────────────────────────────────────────────────────────
\newpage
\thispagestyle{empty}
\begin{center}
{\Large\bfseries سپاسگزاری}
\end{center}
\vspace{1.5cm}

{\normalsize
سپاس بی‌کران از درگاه خداوند متعال که توفیق بهره‌مندی از دانش را عنایت فرمود.

با کمال احترام و ادب، مراتب سپاسگزاری خویش را از استاد گرانقدر، جناب آقای \textbf{سیدرضا وزیری یزدی}، که با راهنمایی‌های ارزشمند، صبر و دقت علمی خود در تمامی مراحل این پژوهش همراهم بودند، ابراز می‌دارم.

همچنین از تمامی اساتید گروه مهندسی کامپیوتر دانشکده فنی و مهندسی دانشگاه آزاد اسلامی واحد تهران مرکزی که در طول دوران تحصیل بنیان‌های علمی‌ام را استوار ساختند، صمیمانه سپاسگزارم.

از خانواده عزیزم که در تمامی لحظه‌های این مسیر، پشتوانه معنوی و عاطفی‌ام بودند، قدردانی می‌کنم.
}

\vspace{3cm}
\begin{flushright}
{\normalsize ملیکا مقدم‌کیا}\\
{\normalsize زمستان ۱۴۰۴}
\end{flushright}
\newpage

%% ─────────────────────────────────────────────────────────────────────────
%% صفحه و — تقدیم
%% ─────────────────────────────────────────────────────────────────────────
\newpage
\thispagestyle{empty}
\vspace*{3cm}
\begin{center}
{\Large\bfseries تقدیم به}\\[2cm]

{\large
پدر و مادر عزیزم\\[0.5cm]
که نور راهنمای هر قدمم بوده‌اند\\[1.5cm]
و به تمامی کسانی\\[0.3cm]
که دانش را در خدمت صلح و امنیت بشر می‌دانند
}
\end{center}
\newpage

%% ─────────────────────────────────────────────────────────────────────────
%% صفحه ز — فرم نمره
%% ────────────────────────────────────────────────────────────────
%% ─────────────────────────────────────────────────────────────────────────
%% صفحه ح — فرم ب (اطلاعات پایان‌نامه و چکیده فارسی)
%% ─────────────────────────────────────────────────────────────────────────
\newpage
\thispagestyle{empty}
\begin{center}
{\Large\bfseries دانشگاه آزاد اسلامی  واحد تهران مرکزی}\\[0.3cm]
{\normalsize دانشکده فنی و مهندسی  گروه مهندسی کامپیوتر}\\[0.8cm]
{\Large\bfseries فرم ب: چکیده پایان‌نامه کارشناسی}
\end{center}
\vspace{0.5cm}

\noindent
\begin{tabular}{rp{9cm}}
\bfseries عنوان فارسی: & طراحی معماری امنیتی ترکیبی پساکوانتومی \lr{(PQC)} و توزیع کلید کوانتومی \lr{(QKD)} در ارتباطات ماهواره‌های مدار پایین زمین \\[0.3cm]
\bfseries عنوان انگلیسی: & \lr{Design of a Hybrid Post-Quantum Cryptography (PQC) and Quantum Key Distribution (QKD) Security Architecture for Low Earth Orbit (LEO) Satellite Communications} \\[0.3cm]
\bfseries نگارنده: & ملیکا مقدم‌کیا \\[0.2cm]
\bfseries استاد راهنما: & سیدرضا وزیری یزدی \\[0.2cm]
\bfseries مقطع تحصیلی: & کارشناسی مهندسی کامپیوتر \\[0.2cm]
\bfseries تاریخ دفاع: & زمستان ۱۴۰۴ \\
\end{tabular}



%% ─────────────────────────────────────────────────────────────────────────
%% فهرست مطالب — صفحه ط
%% ─────────────────────────────────────────────────────────────────────────
\newpage
\pagestyle{fancy}
\renewcommand{\contentsname}{\begin{center}{\Large\bfseries فهرست مطالب}\end{center}}
\tableofcontents
\newpage

%% ─────────────────────────────────────────────────────────────────────────
%% فهرست شکل‌ها — صفحه ی
%% ─────────────────────────────────────────────────────────────────────────
\newpage
\renewcommand{\listfigurename}{\begin{center}{\Large\bfseries فهرست شکل‌ها}\end{center}}
\listoffigures
\newpage

%% ─────────────────────────────────────────────────────────────────────────
%% فهرست جدول‌ها — صفحه ک
%% ─────────────────────────────────────────────────────────────────────────
\newpage
\renewcommand{\listtablename}{\begin{center}{\Large\bfseries فهرست جدول‌ها}\end{center}}
\listoftables
\newpage

%%%%%%%%%%%%%%%%%%%%%%%%%%%%%%%%%%%%%%%%%%%%%%%%%%%%%%%%%%%%%%%%%%%%%%%%%%%%%
%%  آغاز شماره‌گذاری عددی — از اینجا صفحه ۱ شروع می‌شود
%%%%%%%%%%%%%%%%%%%%%%%%%%%%%%%%%%%%%%%%%%%%%%%%%%%%%%%%%%%%%%%%%%%%%%%%%%%%%
\newpage
\pagenumbering{arabic}
\setcounter{page}{1}
\pagestyle{fancy}
\thispagestyle{plain}

\begin{center}
{\Large\bfseries چکیده}
\end{center}
\vspace{0.5cm}
%%%%%%%%%%%%%%%%%%%%%%%%%%%%%%%%%%%%%%%%%%%%%%%%%%%%%%%%%%%%%%%%%%%%%%%%%%%%%
%%  abstract_fa.tex — چکیده فارسی
%%  محل قرارگیری: در main.tex، بلافاصله پیش از \chapter{کلیات تحقیق}
%%  (یعنی پس از فهرست جدول‌ها و قبل از \pagenumbering{arabic})
%%%%%%%%%%%%%%%%%%%%%%%%%%%%%%%%%%%%%%%%%%%%%%%%%%%%%%%%%%%%%%%%%%%%%%%%%%%%%

%% ─── چکیده فارسی ────────────────────────────────────────────────────────────
%% طبق فرمت IAU تهران مرکزی، چکیده فارسی اولین صفحه‌ای است که
%% شماره‌گذاری عددی (1) دارد. بنابراین \pagenumbering{arabic} را
%% دقیقاً قبل از این محیط قرار دهید.
\newpage
\pagenumbering{arabic}
\setcounter{page}{1}
\pagestyle{fancy}
\thispagestyle{plain}

\begin{center}
{\Large\bfseries چکیده}
\end{center}

\vspace{0.5cm}

{\normalsize
گسترش سریع ارتباطات ماهواره‌ای در مدار پایین زمین \lr{(Low Earth Orbit --- LEO)} و پیشرفت متزامن محاسبات کوانتومی \lr{(Quantum Computing)}، تهدیدی بنیادین برای زیرساخت‌های رمزنگاری کلاسیک مبتنی بر \lr{RSA} و رمزنگاری منحنی بیضوی \lr{(ECC)} پدید آورده است. الگوریتم شور \lr{(Shor's Algorithm)} قادر است این طرح‌ها را بر روی یک کامپیوتر کوانتومی مقیاس‌کافی در زمان چندجمله‌ای بشکند؛ این واقعیت، همراه با تهدید «ضبط‌کن اکنون، رمزگشا بعداً» \lr{(Harvest Now, Decrypt Later --- HNDL)}، امنیت بلندمدت داده‌های ماهواره‌ای را با عمر عملیاتی ۱۰ تا ۱۵ ساله به‌شدت در معرض خطر قرار می‌دهد. پژوهش حاضر در پاسخ به این شکاف امنیتی، یک معماری امنیتی ترکیبی سه‌لایه برای ارتباطات ماهواره‌های \lr{LEO} طراحی و ارزیابی می‌کند.

لایه اول معماری پیشنهادی، پروتکل توزیع کلید کوانتومی \lr{(Quantum Key Distribution --- QKD)} مبتنی بر \lr{BB84} با روش‌های فریب‌حالت \lr{(Decoy-State BB84)} در طول موج ۸۵۰ نانومتر را برای تأمین محرمانگی مستحکم با پشتوانه اطلاعات‌تئوریک \lr{(Information-Theoretic Security)} به کار می‌گیرد. لایه دوم، استاندارد \lr{ML-KEM-768 (FIPS 203)} را برای تبادل کلید مقاوم در برابر تهدیدات کوانتومی بر بستر کانال کلاسیک مستقر می‌سازد. لایه سوم، الگوریتم \lr{ML-DSA-65 (FIPS 204)} را برای احراز هویت دیجیتال و امضای تلگرام‌های کنترلی به کار می‌گیرد. مکانیزم \lr{Fail-Safe} بر پایه یک ماشین حالت محدود \lr{(Finite State Machine --- FSM)} دو لایه اول را در یک ساختار انعطاف‌پذیر هماهنگ می‌کند تا در صورت افت کیفیت کانال \lr{QKD}، سیستم به‌طور خودکار به حالت \lr{PQC-only} بازگردد.

ارزیابی عملکرد با کتابخانه \lr{liboqs-python} و مقایسه با مبناهای کلاسیک نشان داد که \lr{ML-KEM-768} در تولید کلید بیش از ۳۰۰۰ برابر از \lr{RSA-3072} سریع‌تر است و فرضیه ناتوانی محاسباتی الگوریتم‌های پساکوانتومی مبتنی بر شبکه را به‌طور قطعی رد می‌کند. تأخیر کل برقراری نشست امن در معماری پیشنهادی حدود ۳۵ میلی‌ثانیه است که در مقایسه با پنجره دید ۳۰۰ ثانیه‌ای ماهواره \lr{LEO} در ارتفاع ۵۰۰ کیلومتر کاملاً ناچیز است. مدل تحلیلی بودجه لینک \lr{(Link Budget)} کانال \lr{QKD} نشان داد که نرخ خطای بیت کوانتومی \lr{(Quantum Bit Error Rate --- QBER)} در زاویه فراز بهینه در محدوده ۱ تا ۳ درصد باقی می‌ماند و آستانه امنیتی ۱۱ درصد پروتکل \lr{BB84} را نقض نمی‌کند. نتایج کلی نشان می‌دهند که معماری پیشنهادی یک دژ دفاعی مصالحه‌آگاه در برابر هر دو نسل کامپیوترهای کلاسیک و کوانتومی فراهم می‌آورد.
}

\vspace{0.8cm}
\noindent\textbf{کلیدواژه‌ها:} ارتباطات مدار پایین زمین، رمزنگاری پساکوانتومی، توزیع کلید کوانتومی، نرخ خطای بیت کوانتومی، بودجه لینک

%%%%%%%%%%%%%%%%%%%%%%%%%%%%%%%%%%%%%%%%%%%%%%%%%%%%%%%%%%%%%%%%%%%%%%%%%%%%%
%%  فصل اول: کلیات تحقیق
%%%%%%%%%%%%%%%%%%%%%%%%%%%%%%%%%%%%%%%%%%%%%%%%%%%%%%%%%%%%%%%%%%%%%%%%%%%%%
\chapter{کلیات تحقیق}
\setcounter{figure}{0}
\setcounter{table}{0}
\setcounter{equation}{0}

%% ──────────────────────────────────────────────────────────────────────────
\section{مقدمه}
%% ──────────────────────────────────────────────────────────────────────────

در دنیای امروز، ارتباطات ماهواره‌ای \lr{(Satellite Communications)} به یکی از ستون‌های اساسی زیرساخت اطلاعاتی جهان تبدیل شده‌اند. از انتقال داده‌های سازمان‌های دولتی و نظامی گرفته تا پشتیبانی از شبکه‌های تلفن همراه نسل پنجم \lr{(5G)}، ناوبری جهانی \lr{(GPS)} و اینترنت ماهواره‌ای، این سیستم‌ها هر روز حجم عظیمی از اطلاعات حساس را بین نقاط مختلف کره زمین جابه‌جا می‌کنند. امنیت این کانال‌های ارتباطی، شرط بقای بسیاری از خدمات حیاتی است. تاکنون، پایه امنیت این ارتباطات بر رمزنگاری کلید عمومی \lr{(Public Key Cryptography)} متکی بوده که شامل الگوریتم‌هایی مانند \lr{RSA} و رمزنگاری منحنی بیضوی \lr{(Elliptic Curve Cryptography --- ECC)} می‌شود. این الگوریتم‌ها امنیت خود را از سختی مسائل ریاضی خاصی مانند تجزیه اعداد بزرگ به عوامل اول  می‌گیرند \{1\}.

اما در دهه اخیر، پیشرفت شتابان در زمینه محاسبات کوانتومی \lr{(Quantum Computing)} این بنیان را به چالش کشیده است. در سال ۱۹۹۴، پیتر شور \lr{(Peter Shor)} ثابت کرد که یک کامپیوتر کوانتومی به اندازه کافی بزرگ می‌تواند با اجرای الگوریتم شور \lr{(Shor's Algorithm)}، مسئله تجزیه اعداد را در زمانی چندجمله‌ای حل کند یعنی هر دو الگوریتم \lr{RSA} و \lr{ECC} عملاً شکسته خواهند شد \{1\}. این تهدید صرفاً نظری نیست. پژوهشگران و آژانس‌های اطلاعاتی نسبت به یک سناریوی خطرناک هشدار داده‌اند که تحت عنوان «اکنون ضبط کن، بعداً رمزگشایی کن» \lr{(Harvest Now, Decrypt Later --- HNDL)} شناخته می‌شود: مهاجمان از همین امروز ارتباطات رمزشده را ضبط می‌کنند تا در آینده، پس از دسترسی به کامپیوترهای کوانتومی قوی، آن‌ها را بازگشایی کنند \{1\}. این مسئله برای ماهواره‌هایی که چرخه عمر عملیاتی ۱۰ تا ۱۵ ساله دارند، بسیار جدی‌تر است؛ چرا که داده‌های رد و بدل‌شده امروز ممکن است در آینده‌ای نه چندان دور در معرض خطر قرار گیرند \{22\}.

در پاسخ به این تهدید، جامعه علمی دو رویکرد مکمل را توسعه داده است. رویکرد اول، رمزنگاری پساکوانتومی \lr{(Post-Quantum Cryptography --- PQC)} است؛ دسته‌ای از الگوریتم‌های کلاسیک که امنیت آن‌ها بر پایه مسائل ریاضی‌ای استوار است که حتی کامپیوترهای کوانتومی نیز در حل آن‌ها ناتوان هستند. مؤسسه ملی استاندارد و فناوری آمریکا \lr{(National Institute of Standards and Technology --- NIST)} پس از یک فرآیند ارزیابی هشت‌ساله، در سال ۲۰۲۴ استانداردهای \lr{FIPS 203} و \lr{FIPS 204} را تحت عنوان \lr{ML-KEM} و \lr{ML-DSA} منتشر کرد که هر دو بر پایه ریاضیات شبکه مشبک \lr{(Lattice-based Cryptography)} بنا شده‌اند \{3, 9\}. رویکرد دوم، توزیع کلید کوانتومی \lr{(Quantum Key Distribution --- QKD)} است که با بهره‌گیری از اصول بنیادی فیزیک کوانتومی، کانالی برای تبادل کلید فراهم می‌کند که هرگونه استراق‌سمع را آشکار می‌سازد. پروتکل \lr{BB84} به‌عنوان پایه‌ای‌ترین پروتکل \lr{QKD}، امنیت اطلاعاتی \lr{(Information-Theoretic Security)} را تضمین می‌کند  امنیتی که حتی با در اختیار داشتن توان محاسباتی نامحدود نیز قابل شکستن نیست \{4\}.

با وجود این، هر یک از دو رویکرد به‌تنهایی محدودیت‌هایی دارد. \lr{QKD} ماهواره‌ای با چالش‌هایی مانند نرخ کلید پایین در شرایط جوی نامناسب، پنجره‌های دید کوتاه (به‌ویژه در مدار \lr{LEO}) و هزینه بالای پیاده‌سازی روبرو است \{23\}؛ در حالی که \lr{PQC} هنوز یک رویکرد نوپا با چالش‌های عملی در محیط‌های محاسباتی محدود ماهواره‌ای محسوب می‌شود \{15\}. از این رو، ترکیب هوشمندانه این دو رویکرد در یک معماری یکپارچه که بتواند نقاط ضعف هر یک را با نقاط قوت دیگری جبران کند، به یکی از مهم‌ترین جهت‌های پژوهشی در حوزه امنیت ارتباطات ماهواره‌ای تبدیل شده است \{19\}.

پایان‌نامه حاضر در همین راستا تدوین شده است. هدف اصلی، بررسی تهدیدات کوانتومی علیه زیرساخت‌های ماهواره‌ای، تحلیل معماری‌های امنیتی موجود و پیشنهاد یک معماری ترکیبی \lr{PQC+QKD} متناسب با محدودیت‌های واقعی ارتباطات ماهواره‌ای در مدار پایین زمین \lr{(Low Earth Orbit --- LEO)} است. در فصل‌های پیش‌رو، این مسئله به‌صورت گام‌به‌گام از مبانی نظری تا تحلیل عملکرد بررسی خواهد شد.

%% ──────────────────────────────────────────────────────────────────────────
\section{بیان مسئله}
%% ──────────────────────────────────────────────────────────────────────────

امنیت ارتباطات ماهواره‌ای در دنیای امروز بر پایه دو ستون اصلی استوار است: رمزنگاری کلید عمومی برای تبادل امن کلید، و زیرساخت گواهی دیجیتال برای احراز هویت. این دو ستون در حال حاضر و پیش از ظهور کامپیوترهای کوانتومی قدرتمند، کارایی خود را ثابت کرده‌اند. اما پیشرفت‌های اخیر در زمینه محاسبات کوانتومی نشان می‌دهد که این پایه‌های امنیتی به‌زودی دچار شکست اساسی خواهند شد. در ادامه، دو بُعد اصلی این مسئله تشریح می‌شود.

%% ────────────────────────────────────
\subsection{آسیب‌پذیری زیرساخت کلید عمومی \lr{(PKI)}}
%% ────────────────────────────────────

زیرساخت کلید عمومی \lr{(Public Key Infrastructure --- PKI)} سیستمی است که با استفاده از گواهی‌های دیجیتال، هویت طرفین ارتباط را تأیید می‌کند و تبادل امن کلیدهای رمزنگاری را ممکن می‌سازد. قلب این سیستم، الگوریتم‌های رمزنگاری کلید عمومی مانند \lr{RSA} و رمزنگاری منحنی بیضوی \lr{(Elliptic Curve Cryptography --- ECC)} هستند که هر دو امنیت خود را از سختی برخی مسائل ریاضی می‌گیرند: \lr{RSA} به سختی تجزیه اعداد بزرگ به عوامل اول تکیه می‌کند و \lr{ECC} بر اساس دشواری محاسبه لگاریتم گسسته روی منحنی‌های بیضوی بنا شده است \{1\}.

مشکل اینجاست که هر دو این مسائل ریاضی توسط کامپیوتر کوانتومی قابل حل هستند. در سال ۱۹۹۴، پیتر شور ثابت کرد که الگوریتم شور \lr{(Shor's Algorithm)} می‌تواند هر دوی این مسائل را در زمانی چندجمله‌ای  یعنی بسیار سریع‌تر از هر الگوریتم کلاسیک موجود  حل کند \{1\}. به‌عبارت ساده‌تر، یک کامپیوتر کوانتومی به اندازه کافی قوی می‌تواند در چند ساعت کاری را انجام دهد که یک ابرکامپیوتر کلاسیک به میلیون‌ها سال نیاز دارد. بر اساس نتایج مستند در ادبیات علمی، \lr{RSA-3072} و \lr{ECDH-256} که امروز استانداردهای امنیتی معتبر محسوب می‌شوند، در برابر الگوریتم شور کاملاً شکسته خواهند شد \{1\}.

شکستن \lr{PKI} پیامدهای بسیار گسترده‌ای برای ارتباطات ماهواره‌ای دارد. تمام فرآیندهای احراز هویت \lr{(Authentication)} میان ماهواره و ایستگاه زمینی، کلیه پروتکل‌های تبادل کلید جلسه \lr{(Session Key Exchange)}، و حتی تأیید صحت به‌روزرسانی‌های نرم‌افزاری که از زمین برای ماهواره ارسال می‌شوند، همگی بر \lr{PKI} متکی هستند \{15\}. اگر \lr{PKI} بشکند، یک مهاجم می‌تواند هویت ایستگاه زمینی مجاز را جعل کند، دستورات مخرب به ماهواره ارسال نماید یا کل ترافیک رمزشده را رهگیری کند. علاوه بر این، یک سناریوی خطرناک‌تر نیز وجود دارد که به «اکنون ضبط کن، بعداً رمزگشایی کن» \lr{(Harvest Now, Decrypt Later --- HNDL)} معروف است: مهاجمان هم‌اکنون ارتباطات رمزشده ماهواره‌ای را ضبط می‌کنند تا در آینده پس از دسترسی به کامپیوتر کوانتومی، آن‌ها را بازگشایی کنند \{1\}. این تهدید با توجه به اینکه ماهواره‌ها عمر عملیاتی ۱۰ تا ۱۵ ساله دارند، بسیار جدی‌تر از آن چیزی است که در نگاه اول به نظر می‌رسد \{22\}.

\begin{figure}[H]
\centering
\includegraphics[width=0.8\textwidth]{f1p1.png}
\caption{نمای مفهومی استراتژی «اکنون ضبط کن، بعداً رمزگشایی کن» \lr{(Harvest Now, Decrypt Later --- HNDL)} و تهدید آن برای ارتباطات ماهواره‌ای}
\label{fig:1-1}
\end{figure}

%% ────────────────────────────────────
\subsection{چالش‌های ارتباطات ماهواره‌ای}
%% ────────────────────────────────────

حتی اگر راهکارهای امنیتی نوین مانند \lr{PQC} و \lr{QKD} وجود داشته باشند، پیاده‌سازی آن‌ها بر روی ماهواره‌ها --- به‌ویژه ماهواره‌های در مدار پایین زمین \lr{(Low Earth Orbit --- LEO)} --- با چالش‌های ذاتی و جدی روبروست.

نخستین چالش، محدودیت‌های سخت‌افزاری ماهواره است. پردازنده‌های مقاوم در برابر تشعشع \lr{(Radiation-Hardened Processors)} که برای محیط فضا طراحی می‌شوند، معمولاً سرعت پردازشی در حدود چند ده تا چند صد مگاهرتز دارند و حافظه دسترسی تصادفی \lr{(RAM)} آن‌ها در مقیاس کیلوبایت تا مگابایت محدود است \{22\}. الگوریتم‌های \lr{PQC} مانند \lr{ML-KEM} و \lr{ML-DSA} به دلیل استفاده از ضرب‌های چندجمله‌ای در ابعاد بالا، سربار محاسباتی \lr{(Computational Overhead)} قابل توجهی دارند که اجرای آن‌ها بر روی چنین سخت‌افزار محدودی یک چالش مهندسی واقعی است \{15\}. محدودیت توان مصرفی نیز بر این مشکل می‌افزاید؛ زیرا ماهواره‌ها منبع انرژی محدودی دارند و هر عملیات محاسباتی اضافی به معنای مصرف برق بیشتر و در نتیجه کاهش عمر مفید باتری‌ها یا کاهش توان قابل تخصیص به سایر سیستم‌هاست.

دومین چالش، مشکلات فیزیکی کانال ارتباطی است. ماهواره‌های \lr{LEO} در ارتفاعی بین ۴۰۰ تا ۱۲۰۰ کیلومتری از سطح زمین قرار دارند و با سرعت بالا حرکت می‌کنند. این امر باعث می‌شود پنجره دید \lr{(Visibility Window)} هر ماهواره از یک ایستگاه زمینی تنها حدود ۵ تا ۱۰ دقیقه در هر عبور باشد \{5\}. در مورد \lr{QKD}، این محدودیت به‌ویژه بحرانی است؛ زیرا تولید کلید کوانتومی در چنین فاصله‌ای با اتلاف شدید فوتون‌ها \lr{(Photon Loss)} در مسیر انتشار فضایی و لایه‌های جوی روبروست. بودجه لینک اپتیکی \lr{(Optical Link Budget)} در یک پاس نمونه می‌تواند از حدود ۳۰ تا بیش از ۴۵ دسیبل اتلاف داشته باشد \{23\} و تداخلات جوی مانند آشفتگی‌های اتمسفری \lr{(Atmospheric Turbulence)} نرخ خطای بیت کوانتومی \lr{(Quantum Bit Error Rate --- QBER)} را به‌شدت افزایش می‌دهد. نرخ کلید تولیدشده توسط سیستم‌های \lr{QKD} ماهواره‌ای در بهترین شرایط در حد ده‌ها کیلوبیت بر ثانیه است \{5\}، که برای پوشش نیازهای ارتباطی پیوسته به هیچ وجه کافی نیست.

در مجموع، نه \lr{PQC} به‌تنهایی و نه \lr{QKD} به‌تنهایی نمی‌توانند پاسخ کاملی به تمام نیازهای امنیتی ارتباطات ماهواره‌ای در عصر کوانتومی بدهند. این شکاف، ضرورت طراحی یک معماری ترکیبی هوشمند را که نقاط قوت هر دو رویکرد را هم‌زمان بهره‌برداری کند، آشکار می‌سازد \{19\}.
%% ──────────────────────────────────────────────────────────────────────────
\section{اهمیت و ضرورت تحقیق}
%% ──────────────────────────────────────────────────────────────────────────

%% ────────────────────────────────────
\subsection{تهدید \lr{HNDL} و فوریت اقدام}
%% ────────────────────────────────────

یکی از مهم‌ترین دلایلی که نمی‌توان برای مقابله با تهدیدات کوانتومی صبر کرد، استراتژی «اکنون ضبط کن، بعداً رمزگشایی کن» \lr{(Harvest Now, Decrypt Later --- HNDL)} است. در این استراتژی، مهاجمان از همین امروز ترافیک رمزشده ارتباطات ماهواره‌ای را ضبط و ذخیره می‌کنند، بدون آنکه در حال حاضر توانایی رمزگشایی آن را داشته باشند \{1\}. هدف این است که پس از دسترسی به کامپیوترهای کوانتومی قدرتمند در آینده، این داده‌ها بازگشایی شوند.

آنچه این تهدید را به‌ویژه برای ماهواره‌ها خطرناک می‌کند، عمر عملیاتی بلند آن‌هاست. یک ماهواره که امروز پرتاب می‌شود ممکن است ۱۰ تا ۱۵ سال در مدار فعال باشد \{22\}. این یعنی داده‌هایی که امروز با \lr{RSA} یا \lr{ECC} رمزنگاری می‌شوند، تا ۱۵ سال دیگر نیز باید محرمانه بمانند. اما اگر در این بازه زمانی کامپیوترهای کوانتومی به بلوغ برسند، تمام آن داده‌های ضبط‌شده قابل رمزگشایی خواهند بود. از آنجا که به‌روزرسانی سیستم رمزنگاری یک ماهواره پس از پرتاب بسیار دشوار و گران است، انتخاب الگوریتم‌های مقاوم در برابر کوانتوم پیش از پرتاب یک الزام غیرقابل اجتناب است \{22\}. به‌عبارت دیگر، اگر امروز اقدام نشود، هیچ راه بازگشتی وجود نخواهد داشت.

%% ────────────────────────────────────
\subsection{اهمیت راهبردی ایمن‌سازی ماهواره‌ها}
%% ────────────────────────────────────

امنیت ارتباطات ماهواره‌ای صرفاً یک مسئله فنی نیست؛ بلکه مستقیماً با امنیت ملی و زیرساخت‌های حیاتی جهانی پیوند دارد. ماهواره‌ها امروز ستون فقرات سیستم‌های ناوبری جهانی (مانند \lr{GPS})، شبکه‌های ارتباطی نظامی، انتقال داده‌های مالی و بانکی بین‌المللی، و پایش محیط‌زیست و هواشناسی هستند \{14\}. اختلال یا نفوذ در هر یک از این زیرساخت‌ها می‌تواند پیامدهای فاجعه‌بار اقتصادی، نظامی و اجتماعی در مقیاس جهانی داشته باشد.

از این منظر، ایمن‌سازی ارتباطات ماهواره‌ای در برابر تهدیدات کوانتومی یک اولویت راهبردی \lr{(Strategic Priority)} برای دولت‌ها، سازمان‌های نظامی و شرکت‌های فناوری محسوب می‌شود. پژوهش‌های اخیر نشان می‌دهند که یک معماری ترکیبی از \lr{PQC} و \lr{QKD} می‌تواند هم از نظر محاسباتی و هم از نظر فیزیکی، سطحی از امنیت فراهم کند که هیچ‌یک از این دو رویکرد به‌تنهایی قادر به ارائه آن نیستند \{13\}. پایان‌نامه حاضر در همین راستا تلاش می‌کند تا با بررسی دقیق این معماری‌ها، گامی عملی در مسیر ایمن‌سازی زیرساخت‌های فضایی در برابر تهدیدات عصر کوانتومی بردارد.

%% ──────────────────────────────────────────────────────────────────────────
\section{اهداف تحقیق}
%% ──────────────────────────────────────────────────────────────────────────

\textbf{هدف اصلی}

هدف اصلی این پایان‌نامه، پیشنهاد، طراحی و ارزیابی یک معماری امنیتی ترکیبی \lr{(Hybrid Security Architecture)} مبتنی بر رمزنگاری پساکوانتومی \lr{(Post-Quantum Cryptography --- PQC)} و توزیع کلید کوانتومی \lr{(Quantum Key Distribution --- QKD)} است، به گونه‌ای که بتواند ارتباطات ماهواره‌ای --- به‌ویژه در مدار پایین زمین \lr{(Low Earth Orbit --- LEO)}  را در برابر تهدیدات کوانتومی کنونی و آینده به‌صورت همزمان ایمن سازد.

\medskip
\textbf{اهداف فرعی}

\textbf{شناسایی و تحلیل آسیب‌پذیری‌های موجود:} بررسی نقاط ضعف زیرساخت‌های امنیتی فعلی ماهواره‌ها، از جمله پروتکل‌های احراز هویت و تبادل کلید مبتنی بر \lr{RSA} و \lr{ECC}، در برابر حملات کوانتومی نظیر الگوریتم شور \lr{(Shor's Algorithm)}.

\textbf{بررسی محدودیت‌های ذاتی ماهواره‌های \lr{LEO}:} تحلیل چالش‌های سخت‌افزاری و فیزیکی محیط فضایی شامل توان پردازشی محدود، مصرف انرژی، پنجره دید کوتاه و بودجه لینک اپتیکی \lr{(Optical Link Budget)}، به‌منظور درک الزامات واقعی پیاده‌سازی الگوریتم‌های امنیتی روی ماهواره.

\textbf{ارزیابی عملکردی الگوریتم‌های استاندارد \lr{PQC}:} سنجش و مقایسه الگوریتم‌های \lr{ML-KEM} و \lr{ML-DSA} که توسط مؤسسه ملی استاندارد و فناوری \lr{(NIST)} استانداردسازی شده‌اند  از نظر اندازه کلید \lr{(Key Size)}، زمان اجرا \lr{(Execution Time)} و سربار پهنای باند \lr{(Bandwidth Overhead)} در مقایسه با الگوریتم‌های کلاسیک \lr{RSA} و \lr{ECC}.

\textbf{تحلیل چالش‌های پیاده‌سازی \lr{QKD} در فضا:} بررسی عوامل مؤثر بر کیفیت کانال کوانتومی ماهواره‌ای، از جمله نرخ خطای بیت کوانتومی \lr{(Quantum Bit Error Rate --- QBER)}، اتلاف فوتون در مسیر انتشار فضایی، و تأثیر اغتشاشات جوی بر نرخ تولید کلید، به‌منظور تعیین محدوده عملکردی واقعی \lr{QKD} ماهواره‌ای.

%% ──────────────────────────────────────────────────────────────────────────
\section{سؤالات تحقیق}
%% ──────────────────────────────────────────────────────────────────────────

\textbf{سؤال اصلی}

چگونه می‌توان یک معماری امنیتی ترکیبی \lr{(Hybrid Security Architecture)} مبتنی بر رمزنگاری پساکوانتومی \lr{(Post-Quantum Cryptography --- PQC)} و توزیع کلید کوانتومی \lr{(Quantum Key Distribution --- QKD)} برای ارتباطات ماهواره‌ای در مدار پایین زمین \lr{(Low Earth Orbit --- LEO)} طراحی کرد که ضمن تأمین امنیت در برابر تهدیدات کامپیوترهای کوانتومی، با محدودیت‌های پردازشی، انرژی و فیزیکی ذاتی ماهواره‌ها سازگار باشد؟

\medskip
\textbf{سؤالات فرعی}

\textbf{سؤال فرعی اول:} آسیب‌پذیری‌های اصلی زیرساخت کلید عمومی \lr{(Public Key Infrastructure --- PKI)} فعلی ماهواره‌ها در برابر حملات کوانتومی نظیر الگوریتم شور کدامند و شکست این زیرساخت چه پیامدهایی برای امنیت ارتباطات ماهواره‌ای خواهد داشت؟

\textbf{سؤال فرعی دوم:} الگوریتم‌های استاندارد \lr{PQC} از جمله \lr{ML-KEM} و \lr{ML-DSA} از نظر اندازه کلید، زمان اجرا و سربار پهنای باند چه تفاوت عملکردی نسبت به الگوریتم‌های کلاسیک \lr{RSA} و \lr{ECC} در محیط محدود ماهواره‌ای دارند؟

\textbf{سؤال فرعی سوم:} محدودیت‌های فیزیکی پیاده‌سازی \lr{QKD} در کانال‌های ماهواره‌ای \lr{LEO}  از جمله نرخ خطای بیت کوانتومی \lr{(QBER)}، اتلاف فوتون و کوتاهی پنجره دید  تا چه حد نرخ تولید کلید را محدود می‌سازند و آیا \lr{QKD} به‌تنهایی برای تأمین نیازهای عملیاتی کافی است؟

%% ──────────────────────────────────────────────────────────────────────────
\section{فرضیه‌های تحقیق}
%% ──────────────────────────────────────────────────────────────────────────

\textbf{فرضیه اصلی}

به نظر می‌رسد طراحی یک معماری ترکیبی که در آن الگوریتم‌های \lr{PQC}  به‌ویژه \lr{ML-KEM} برای کپسوله‌سازی کلید \lr{(Key Encapsulation)} و \lr{ML-DSA} برای احراز هویت دیجیتال \lr{(Digital Authentication)}  مسئولیت لایه رمزنگاری محاسباتی را بر عهده داشته باشند و \lr{QKD} به‌عنوان یک لایه امنیت فیزیکی مکمل برای تولید و توزیع کلیدهای اطلاعات‌تئوریک \lr{(Information-Theoretically Secure Keys)} عمل کند، می‌تواند یک راهکار امن، عملی و سازگار با محدودیت‌های ماهواره‌های \lr{LEO} فراهم آورد.

\medskip
\textbf{فرضیه فرعی اول:} به نظر می‌رسد زیرساخت‌های امنیتی فعلی ماهواره‌ها که بر الگوریتم‌های \lr{RSA} و \lr{ECC} متکی هستند، در برابر الگوریتم شور به‌طور کامل شکست‌پذیر بوده و در صورت دسترسی مهاجم به یک کامپیوتر کوانتومی کافی بزرگ، تمام مکانیزم‌های احراز هویت و تبادل کلید مبتنی بر این الگوریتم‌ها بی‌اعتبار خواهند شد.

\textbf{فرضیه فرعی دوم:} به نظر می‌رسد الگوریتم‌های \lr{PQC} در مقایسه با معادل‌های کلاسیک خود دارای سربار محاسباتی و اندازه کلید بیشتری هستند، اما با توجه به پیشرفت‌های پیاده‌سازی نرم‌افزاری و بهینه‌سازی‌های موجود، اجرای آن‌ها در محدوده قابل قبول برای سخت‌افزار ماهواره‌های \lr{LEO} امکان‌پذیر خواهد بود.

\textbf{فرضیه فرعی سوم:} به نظر می‌رسد \lr{QKD} ماهواره‌ای به دلیل محدودیت‌های فیزیکی  از جمله اتلاف شدید فوتون در کانال فضایی و کوتاه بودن پنجره دید در مدار \lr{LEO}  به‌تنهایی نمی‌تواند نیاز به ارتباط امن پیوسته را تأمین کند، اما به‌عنوان یک لایه امنیتی مکمل در کنار \lr{PQC}، سطح اطمینان کلی سیستم را به‌طور معناداری افزایش می‌دهد.

%% ──────────────────────────────────────────────────────────────────────────
\section{نوآوری تحقیق}
%% ──────────────────────────────────────────────────────────────────────────

مرور ادبیات موجود نشان می‌دهد که پژوهش‌های پیشین در حوزه امنیت ارتباطات ماهواره‌ای عمدتاً در یکی از سه دسته محدود قرار می‌گیرند: گروه اول، امنیت ماهواره‌ها را صرفاً از منظر رمزنگاری کلاسیک \lr{(Classical Cryptography)} بررسی کرده و به تهدیدات کوانتومی توجهی نداشته‌اند. گروه دوم، بر الگوریتم‌های \lr{PQC} تمرکز کرده‌اند اما آن‌ها را عمدتاً در محیط‌های زمینی با منابع پردازشی فراوان ارزیابی نموده‌اند و واقعیت‌های سخت‌افزاری ماهواره را نادیده گرفته‌اند. گروه سوم نیز \lr{QKD} ماهواره‌ای را به‌صورت مستقل و صرفاً از دیدگاه فیزیک کوانتومی مطرح کرده‌اند، بدون آنکه آن را در قالب یک معماری امنیتی یکپارچه با لایه‌های نرم‌افزاری قرار دهند.

نوآوری اصلی پایان‌نامه حاضر در آن است که این سه دیدگاه را در قالب یک چارچوب تحلیلی جامع و یکپارچه گرد هم می‌آورد. به‌طور مشخص، این پژوهش یک معماری ترکیبی \lr{(Hybrid Architecture)} از \lr{PQC} و \lr{QKD} پیشنهاد می‌دهد که محدودیت‌های واقعی و عملیاتی ماهواره‌های \lr{LEO}  از جمله توان پردازشی محدود، مصرف انرژی، کوتاهی پنجره دید و افت کانال اپتیکی  را به‌صورت همزمان در نظر می‌گیرد. علاوه بر این، از طریق شبیه‌سازی مقایسه‌ای الگوریتم‌های \lr{PQC} با معادل‌های کلاسیک آن‌ها، یک ارزیابی کمّی از قابلیت اجرای این معماری در محیط ماهواره‌ای ارائه می‌شود.

%% ──────────────────────────────────────────────────────────────────────────
\section{قلمرو تحقیق}
%% ──────────────────────────────────────────────────────────────────────────

برای حفظ انسجام و عمق تحلیل، دامنه این پژوهش در سه بُعد زیر تعریف و محدود شده است:

\textbf{قلمرو موضوعی \lr{(Thematic Scope)}:} این پژوهش در حوزه امنیت سایبری \lr{(Cybersecurity)} ارتباطات ماهواره‌ای قرار می‌گیرد و تمرکز ویژه آن بر دو حوزه رمزنگاری پساکوانتومی \lr{(Post-Quantum Cryptography --- PQC)} و توزیع کلید کوانتومی \lr{(Quantum Key Distribution --- QKD)} است. مسائل مرتبط با لایه فیزیکی رادیویی، پروتکل‌های شبکه ماهواره‌ای یا امنیت فیزیکی سخت‌افزار خارج از محدوده این پژوهش قرار دارند.

\textbf{قلمرو مکانی و فضایی \lr{(Spatial Scope)}:} بررسی‌های این پایان‌نامه محدود به ماهواره‌های مستقر در مدار پایین زمین \lr{(Low Earth Orbit --- LEO)} در محدوده ارتفاعی ۴۰۰ تا ۱۲۰۰ کیلومتری از سطح زمین است. مدارهای میانی \lr{(Medium Earth Orbit --- MEO)} و زمین‌آهنگ \lr{(Geostationary Orbit --- GEO)} به‌عنوان مقایسه پس‌زمینه‌ای مطرح می‌شوند اما محور اصلی تحلیل نیستند.

\textbf{قلمرو زمانی \lr{(Temporal Scope)}:} این پژوهش بر استانداردها و پژوهش‌های منتشرشده تا پایان سال ۲۰۲۴ تکیه می‌کند. به‌ویژه، استانداردهای \lr{FIPS 203} و \lr{FIPS 204} منتشرشده توسط مؤسسه ملی استاندارد و فناوری \lr{(NIST)} در سال ۲۰۲۴ به‌عنوان مبنای اصلی ارزیابی الگوریتم‌های \lr{PQC} در نظر گرفته شده‌اند.

%% ──────────────────────────────────────────────────────────────────────────
\section{روش‌شناسی کلی تحقیق}
%% ──────────────────────────────────────────────────────────────────────────

این پژوهش از یک رویکرد ترکیبی توصیفی-تحلیلی و کاربردی \lr{(Descriptive-Analytical and Applied Approach)} بهره می‌گیرد. به این معنا که در کنار مرور و تحلیل ادبیات علمی موجود، از ابزارهای شبیه‌سازی نرم‌افزاری و مدل‌سازی ریاضی نیز برای تولید داده‌های ارزیابی استفاده می‌شود. جزئیات این رویکرد در دو محور اصلی زیر تعریف شده است:

\textbf{ارزیابی نرم‌افزاری الگوریتم‌های \lr{PQC}:} برای سنجش عملکرد الگوریتم‌های \lr{ML-KEM} و \lr{ML-DSA} در محیط محدود ماهواره‌ای، از شبیه‌سازی نرم‌افزاری \lr{(Software Simulation)} با زبان برنامه‌نویسی پایتون \lr{(Python)} و کتابخانه \lr{liboqs} استفاده می‌شود. سنجه‌های \lr{(Metrics)} اصلی مورد اندازه‌گیری شامل اندازه کلید \lr{(Key Size)}، زمان اجرای عملیات رمزنگاری \lr{(Execution Time)} و سربار پهنای باند \lr{(Bandwidth Overhead)} در مقایسه با الگوریتم‌های کلاسیک \lr{RSA} و \lr{ECC} خواهند بود \{15\}.

\textbf{مدل‌سازی تحلیلی کانال \lr{QKD}:} به دلیل محدودیت دسترسی به سخت‌افزار فضایی واقعی، ارزیابی عملکرد \lr{QKD} ماهواره‌ای از طریق مدل‌سازی ریاضی \lr{(Mathematical Modeling)} و تحلیل داده‌های استخراج‌شده از پژوهش‌های معتبر پیشین انجام می‌شود. در این بخش، پارامترهایی مانند بودجه لینک اپتیکی \lr{(Optical Link Budget)} و نرخ خطای بیت کوانتومی \lr{(Quantum Bit Error Rate --- QBER)} برای سناریوی مدار \lr{LEO} بررسی خواهند شد \{23\}.

در نهایت، خروجی‌های هر دو محور فوق با یکدیگر تلفیق شده و مبنای طراحی معماری ترکیبی نهایی \lr{(Hybrid PQC+QKD Architecture)} پیشنهادی این پایان‌نامه را تشکیل می‌دهند \{19\}.

%% ──────────────────────────────────────────────────────────────────────────
\section{تعریف واژگان کلیدی}
%% ──────────────────────────────────────────────────────────────────────────

\textbf{توزیع کلید کوانتومی \lr{(Quantum Key Distribution --- QKD)}:} روشی مبتنی بر اصول فیزیک کوانتومی برای تبادل امن کلید رمزنگاری میان دو طرف ارتباط است، به گونه‌ای که هرگونه تلاش برای استراق‌سمع \lr{(Eavesdropping)} به‌طور اجتناب‌ناپذیر اختلالی قابل تشخیص در کانال ایجاد می‌کند و دو طرف را از وجود شنودگر آگاه می‌سازد.

\textbf{رمزنگاری پساکوانتومی \lr{(Post-Quantum Cryptography --- PQC)}:} مجموعه‌ای از الگوریتم‌های رمزنگاری نرم‌افزاری کلاسیک است که بر پایه مسائل ریاضی‌ای طراحی شده‌اند که حتی کامپیوترهای کوانتومی قدرتمند نیز قادر به حل مؤثر آن‌ها نیستند و بنابراین در برابر حملات کوانتومی مقاوم \lr{(Quantum-Resistant)} هستند.

\textbf{مدار پایین زمین \lr{(Low Earth Orbit --- LEO)}:} ناحیه‌ای از فضا در ارتفاع تقریبی ۲۰۰ تا ۲۰۰۰ کیلومتری از سطح زمین است (که قلمرو این پژوهش بازهٔ ۴۰۰ تا ۱۲۰۰ کیلومتر آن است) که ماهواره‌های مستقر در آن تأخیر انتشار سیگنال \lr{(Propagation Delay)} بسیار پایین‌تری نسبت به مدارهای بالاتر دارند و برای کاربردهایی که به زمان‌بندی دقیق نیاز دارند مناسب‌ترند.

\textbf{اکنون ضبط کن، بعداً رمزگشایی کن \lr{(Harvest Now, Decrypt Later --- HNDL)}:} استراتژی حمله‌ای که در آن مهاجم ارتباطات رمزشده کنونی را ضبط و ذخیره می‌کند تا در آینده، پس از دسترسی به کامپیوترهای کوانتومی، آن‌ها را رمزگشایی نماید؛ تهدیدی که اقدام فوری برای انتقال به رمزنگاری مقاوم در برابر کوانتوم را ضروری می‌سازد.

\textbf{ماشین حالت متناهی امنیتی \lr{(Security Finite State Machine --- FSM)}:} مدل محاسباتی است که در معماری پیشنهادی این پایان‌نامه برای مدیریت لایه‌های امنیتی و انتقال کنترل‌شده میان حالت‌های عملیاتی مختلف سیستم --- از جمله اجرای عادی \lr{QKD}، افت کیفیت کانال، و بازگشت به لایه \lr{PQC} --- به کار می‌رود.
%% ──────────────────────────────────────────────────────────────────────────
\section{ساختار پایان‌نامه}
%% ──────────────────────────────────────────────────────────────────────────

پایان‌نامه حاضر در پنج فصل سازمان‌دهی شده است. فصل اول که همین فصل است، با معرفی کلیات تحقیق شامل مقدمه، بیان مسئله، اهمیت و ضرورت، اهداف، سؤالات و فرضیه‌ها، پایه مفهومی پژوهش را بنا نهاده است.

فصل دوم به مرور جامع ادبیات و پیشینه تحقیق \lr{(Literature Review)} اختصاص دارد. در این فصل، مبانی نظری رمزنگاری مدرن، تهدیدات کوانتومی، الگوریتم‌های \lr{PQC}، پروتکل‌های \lr{QKD} و پژوهش‌های پیشین در حوزه امنیت ماهواره‌ای به‌تفصیل بررسی می‌شوند.

فصل سوم روش‌شناسی تحقیق \lr{(Research Methodology)} و معماری ترکیبی پیشنهادی را تشریح می‌کند. در این فصل، مدل سیستم، توپولوژی شبکه، پروتکل‌های تبادل کلید و طراحی شبیه‌سازی نرم‌افزاری به‌صورت کامل ارائه می‌شوند.

فصل چهارم به تجزیه و تحلیل داده‌ها \lr{(Data Analysis)} می‌پردازد و نتایج شبیه‌سازی الگوریتم‌های \lr{PQC}، ارزیابی امنیتی معماری و تحلیل چالش‌های \lr{QKD} ماهواره‌ای را با جداول و نمودارهای مقایسه‌ای ارائه می‌دهد.

فصل پنجم به نتیجه‌گیری \lr{(Conclusion)} و پیشنهادات آتی اختصاص دارد و با جمع‌بندی یافته‌های کلیدی، پاسخ به سؤالات تحقیق و ارائه راهکارهایی برای پژوهش‌های آینده، پایان می‌پذیرد.

\begin{figure}[H]
\centering
\includegraphics[width=0.75\textwidth]{f1p2.png}
\caption{نقشه راه ساختار فصل‌های پایان‌نامه}
\label{fig:1-2}
\end{figure}

%%%%%%%%%%%%%%%%%%%%%%%%%%%%%%%%%%%%%%%%%%%%%%%%%%%%%%%%%%%%%%%%%%%%%%%%%%%%%
%%  — پایان فصل اول —
%%  فصل‌های بعدی (۲، ۳، ۴، ۵) در گام بعدی افزوده می‌شوند
%%%%%%%%%%%%%%%%%%%%%%%%%%%%%%%%%%%%%%%%%%%%%%%%%%%%%%%%%%%%%%%%%%%%%%%%%%%%%
%%%%%%%%%%%%%%%%%%%%%%%%%%%%%%%%%%%%%%%%%%%%%%%%%%%%%%%%%%%%%%%%%%%%%%%%%%%%%
%%  chapter2.tex — فصل دوم: تاریخچه و پیشینه تحقیق
%%  این فایل را در main.tex پیش از \end{document} paste کنید
%%  (بعد از پایان فصل اول و قبل از \end{document})
%%%%%%%%%%%%%%%%%%%%%%%%%%%%%%%%%%%%%%%%%%%%%%%%%%%%%%%%%%%%%%%%%%%%%%%%%%%%%

%%%%%%%%%%%%%%%%%%%%%%%%%%%%%%%%%%%%%%%%%%%%%%%%%%%%%%%%%%%%%%%%%%%%%%%%%%%%%
%%  فصل دوم: تاریخچه و پیشینه تحقیق
%%%%%%%%%%%%%%%%%%%%%%%%%%%%%%%%%%%%%%%%%%%%%%%%%%%%%%%%%%%%%%%%%%%%%%%%%%%%%
\chapter{تاریخچه و پیشینه تحقیق}
\setcounter{figure}{0}
\setcounter{table}{0}
\setcounter{equation}{0}

%% ──────────────────────────────────────────────────────────────────────────
\section{مقدمه}
%% ──────────────────────────────────────────────────────────────────────────

درک عمیق چالش‌های امنیتی در ارتباطات ماهواره‌ای مستلزم آن است که پیش از هر بحثی، پایه‌های نظری موضوع به‌دقت بررسی شوند. این فصل با هدف ارائه بستر مفهومی لازم برای فهم معماری‌های امنیتی پیشنهادی، در قالب یک مسیر پیوسته و منطقی تنظیم شده است. نخست، مبانی رمزنگاری مدرن \lr{(Modern Cryptography)} شامل الگوریتم‌های متقارن و نامتقارن متداول معرفی می‌شوند. سپس تهدیداتی که محاسبات کوانتومی \lr{(Quantum Computing)} برای این سیستم‌ها ایجاد می‌کنند، از جمله الگوریتم‌های شور \lr{(Shor)} و گرور \lr{(Grover)}، تحلیل می‌گردند. در ادامه، دو راهکار اصلی برای مقابله با این تهدیدات، یعنی رمزنگاری پساکوانتومی \lr{(Post-Quantum Cryptography)} و توزیع کلید کوانتومی \lr{(Quantum Key Distribution)}، به تفصیل بررسی می‌شوند. سرانجام، پژوهش‌های پیشین در حوزه ارتباطات ماهواره‌ای ایمن مرور می‌شوند تا شکاف‌های پژوهشی \lr{(Research Gaps)} موجود مشخص و زمینه برای فصل‌های بعدی فراهم گردد.

%% ──────────────────────────────────────────────────────────────────────────
\section{مبانی رمزنگاری مدرن}
%% ──────────────────────────────────────────────────────────────────────────

رمزنگاری مدرن \lr{(Modern Cryptography)} بر دو پایه اصلی استوار است: رمزنگاری نامتقارن \lr{(Asymmetric Cryptography)} برای تبادل کلید و احراز هویت، و رمزنگاری متقارن \lr{(Symmetric Cryptography)} برای رمزنگاری حجم بالای داده. درک نقاط قوت و ضعف هر یک از این دسته‌ها برای ارزیابی تهدیدات کوانتومی و ضرورت گذار به رمزنگاری پساکوانتومی ضروری است.

در رمزنگاری نامتقارن، هر کاربر دارای یک جفت کلید عمومی \lr{(Public Key)} و خصوصی \lr{(Private Key)} است. پرکاربردترین این الگوریتم‌ها، \lr{RSA (Rivest--Shamir--Adleman)}، است که امنیت آن بر پایه دشواری محاسباتی تجزیه عدد صحیح بزرگ \lr{N = p × q} به عوامل اول بنا شده است. شکستن \lr{RSA} مستلزم یافتن \lr{p} و \lr{q} از روی \lr{N} است، مسئله‌ای که برای کلیدهای بزرگ با محاسبات کلاسیک غیرعملی است \{1\}. \lr{RSA} به‌طور گسترده در پروتکل‌های \lr{TLS}، \lr{SSH} و \lr{IPsec} که زیرساخت ارتباطات ماهواره‌ای بر آن‌ها متکی است استفاده می‌شود \{2\}. با این حال، نقطه ضعف بنیادین \lr{RSA} طول کلید بسیار بلند آن است؛ دستیابی به امنیت ۱۲۸ بیتی به کلید ۳۰۷۲ بیتی نیاز دارد که سربار محاسباتی و ارتباطی قابل توجهی ایجاد می‌کند \{22\}.

\lr{ECC (Elliptic Curve Cryptography)} رویکردی بهینه‌تر دارد و امنیت خود را از دشواری مسئله لگاریتم گسسته روی منحنی‌های بیضوی \lr{(Elliptic Curve Discrete Logarithm Problem)} می‌گیرد \{1\}. مزیت کلیدی \lr{ECC} در کوتاهی کلیدهاست؛ یک کلید ۲۵۶ بیتی \lr{ECC} امنیتی معادل کلید ۳۰۷۲ بیتی \lr{RSA} ارائه می‌دهد \{2\}. این ویژگی \lr{ECC} را برای محیط‌های با منابع محدود مانند ماهواره‌ها بسیار جذاب کرده است. در پروتکل‌های احراز هویت ماهواره‌ای، \lr{ECDSA (Elliptic Curve Digital Signature Algorithm)} جایگاه ویژه‌ای دارد، اگرچه پروتکل‌های مبتنی بر \lr{ECDSA} در برابر حملاتی نظیر \lr{signal leakage} که از پیاده‌سازی‌های ناامن نشأت می‌گیرند، آسیب‌پذیر بوده‌اند \{18\}. هر دوی این الگوریتم‌ها عمدتاً برای تبادل کلید \lr{(Key Exchange)} و امضای دیجیتال \lr{(Digital Signature)} به‌کار می‌روند و نه رمزنگاری مستقیم داده.

برای رمزنگاری حجم بالای داده پس از برقراری کانال امن، از الگوریتم‌های متقارن استفاده می‌شود. استاندارد \lr{AES (Advanced Encryption Standard)} با طول‌های کلید ۱۲۸، ۱۹۲ یا ۲۵۶ بیتی، رایج‌ترین انتخاب در این حوزه است. \lr{AES} سرعت پردازشی بسیار بالاتری نسبت به الگوریتم‌های نامتقارن دارد و در معماری‌های ارتباطات ماهواره‌ای، پس از تبادل کلید جلسه \lr{(Session Key)} توسط \lr{RSA} یا \lr{ECC}، رمزنگاری داده‌های تله‌متری و فرمان را بر عهده می‌گیرد \{14\}. در چارچوب‌های ترکیبی پیشرفته‌تر، \lr{AES-256} به‌عنوان لایه پشتیبان در کنار لایه کوانتومی نیز به‌کار می‌رود \{13\}. اما به رغم کارایی \lr{AES}، امنیت کل زنجیره ارتباطی تنها به اندازه ضعیف‌ترین حلقه آن است — و آن حلقه همان الگوریتم‌های نامتقارن \lr{RSA} و \lr{ECC} هستند که در بخش بعدی نشان داده می‌شود چرا در برابر کامپیوترهای کوانتومی کاملاً آسیب‌پذیرند \{1\}.

%% ─── جدول ۲-۱ ──────────────────────────────────────────────────────────────
\begin{table}[H]
\caption{مقایسه ویژگی‌های الگوریتم‌های کلاسیک \lr{RSA} و \lr{ECC}}
\label{tab:2-1}
\centering
\small
\begin{tabular}{|r|p{4.5cm}|p{4.5cm}|}
\hline
\bfseries ویژگی & \bfseries الگوریتم \lr{RSA} & \bfseries الگوریتم \lr{ECC} \\
\hline
مبنای ریاضی &
دشواری تجزیه عدد صحیح بزرگ \lr{(Integer Factorization)} &
دشواری لگاریتم گسسته روی منحنی بیضوی \lr{(ECDLP)} \\
\hline
طول کلید برای امنیت ۱۲۸ بیتی & ۳۰۷۲ بیت & ۲۵۶ بیت \\
\hline
طول کلید برای امنیت ۱۹۲ بیتی & ۷۶۸۰ بیت & ۳۸۴ بیت \\
\hline
سرعت تولید کلید &
کند \lr{(\textasciitilde{}240 ms)} &
سریع‌تر \lr{(\textasciitilde{}0.15 ms)} \\
\hline
حجم سربار ارتباطی &
زیاد (کلید عمومی \textasciitilde{}384 بایت برای \lr{RSA-3072}) &
کم (کلید عمومی \textasciitilde{}64 بایت برای \lr{P-256}) \\
\hline
مناسب برای ماهواره & ضعیف  کلیدهای حجیم & مناسب‌تر  فشردگی کلید \\
\hline
مقاومت کوانتومی &
ندارد (الگوریتم شور آن را کاملاً می‌شکند) &
ندارد (الگوریتم شور آن را کاملاً می‌شکند) \\
\hline
کاربرد رایج در ماهواره &
\lr{TLS}، \lr{IPsec}، احراز هویت زمین-ماهواره &
\lr{ECDSA} در پروتکل‌های احراز هویت \lr{LEO} \\
\hline
\end{tabular}
\end{table}

\begin{figure}[H]
\centering
\includegraphics[width=0.78\textwidth]{f2p1.png}
\caption{مقایسه معماری رمزنگاری نامتقارن \lr{(RSA/ECC)} و متقارن \lr{(AES)} در ارتباطات امن}
\label{fig:2-1}
\end{figure}

%% ──────────────────────────────────────────────────────────────────────────
\section{تهدیدات کوانتومی}
%% ──────────────────────────────────────────────────────────────────────────

ظهور محاسبات کوانتومی \lr{(Quantum Computing)} نه یک تهدید دور و انتزاعی، بلکه یک بحران امنیتی در حال وقوع است. در سال ۲۰۲۶، اجماع پژوهشگران بر این است که دستیابی به یک کامپیوتر کوانتومی از نظر رمزنگاری مرتبط \lr{(Cryptographically Relevant Quantum Computer)} دیگر مسئله‌ای با افق ده‌ها ساله نیست، و این تهدید از سه بُعد اصلی قابل بررسی است.

\textbf{الگوریتم شور \lr{(Shor's Algorithm)}} که در ۱۹۹۴ ارائه شد، ثابت کرد که یک کامپیوتر کوانتومی کافی می‌تواند مسئله تجزیه عدد صحیح بزرگ و مسئله لگاریتم گسسته را — پایه ریاضی \lr{RSA} و \lr{ECC} — در زمان چندجمله‌ای \lr{(Polynomial Time)} با پیچیدگی \lr{O(n³ log n)} حل کند \{1\}. این یعنی تمام پروتکل‌هایی که زیرساخت ارتباطات ماهواره‌ای بر آن‌ها متکی است — از \lr{TLS} و \lr{IPsec} برای احراز هویت فرمان‌ها تا \lr{ECDSA} برای امضای به‌روزرسانی‌های نرم‌افزاری — به‌طور همزمان در معرض خطر قرار خواهند گرفت \{2\}. با توجه به اینکه استانداردهای موجود رمزنگاری کلید عمومی مانند \lr{RSA}، \lr{DSA} و \lr{ECDSA} در پروتکل‌هایی مثل \lr{SSH}، \lr{IKE} و \lr{DNSSEC} استفاده می‌شوند، الگوریتم شور تهدیدی وجودی برای کل زیرساخت کلید عمومی \lr{(Public Key Infrastructure)} جهانی محسوب می‌شود \{2\}.

\textbf{الگوریتم گرور \lr{(Grover's Algorithm)}} تهدیدی متفاوت اما قابل مدیریت‌تر است. این الگوریتم فضای جستجوی \lr{N} عنصری را با تقریباً \lr{√N} ارزیابی کوانتومی پیمایش می‌کند و به این ترتیب طول کلید مؤثر رمزنگاری‌های متقارن را به نصف کاهش می‌دهد \{1\}. عملاً \lr{AES-128} در برابر یک مهاجم مجهز به \lr{CRQC} به امنیتی معادل ۶۴ بیت تقلیل می‌یابد که قابل قبول نیست؛ در مقابل، \lr{AES-256} با همان مهاجم همچنان ۱۲۸ بیت امنیت مؤثر خواهد داشت. برخلاف الگوریتم شور که رمزنگاری نامتقارن را کاملاً می‌شکند، راهکار مقابله با گرور ساده است: دوبرابر کردن طول کلیدهای متقارن، چنانکه استفاده از \lr{AES-256} به‌جای \lr{AES-128} در معماری‌های پیشنهادی برای کانال‌های ماهواره‌ای اکنون اجماع پژوهشگران است \{22\}.

اما خطرناک‌ترین بُعد این تهدید در سال ۲۰۲۶، نه لزوماً روز کوانتوم \lr{(Q-Day)} به‌عنوان یک رخداد آینده، بلکه حمله‌ای است که هم‌اکنون در جریان است: \textbf{«اکنون ضبط کن، بعداً رمزگشایی کن»} \lr{(Harvest Now, Decrypt Later --- HNDL)}. در این استراتژی، مهاجمان ترافیک رمزشده کنونی را انبوه‌سازی می‌کنند تا پس از دستیابی به \lr{CRQC} آن را رمزگشایی نمایند \{1\}. این تهدید برای ارتباطات ماهواره‌ای از ویژگی منحصربه‌فردی برخوردار است: ماهواره‌ها عمر عملیاتی ۱۰ تا ۱۵ ساله دارند \{22\} و تله‌متری، فرمان‌های کنترل و داده‌های حساسی که امروز با \lr{RSA/ECC} رمزنگاری می‌شوند، ممکن است با ارزش استراتژیک بلندمدتی باشند که ضبط آن‌ها برای رمزگشایی آینده کاملاً توجیه‌پذیر است \{15\}. ماهواره‌هایی که امروز پرتاب می‌شوند و با رمزنگاری‌های کلاسیک محافظت می‌شوند، پیش از پایان عمر عملیاتی‌شان می‌توانند در بازه زمانی وقوع \lr{Q-Day} قرار گیرند \{18\}. از این رو، در سال ۲۰۲۶ دیگر نمی‌توان گذار به رمزنگاری پساکوانتومی را به آینده موکول کرد؛ این گذار یک ضرورت فوری است.

\begin{figure}[H]
\centering
\fbox{\parbox{0.78\textwidth}{\centering\vspace{1.5cm}
[نمایی از عملکرد الگوریتم شور \lr{(Shor)} بر شکستن \lr{RSA} و تهدید الگوریتم گرور \lr{(Grover)} بر رمزنگاری متقارن]
\vspace{1.5cm}}}
\caption{نمایی از عملکرد الگوریتم شور \lr{(Shor)} بر شکستن \lr{RSA} و تهدید الگوریتم گرور \lr{(Grover)} بر رمزنگاری متقارن}
\label{fig:2-2}
\end{figure}

%% ──────────────────────────────────────────────────────────────────────────
\section{رمزنگاری پساکوانتومی}
%% ──────────────────────────────────────────────────────────────────────────

رمزنگاری پساکوانتومی \lr{(Post-Quantum Cryptography)} پاسخ جهان رمزنگاری به تهدیدات محاسبات کوانتومی است، اما از یک منظر اساسی با توزیع کلید کوانتومی متمایز است: \lr{PQC} کاملاً بر روی ریاضیات کلاسیک استوار است و نیازی به هیچ زیرساخت اپتیکی یا کوانتومی ندارد. این الگوریتم‌ها بر روی همان پردازنده‌ها و شبکه‌های ارتباطی موجود قابل اجرا هستند، اما پایه ریاضی آن‌ها از مسائلی انتخاب شده که حتی کامپیوترهای کوانتومی توانمند آینده نیز قادر به حل آن‌ها در زمان چندجمله‌ای نیستند \{1\}. این ویژگی، \lr{PQC} را به یک راهکار قابل استقرار فوری و مستقل از فناوری‌های نوظهور تبدیل می‌کند — مزیتی که در محیط‌های عملیاتی محدود نظیر ماهواره‌ها اهمیت حیاتی دارد \{22\}.

فرآیند استانداردسازی مؤسسه ملی استانداردها و فناوری آمریکا \lr{(NIST)} که از سال �2016 آغاز شد و طی آن از میان ۸۲ نامزد دریافتی از سراسر جهان، ۶۹ الگوریتم در سه دور رقابتی فشرده ارزیابی شدند، در اوت ۲۰۲۴ به مرحله نهایی خود رسید \{2\}. در سال ۲۰۲۶، این استانداردها دیگر اسناد آزمایشی نیستند؛ آن‌ها الزامات عملیاتی هستند که سازمان‌ها و دولت‌ها به‌طور فعال در حال مهاجرت به آن‌ها هستند. \textbf{\lr{FIPS 203}} با عنوان \lr{ML-KEM (Module-Lattice-Based Key-Encapsulation Mechanism)}، جایگزین رسمی \lr{ECDH} و \lr{RSA} برای تبادل کلید است و در سه سطح \lr{ML-KEM-512}، \lr{ML-KEM-768} و \lr{ML-KEM-1024} عرضه می‌شود \{3\}. \textbf{\lr{FIPS 204}} با عنوان \lr{ML-DSA (Module-Lattice-Based Digital Signature Algorithm)}، جایگزین \lr{ECDSA} و \lr{RSA-Sign} برای احراز هویت و امضای دیجیتال است و به‌ویژه برای امضای فرمان‌های کنترلی ماهواره و تأیید به‌روزرسانی‌های نرم‌افزاری کاربرد مستقیم دارد \{9\}. \textbf{\lr{FIPS 205}} با عنوان \lr{SLH-DSA (Stateless Hash-Based Digital Signature)}، یک لایه پشتیبان است که امنیت آن صرفاً بر پایه امنیت توابع درهم‌ساز \lr{(Hash Functions)} بنا شده و نه هیچ فرض ریاضی دیگری — این ویژگی آن را به بیمه‌ای بلندمدت در برابر احتمال شکستن الگوریتم‌های مبتنی بر شبکه تبدیل می‌کند \{12\}.

انتخاب رمزنگاری مبتنی بر شبکه مشبک \lr{(Lattice-Based Cryptography)} به‌عنوان پایه هر دو استاندارد اصلی \lr{ML-KEM} و \lr{ML-DSA} بر اساس دو عامل بنیادین بوده است. نخست، امنیت این الگوریتم‌ها بر سختی مسائلی چون یادگیری با خطای ماژولار \lr{(Module Learning With Errors --- MLWE)} استوار است که نه تنها در برابر کامپیوترهای کلاسیک، بلکه در برابر بهترین الگوریتم‌های کوانتومی شناخته‌شده نیز مقاوم است \{3\}. دوم، عملیات محوری این الگوریتم‌ها --- یعنی ضرب ماتریس و بردار در حلقه‌های جبری --- با استفاده از تبدیل نظریه اعداد \lr{(Number Theoretic Transform --- NTT)} به‌شدت بهینه‌پذیر هستند؛ به گونه‌ای که \lr{ML-KEM-768} بر روی پردازنده‌های معمولی در حدود ۰.۰۵ تا ۰.۰۷ میلی‌ثانیه اجرا می‌شود که از \lr{RSA-3072} بیش از سه‌هزار برابر سریع‌تر است \{22\}. برای ماهواره‌هایی که دارای محدودیت‌های توان، حافظه و محاسبات هستند، این کارایی محاسباتی بالا در کنار امنیت اثبات‌پذیر، معماری‌های مبتنی بر شبکه مشبک را به گزینه‌ای اجتناب‌ناپذیر برای استقرار عملی تبدیل می‌کند \{24\}. تأخیر محاسباتی \lr{ML-KEM} --- در حد چند صدم میلی‌ثانیه --- تنها کسر ناچیزی (کمتر از ۰٫۲ درصد) از تأخیر کل برقراری نشست امن (حدود ۳۵ میلی‌ثانیه) را تشکیل می‌دهد و عملاً هیچ گلوگاه \lr{(Bottleneck)} عملیاتی ایجاد نمی‌کند؛ برای مقایسه، تأخیر انتشار یک‌طرفهٔ لینک \lr{LEO} در ارتفاع ۵۰۰ کیلومتر تنها حدود ۱٫۷ میلی‌ثانیه است \{22\}.

\begin{figure}[H]
\centering
\fbox{\parbox{0.78\textwidth}{\centering\vspace{1.5cm}
[طبقه‌بندی درختی خانواده‌های الگوریتم‌های رمزنگاری پساکوانتومی و استانداردهای \lr{NIST FIPS}]
\vspace{1.5cm}}}
\caption{طبقه‌بندی درختی خانواده‌های الگوریتم‌های رمزنگاری پساکوانتومی و استانداردهای \lr{NIST FIPS}}
\label{fig:2-3}
\end{figure}
%% ──────────────────────────────────────────────────────────────────────────
\section{توزیع کلید کوانتومی}
%% ──────────────────────────────────────────────────────────────────────────

توزیع کلید کوانتومی \lr{(Quantum Key Distribution --- QKD)} از منظر مهندسی ارتباطات یک پارادایم کاملاً متفاوت از رمزنگاری کلاسیک است؛ امنیت آن نه از دشواری محاسباتی یک مسئله ریاضی، بلکه از قوانین بنیادین فیزیک کوانتوم نشأت می‌گیرد \{4\}. دو اصل محوری این بنیان را تشکیل می‌دهند. نخست، \textbf{اصل عدم قطعیت هایزنبرگ} \lr{(Heisenberg Uncertainty Principle)} که اندازه‌گیری همزمان دقیق جفت‌های متقارن کوانتومی --- مانند پلاریزاسیون فوتون در دو پایه ناسازگار --- را ممنوع می‌کند. دوم، \textbf{قضیه عدم کپی‌برداری} \lr{(No-Cloning Theorem)} که تکثیر دقیق یک حالت کوانتومی ناشناخته را فیزیکاً غیرممکن می‌سازد \{4\}. ترکیب این دو اصل یک نتیجه مهندسی حیاتی به دنبال دارد: هر تلاشی برای استراق‌سمع بر روی کانال کوانتومی، ناگزیر اثر آماری قابل شناسایی بر روی داده‌های دریافتی به‌جا می‌گذارد. این ویژگی، \lr{QKD} را قادر می‌سازد تا \textbf{امنیت اطلاعاتی-تئوریک} \lr{(Information-Theoretic Security)} را تضمین کند --- نوعی از امنیت که حتی مهاجمی با توان محاسباتی نامحدود نمی‌تواند آن را بشکند.

\textbf{پروتکل \lr{BB84}} که در ۱۹۸۴ توسط بنت و براسار ارائه شد، مبنایی‌ترین و اثبات‌شده‌ترین پروتکل \lr{QKD} است. در این پروتکل، آلیس بیت‌های خود را در قالب قطبش فوتون‌ها \lr{(Photon Polarization)} در یکی از دو پایه اندازه‌گیری متمایز کدگذاری و ارسال می‌کند. باب برای هر فوتون به‌صورت تصادفی یکی از دو پایه را برای اندازه‌گیری انتخاب می‌کند \{6\}. پس از تبادل، از طریق یک کانال کلاسیک عمومی، پایه‌های انتخابی مقایسه و تنها بیت‌هایی که با پایه یکسان اندازه‌گیری شده‌اند نگه داشته می‌شوند --- فرآیندی که «پالایش پایه» \lr{(Basis Sifting)} نام دارد. در پیاده‌سازی‌های عملی ماهواره‌ای مانند \lr{Micius}، این پروتکل در قالب \textbf{\lr{Decoy-State BB84}} اجرا می‌شود که در آن فوتون‌های سیگنال اصلی با فوتون‌های پوشش با شدت متفاوت آمیخته می‌شوند تا حملات تجزیه عدد فوتون \lr{(Photon Number Splitting)} را قابل شناسایی کنند \{6\}. \textbf{پروتکل \lr{E91}} رویکردی متفاوت دارد و از جفت فوتون‌های \textbf{درهم‌تنیده} \lr{(Entangled Photons)} بهره می‌گیرد. در این پروتکل، ماهواره جفت فوتون‌های درهم‌تنیده را به دو ایستگاه زمینی ارسال می‌کند و همبستگی‌های کوانتومی میان آن‌ها پایه تولید کلید مشترک را تشکیل می‌دهد؛ ویژگی کلیدی این معماری در آن است که ماهواره دیگر نیازی به نقش گره قابل اعتماد \lr{(Trusted Node)} ندارد \{7\}.

\begin{figure}[H]
\centering
\fbox{\parbox{0.78\textwidth}{\centering\vspace{1.5cm}
[فلوچارت مراحل پروتکل \lr{BB84} از ارسال فوتون تا استخراج کلید نهایی]
\vspace{1.5cm}}}
\caption{فلوچارت مراحل پروتکل \lr{BB84} از ارسال فوتون تا استخراج کلید نهایی}
\label{fig:2-4}
\end{figure}

از دیدگاه مهندسی ارتباطات، \textbf{نرخ خطای بیت کوانتومی} \lr{(Quantum Bit Error Rate --- QBER)} مهم‌ترین سنجه کیفیت کانال \lr{QKD} است. \lr{QBER} نسبت بیت‌های اشتباه دریافتی به کل بیت‌های دریافتی را بیان می‌کند و ترکیبی از دو منشأ مستقل است: نویز ذاتی کانال (شامل \lr{dark count} آشکارسازها، نویز پس‌زمینه محیط، و خطای تراز اپتیکی) و اثر اندازه‌گیری مهاجم \{5\}. نکته مهندسی اساسی اینجاست که از منظر نظریه اطلاعات، \lr{QBER} یک معیار یکپارچه است که نمی‌توان دو منشأ آن را از یکدیگر تفکیک کرد؛ به همین دلیل، هر افزایش \lr{QBER} --- چه از نویز کانال باشد و چه از استراق‌سمع --- باید محافظه‌کارانه به مهاجم نسبت داده شود. در پروتکل \lr{BB84}، آستانه \lr{QBER} برابر ۱۱ درصد است و هرگاه این آستانه تجاوز شود، هیچ کلید امنی قابل استخراج نیست. در آزمایش‌های عملی ماهواره \lr{Micius}، \lr{QBER} در بهترین شرایط حدود ۱.۱ درصد ثبت شده، اما در یک سیستم واقعی با اتلاف کانال ۲۲ دسیبل، نویز همبستگی‌های تصادفی \lr{(Accidental Coincidences)} با رابطه \lr{R = 2 × S₁ × S₂ × τ} به \lr{QBER} مؤثر می‌افزاید \{5\}.

پس از پالایش پایه و تخمین \lr{QBER}، مرحله \textbf{پس‌پردازش} \lr{(Post-Processing)} آغاز می‌شود که یک خط لوله کامل مهندسی است. در گام نخست، \textbf{تطبیق اطلاعات} \lr{(Information Reconciliation)} از طریق کدهای تصحیح خطا مانند \lr{LDPC} اجرا می‌شود که در آن آلیس اطلاعات افزونه‌ای \lr{(Syndrome)} از طریق کانال کلاسیک عمومی ارسال می‌کند تا باب بتواند خطاهای باقی‌مانده را شناسایی و اصلاح کند \{10\}. مصرف این کانال کلاسیک خود باید توسط کدهای احراز هویت امن --- ترجیحاً مبتنی بر \lr{PQC} --- محافظت شود تا جلوی حمله مرد میانی \lr{(Man-in-the-Middle)} گرفته شود \{10\}. در گام دوم، \textbf{تقویت حریم خصوصی} \lr{(Privacy Amplification)} با اعمال یک تابع هش جهانی \lr{(Universal Hash Function)} مانند ماتریس \lr{Toeplitz}، رشته تصحیح‌شده را فشرده می‌کند تا هرگونه اطلاعاتی که شنودگر احتمالی به‌دست آورده به حد ناچیز از نظر محاسباتی برسد \{8\}. خروجی این خط لوله، کلیدی است با امنیت اطلاعاتی-تئوریک اثبات‌پذیر.

نکته مهندسی کلیدی که اغلب در توصیفات ساده‌شده نادیده گرفته می‌شود، محدودیت ذاتی نرخ تولید کلید \lr{QKD} است. برای یک لینک \lr{LEO} معمولی با اتلاف ۲۲ تا ۳۴ دسیبل، نرخ کلید مخفی به سقف بنیادین \lr{C = -log₂(1-η)} محدود است که برای \lr{η ≈ 0.004} تا \lr{0.01}، در حد چند کیلوبیت تا چند مگابیت در هر پاس قرار می‌گیرد \{4\}. این محدودیت ذاتی، \lr{QKD} را برای تأمین کلیدهای رمزنگاری پیوسته با حجم بالا نامناسب می‌کند و توجیه بنیادین استفاده از معماری ترکیبی با لایه \lr{PQC} را فراهم می‌آورد.

در سال ۲۰۲۶، شبکه‌های \lr{QKD} از لینک‌های ساده نقطه-به-نقطه در حال تبدیل شدن به زیرساخت‌های کلید کوانتومی \lr{(Quantum Key Infrastructure --- QKI)} هستند. آزمایش‌های تجربی نشان داده‌اند که یک شبکه ده‌کاربره با معماری \lr{all-pass} که از گواهی‌های دیجیتال \lr{PQC} به‌جای کلیدهای پیش‌توزیع‌شده استفاده می‌کند، می‌تواند نرخ‌های کلید ۷۰ تا ۱۴۰ کیلوبیت در ثانیه را در فواصل ۵۰ تا ۹۰ کیلومتر با پایداری ۳۰ ساعته بدون قطعی حفظ کند \{10\}. در مقیاس قاره‌ای، این زیرساخت‌ها با استفاده از ماهواره‌های \lr{LEO} به‌عنوان \lr{relay} --- که در هر پاس ۴.۵ مگابیت کلید مخفی در یک پاس واحد بلادرنگ تولید می‌کنند \{8\} --- به یک معماری جامع امنیت کوانتومی تبدیل می‌شوند که در آن \lr{PQC} و \lr{QKD} نه رقیب، بلکه مکمل یکدیگرند.

%% ──────────────────────────────────────────────────────────────────────────
\section{ارتباطات ماهواره‌ای و \lr{QKD} ماهواره‌ای}
%% ──────────────────────────────────────────────────────────────────────────

درک معماری ارتباطات ماهواره‌ای از منظر مهندسی، پیش‌نیاز اساسی برای ارزیابی چالش‌های واقعی استقرار \lr{QKD} در فضا است. ماهواره‌ها بر اساس ارتفاع مداری به سه دسته اصلی تقسیم می‌شوند که هر یک مشخصات فیزیکی و تجاری متمایزی دارند. مدار پایین زمین \lr{(Low Earth Orbit --- LEO)} در ارتفاع ۲۰۰ تا ۲۰۰۰ کیلومتر قرار می‌گیرد، مدار میانی \lr{(Medium Earth Orbit --- MEO)} در ارتفاع ۲۰۰۰ تا ۳۵۷۸۶ کیلومتر، و مدار ژئوستاتیک \lr{(Geostationary Orbit --- GEO)} در ارتفاع ثابت ۳۵۷۸۶ کیلومتر \{23\}. از دیدگاه تحلیل بودجه لینک \lr{(Link Budget)}، افت مسیر فضای آزاد با مجذور فاصله مقیاس می‌شود و این رابطه تعیین‌کننده مطلق کارایی لینک کوانتومی است. یک ماهواره \lr{GEO} با فاصله ۳۵۷۸۶ کیلومتر، اتلاف انتشاری بیش از ۶۴ دسیبل را برای یک لینک اپتیکی نمونه تحمیل می‌کند؛ در حالی که همان لینک برای یک ماهواره \lr{LEO} در ارتفاع ۵۰۰ کیلومتر تنها حدود ۲۰ تا ۳۰ دسیبل اتلاف انتشاری دارد \{5\}. از آنجا که پروتکل‌های \lr{DV-QKD} مانند \lr{BB84} تنها تا اتلاف ۶۰ تا ۷۵ دسیبل قابل اجرا هستند، \lr{LEO} تنها گزینه عملیاتی فعلی برای \lr{QKD} ماهواره‌ای است \{23\}.

تفاوت بنیادین دیگری که \lr{QKD} ماهواره‌ای را از ارتباطات ماهواره‌ای کلاسیک جدا می‌کند، ماهیت کانال فیزیکی مورد نیاز است. ارتباطات رادیویی \lr{(Radio Frequency --- RF)} که در اکثر سیستم‌های ماهواره‌ای تجاری برای تله‌متری، فرمان و انتقال داده به‌کار می‌روند، بر امواج الکترومغناطیسی با پراکندگی همه‌سویه استوارند \{22\}. \lr{QKD} از اساس با این رویکرد ناسازگار است؛ ارسال کوانتوم‌های اطلاعات (فوتون‌های منفرد یا درهم‌تنیده) به ارتباطات اپتیکی فضای آزاد \lr{(Free-Space Optics --- FSO)} نیاز دارد که در آن پرتو لیزر به‌شکل مستقیم و باریک میان فرستنده و گیرنده منتشر می‌شود.

اولین و مهم‌ترین چالش مهندسی، سیستم نشانه‌روی، رهیابی و ردیابی \lr{(Pointing, Acquisition and Tracking --- PAT)} است. یک پرتو لیزر \lr{QKD} در طول موج ۷۸۰ تا ۸۵۰ نانومتر از یک تلسکوپ ۳۰ سانتی‌متری (مشخصاتی نمونه‌وار از پایانه‌های کلاس \lr{Micius}، متمایز از پایانهٔ فشردهٔ ۱۰ سانتی‌متری مرجع این پژوهش در فصل سوم)، واگرایی اپتیکی در حدود ۱۰ میکرورادیان دارد. این به معنی آن است که در فاصله ۵۰۰ کیلومتر، قطر پرتو به تنها ۵ تا ۱۲ متر می‌رسد --- دقتی که سیستم‌های \lr{PAT} باید آن را در برابر حرکت سریع ماهواره با سرعت ۷ کیلومتر در ثانیه حفظ کنند \{6\}. ماهواره \lr{Micius} چین در این آزمایش‌ها دقت نشانه‌روی نهایی حدود ۰.۶ میکرورادیان را به دست آورد. یک نکته مهندسی حیاتی اینجاست که سیستم \lr{PAT} باید زاویه پیش‌رو \lr{(Point Ahead Angle --- PAA)} را جبران کند --- چرا که در زمان رفت‌وبرگشت نور، ماهواره \lr{LEO} با سرعت ۷ کیلومتر در ثانیه ۲۰ تا ۵۰ میکرورادیان جابجا می‌شود. از آنجا که این مقدار از زاویه ایزوپلاناتیک جو \lr{(Isoplanatic Angle)} بزرگ‌تر است، این دلیل اصلی برتری آشکار معماری \lr{Downlink} (از ماهواره به زمین) بر \lr{Uplink} در سیستم‌های \lr{QKD} ماهواره‌ای است \{23\}.

دومین چالش اساسی، تلاطم جوی \lr{(Atmospheric Turbulence)} است که در قالب پدیده‌های \lr{Scintillation}، \lr{Beam Wandering}، و \lr{Beam Broadening} بر عملکرد لینک اثر می‌گذارد. از منظر بودجه لینک، این تلاطم یک اتلاف مؤثر حدود ۲ دسیبل به بودجه اضافه می‌کند \{8\}. یک تحلیل دقیق بودجه لینک برای یک سیستم نمونه --- با تلسکوپ فرستنده ۱۰ سانتی‌متری روی ماهواره، تلسکوپ گیرنده ۶۰ سانتی‌متری روی زمین، و ارتفاع مدار ۵۰۰ کیلومتر --- اتلاف کل را در حدود ۳۱ دسیبل برای یک پاس بهینه با زاویه \lr{elevation} نود درجه نشان می‌دهد که در آن ۲ دسیبل به تلاطم، ۱.۲۳ دسیبل به اپتیک گیرنده، ۲.۳۶ دسیبل به آشکارساز، و ۲.۵۱ دسیبل به کوپلینگ فیبر تعلق می‌گیرد \{8\}. همچنین، نور پس‌زمینه --- که در ارتباطات \lr{RF} اصلاً معنایی ندارد --- در \lr{FSO} یک منبع جدی نویز است که عملیات اکثر سیستم‌های \lr{QKD} ماهواره‌ای را به ساعات شبانه محدود می‌کند.

محدودیت سوم ذاتاً با هندسه مداری \lr{LEO} مرتبط است. هر ماهواره \lr{LEO} بر فراز یک ایستگاه زمینی تنها پنجره‌ای ۵ تا ۱۰ دقیقه‌ای در هر پاس دارد --- و تنها حدود ۱۵ پاس در روز در بهترین حالت برای یک ایستگاه خاص \{23\}. این محدودیت مستقیماً نرخ انباشت کلید روزانه را محدود می‌کند؛ به طوری که حتی با بهترین سیستم‌های موجود و در یک پاس کامل ۲۹۴ ثانیه‌ای، تنها ۴.۵۸ مگابیت کلید مخفی تولید می‌شود \{8\}. این رقم، در مقایسه با نیاز ارتباطات ماهواره‌ای پیوسته با پهنای باند گیگابیت، نشان می‌دهد که \lr{QKD} ماهواره‌ای به‌تنهایی قادر به تأمین کل نیاز رمزنگاری سیستم نیست و معماری ترکیبی با لایه \lr{PQC} یک ضرورت مهندسی اجتناب‌ناپذیر است.

در سال ۲۰۲۶، اکوسیستم \lr{QKD} ماهواره‌ای در حال گذار از آزمایش‌های علمی به سرویس‌های عملیاتی است. از ماهواره \lr{Micius} چین به‌عنوان اولین اثبات عملی، به مأموریت‌هایی مانند \lr{Eagle-1} اروپا، \lr{SAGA 1G} در چارچوب \lr{EuroQCI}، و \lr{IRIS²} در قالب زیرساخت ماهواره‌ای امن اتحادیه اروپا رسیده‌ایم \{23\}. این تحول زیرساختی --- از لینک‌های نقطه-به-نقطه به شبکه‌های زیرساخت کلید کوانتومی \lr{(QKI)} --- همراه با ادغام لایه \lr{PQC} برای مدیریت نشست‌های ارتباطی پیوسته، چارچوب معماری‌ای را شکل می‌دهد که پایان‌نامه حاضر در آن فعالیت می‌کند.

\begin{figure}[H]
\centering
\fbox{\parbox{0.78\textwidth}{\centering\vspace{1.5cm}
[معماری سیستم \lr{QKD} ماهواره‌ای در مدار \lr{LEO} و مقایسه پیکربندی‌های \lr{Downlink} و \lr{Trusted-Node Relay}]
\vspace{1.5cm}}}
\caption{معماری سیستم \lr{QKD} ماهواره‌ای در مدار \lr{LEO} و مقایسه پیکربندی‌های \lr{Downlink} و \lr{Trusted-Node Relay}}
\label{fig:2-5}
\end{figure}
%% ──────────────────────────────────────────────────────────────────────────
\section{مرور سیستماتیک پژوهش‌های پیشین}
%% ──────────────────────────────────────────────────────────────────────────

بررسی انتقادی ادبیات پژوهشی در حوزه تقاطع \lr{QKD}، \lr{PQC} و ارتباطات ماهواره‌ای نشان می‌دهد که این جریان‌های پژوهشی از لحاظ تاریخی به‌صورت مستقل حرکت کرده و تنها در سال‌های اخیر به‌سمت یکدیگر همگرا شده‌اند. این بخش، پژوهش‌های موجود را در سه دسته اصلی ارزیابی می‌کند.

\textbf{دسته اول: پژوهش‌های ماهواره‌ای \lr{QKD}} بیشترین پیشرفت تجربی را داشته‌اند. نقطه عطف این حوزه، ماهواره \lr{Micius} چین بود که اولین \lr{QKD} ماهواره-زمین را در مقیاس ۱۲۰۰ کیلومتر با پروتکل \lr{Decoy-State BB84} به اثبات رساند و ۳۰۰۹۳۹ بیت کلید امن در یک پاس ۲۷۳ ثانیه‌ای تولید کرد \{6\}. این کار پیشرفت ساختاری بنیادینی بود. با این حال، وابستگی معماری \lr{BB84} به ماهواره به‌عنوان گره قابل اعتماد \lr{(Trusted Node)} نقطه ضعف امنیتی ذاتی محسوب می‌شد. \lr{Yin} و همکاران یک شبکه بین‌قاره‌ای چین-اتریش گسترش دادند و اولین کنفرانس ویدیویی امن را در فاصله ۷۶۰۰ کیلومتری برگزار کردند \{11\}، اما این پیشرفت هنوز به همان آسیب‌پذیری گره قابل اعتماد متکی بود. گامی بزرگ‌تر در سال ۲۰۲۰ برداشته شد که با پروتکل \lr{BBM92} نشان دادند حتی اگر ماهواره به خطر بیفتد، امنیت کلید بین دو ایستگاه زمینی حفظ می‌شود \{7\}. پژوهش‌های بلادرنگ \lr{(Real-time)} نیز پروتکل کامل \lr{QKD} --- شامل تصحیح خطا و تقویت حریم خصوصی --- را برای اولین بار در یک پاس واحد به‌صورت بلادرنگ اجرا کردند \{8\}. ارزیابی سیستماتیک مرکز فضایی آلمان \lr{(DLR)} توسط \lr{Orsucci} و همکاران در ۲۰۲۵ نیز یک چارچوب تصمیم‌گیری هفت‌مرحله‌ای ارائه داد که نشان می‌دهد برای مأموریت‌های نزدیک آینده، معماری \lr{DV-QKD} در لینک دانلود با پروتکل \lr{BB84} در \lr{LEO} بهینه‌ترین انتخاب است \{23\}. جمع‌بندی نقادانه این دسته آن است که همه این پژوهش‌ها بر بهینه‌سازی لایه فیزیکی و کوانتومی متمرکز بوده‌اند و مسئله امنیت کانال کلاسیک \lr{QKD} در برابر حملات کوانتومی را اغلب به‌عنوان یک فرض سطحی درنظر گرفته‌اند.

\textbf{دسته دوم: پژوهش‌های پیاده‌سازی \lr{PQC}} عمدتاً بر سربار محاسباتی، مصرف حافظه و کارایی الگوریتم‌ها در محیط‌های محدود تمرکز داشته‌اند. انتشار استانداردهای \lr{NIST FIPS 203}، \lr{FIPS 204} و \lr{FIPS 205} \{3, 9, 12\} مرجع پایه این حوزه را تشکیل می‌دهد. در حوزه ارزیابی پیاده‌سازی، پژوهش‌های معیار برای کاربرد ماهواره‌ای نشان دادند که تأخیر محاسباتی \lr{ML-KEM} در مقایسه با تأخیر کل برقراری نشست امن (حدود ۳۵ میلی‌ثانیه) کمتر از ۰٫۲ درصد است --- یعنی هیچ گلوگاه عملیاتی ایجاد نمی‌کند \{22\}. \lr{Dharminder} و همکاران آسیب‌پذیری‌های پروتکل‌های احراز هویت ماهواره \lr{LEO} مبتنی بر \lr{RSA/ECC} را در برابر حملات نشت سیگنال اثبات کردند و یک پروتکل جایگزین مبتنی بر \lr{Ring-LWE} طراحی نمودند \{18\}. نقد اساسی این دسته از پژوهش‌ها آن است که تمرکز آن‌ها بر ارزیابی مستقل الگوریتم‌های \lr{PQC} بوده و این سؤال که \lr{PQC} به‌تنهایی چه نقاط کوری در برابر مهاجمی با \lr{CRQC} دارد کمتر پرسیده شده است.

\textbf{دسته سوم: پژوهش‌های معماری‌های ترکیبی} جدیدترین و از منظر این پایان‌نامه مهم‌ترین دسته هستند. این جریان پژوهشی از این درک بنیادین آغاز می‌شود که \lr{QKD} و \lr{PQC} نه رقیب، بلکه مکمل یکدیگرند: \lr{QKD} امنیت اطلاعاتی-تئوریک ارائه می‌دهد اما نرخ کلید محدود دارد و به زیرساخت اپتیکی پیچیده نیاز دارد؛ \lr{PQC} روی سخت‌افزار موجود اجرا می‌شود اما امنیت محاسباتی دارد. \lr{Wang} و همکاران در ۲۰۲۱ اولین کار تجربی ترکیبی مهم را ارائه دادند که در آن از \lr{PQC} برای احراز هویت کانال کلاسیک \lr{QKD} استفاده شد و ثابت کردند این رویکرد بدون هیچ تأثیری بر نرخ کلید، نیاز به کلیدهای پیش‌توزیع‌شده را حذف می‌کند \{10\}. \lr{Rani} و همکاران سیستمی را گزارش دادند که \lr{QKD} مبتنی بر درهم‌تنیدگی را با \lr{AES-256} در معماری متوالی \lr{(Sequential)} برای کانال‌های ماهواره-زمین ترکیب می‌کند \{13\}. این کار نشان داد که در اتلاف ۱۰.۳ دسیبل، سیستم ترکیبی بدون هیچ کاهش نرخ کلیدی، امنیت مضاعف در برابر شکست احتمالی هر یک از دو لایه را فراهم می‌کند. \lr{Zeng} و همکاران سه معماری هیبرید \lr{PQC-QKD} --- اتصال سری، موازی، و تسهیم راز \lr{(Secret Sharing)} --- را از منظر ریاضی تحلیل کردند \{19\}. نقد تحلیلی این دسته آن است که اکثر این پژوهش‌ها بر ترکیب \lr{QKD} زمینی (فیبری) با \lr{PQC} تمرکز داشته‌اند، و معماری‌های ترکیبی اختصاصی برای بستر ماهواره‌ای --- با محدودیت‌های منحصربه‌فرد آن --- شکافی است که پایان‌نامه حاضر به‌طور مستقیم در پی پر کردن آن است.
%% ──────────────────────────────────────────────────────────────────────────
\section{شکاف پژوهشی}
%% ──────────────────────────────────────────────────────────────────────────

مرور نقادانه ادبیات موجود که در بخش پیشین ارائه شد، تصویر روشنی از سه جریان پژوهشی نسبتاً موازی ترسیم می‌کند که هر یک در قلمرو خود به پیشرفت‌های چشمگیری دست یافته، اما تلاقی هدفمند و جامع این سه جریان --- به‌ویژه در بستر مهندسی واقعی ماهواره‌های مدار پایین زمین --- همچنان یک فضای پژوهشی باز و اثبات‌شده است.

\textbf{بُعد اول شکاف، ماهیت محیط عملیاتی هدف است.} بیشینه پژوهش‌های ترکیبی \lr{QKD} و \lr{PQC} که تاکنون منتشر شده‌اند، بر شبکه‌های فیبری زمینی یا آزمایش‌های آزمایشگاهی با فرض کانال‌های با ویژگی‌های ثابت و قابل کنترل استوار بوده‌اند. محیط عملیاتی یک ماهواره \lr{LEO} اما از اساس متفاوت است: کانال ارتباطی \lr{Free Space Optical} می‌تواند در طول یک پاس واحد ۱۵ تا ۳۰ دسیبل در اتلاف تغییر کند \{8\}، و پنجره عملیاتی مفید آن برای هر ایستگاه زمینی تنها ۵ تا ۱۰ دقیقه در هر گذر است \{23\}. این متغیرهای دینامیکی در ادبیات موجود به‌طور کافی پرداخته نشده‌اند.

\textbf{بُعد دوم شکاف، مسئله محدودیت‌های اندازه، وزن و توان \lr{(SWaP)}} ماهواره است. یک ماهواره \lr{LEO} با محدودیت‌های جدی در توان مصرفی (معمولاً زیر ۵۰ وات)، حافظه (چند مگابایت)، و سرعت پردازنده (محدوده ده‌ها تا صدها مگاهرتز) مواجه است \{5\}. یک معماری ترکیبی \lr{QKD} و \lr{PQC} بار محاسباتی اضافه‌تری ایجاد می‌کند که صرفاً مجموع بارهای جداگانه نیست. تداخل این دو لایه --- به‌ویژه در مرحله پس‌پردازش \lr{QKD} که با تطبیق اطلاعات \lr{(Information Reconciliation)} و تقویت حریم خصوصی \lr{(Privacy Amplification)} حجم بالایی از پردازش را می‌طلبد، و همزمانی آن با احراز هویت \lr{PQC} --- در ادبیات موجود به‌طور جامع مدل‌سازی نشده است.

\textbf{بُعد سوم شکاف، فقدان چارچوب ارزیابی یکپارچه است.} در پژوهش‌های ترکیبی موجود، هر گروه تحقیقاتی معیارهای ارزیابی متفاوتی را انتخاب کرده است؛ برخی بر حذف نیاز به کلید پیش‌توزیع‌شده تمرکز کرده‌اند، برخی بر اثبات تجربی پیاده‌سازی متوالی، و برخی دیگر بستر ماهواره‌ای را به‌عنوان سناریوی هدف مدنظر قرار نداده‌اند \{19\}. این عدم‌انسجام در معیارهای ارزیابی، مقایسه معنادار میان رویکردها را دشوار می‌کند و تصمیم‌گیری مهندسی برای انتخاب بهترین معماری برای یک کاربرد ماهواره‌ای واقعی را عملاً غیرممکن می‌سازد.

\textbf{بُعد چهارم شکاف، مسئله مدیریت چرخه عمر کلید \lr{(Key Lifecycle)}} در سناریوی ماهواره‌ای با عمر ۱۰ تا ۱۵ ساله است \{22\}. در این بازه، چندین تحول کیفی محتمل است: ممکن است \lr{Q-Day} واقع شود؛ ممکن است الگوریتم‌های \lr{PQC} در فرآیندهای آتی شکسته شوند؛ یا ممکن است فناوری سیستم‌های \lr{PAT} و منابع فوتونی ماهواره‌ای به شکلی پیشرفت کنند که نرخ کلید \lr{QKD} مرتبه بزرگی افزایش یابد. هیچ‌یک از معماری‌های پیشنهادی در ادبیات موجود، یک منطق \lr{Fallback} بهینه برای انتقال کنترل‌شده میان این سناریوها ارائه نداده‌اند.

با توجه به این چهار بُعد تعریف‌شده، این پایان‌نامه قصد دارد دقیقاً این جای خالی را پر کند. هدف اصلی آن طراحی و ارزیابی یک معماری ترکیبی لایه‌بندی‌شده است که در آن لایه \lr{ML-KEM} برای تبادل کلید جلسه، لایه \lr{ML-DSA} برای احراز هویت کانال کلاسیک \lr{QKD}، و لایه فیزیکی \lr{QKD} برای تولید کلیدهای اصلی \lr{(Master Keys)} با امنیت اطلاعاتی-تئوریک، به‌شکل هدفمند و با درنظر گرفتن محدودیت‌های منحصربه‌فرد \lr{LEO} ادغام شده‌اند. این معماری با صراحت برای سناریوی اتلاف متغیر کانال \lr{FSO} در مدار \lr{LEO} مدل‌سازی می‌شود، منطق پشتیبان برای شرایطی که \lr{QBER} از آستانه بحرانی تجاوز کند تعریف می‌شود، و سربار محاسباتی کل ترکیب در یک محیط شبیه‌سازی پایتون در برابر الگوریتم‌های کلاسیک ارزیابی می‌شود.

\begin{figure}[H]
\centering
\fbox{\parbox{0.78\textwidth}{\centering\vspace{1.5cm}
[نمای مفهومی معماری ترکیبی لایه‌بندی‌شده \lr{PQC+QKD} برای ارتباطات ایمن ماهواره‌ای]
\vspace{1.5cm}}}
\caption{نمای مفهومی معماری ترکیبی لایه‌بندی‌شده \lr{PQC+QKD} برای ارتباطات ایمن ماهواره‌ای}
\label{fig:2-6}
\end{figure}

%% ──────────────────────────────────────────────────────────────────────────
\section{جمع‌بندی}
%% ──────────────────────────────────────────────────────────────────────────

فصل دوم این پایان‌نامه یک مسیر تحلیلی منسجم را از پایه‌های رمزنگاری مدرن تا مرزهای پژوهشی حوزه امنیت کوانتومی ماهواره‌ای طی کرد. درک این مسیر در کلیت آن، پیش‌نیاز ضروری برای ارزیابی معماری پیشنهادی در فصول بعدی است.

نقطه آغاز این فصل، تشریح زیرساختی بود که دهه‌ها امنیت ارتباطات دیجیتال جهان را تضمین کرده است. الگوریتم‌های \lr{RSA} و \lr{ECC}، با پایه‌های ریاضی دشوار خود در تجزیه اعداد صحیح بزرگ و لگاریتم گسسته روی منحنی‌های بیضوی، برای دهه‌ها ستون‌های اصلی تبادل کلید --- از \lr{TLS} و \lr{SSH} تا پروتکل‌های کنترل ماهواره --- بوده‌اند. \lr{AES} نقش رمزنگاری حجم بالای داده را با کارایی بسیار بالا ایفا کرده است. این معماری دوگانه --- نامتقارن برای تبادل کلید، متقارن برای رمزنگاری داده --- یک چارچوب کارآمد و بلوغ‌یافته است که در آن ضعیف‌ترین حلقه همیشه الگوریتم‌های نامتقارن بوده‌اند، نه \lr{AES}.

سپس این فصل با صراحت نشان داد که این حلقه ضعیف، در سال ۲۰۲۶، دیگر یک تهدید نظری نیست بلکه یک بحران در حال وقوع است. الگوریتم شور با پیچیدگی چندجمله‌ای، \lr{RSA} و \lr{ECC} را به‌طور کامل می‌شکند؛ و آنچه این تهدید را از یک مسئله آینده به یک اضطرار امروزی تبدیل می‌کند، استراتژی \lr{HNDL} است. مهاجمان هم‌اکنون ترافیک رمزشده ارتباطات ماهواره‌ای را ضبط می‌کنند. ماهواره‌هایی که امروز در مدار هستند و با \lr{RSA/ECC} محافظت می‌شوند، با عمر عملیاتی ۱۰ تا ۱۵ ساله، احتمالاً پیش از پایان مأموریتشان در بازه زمانی وقوع \lr{Q-Day} قرار خواهند گرفت. این دیگر یک محاسبه احتمالاتی مبهم نیست --- این یک ریسک مهندسی قابل اندازه‌گیری است.

در پاسخ به این بحران، دو جریان پژوهشی موازی و مستقل شکل گرفته‌اند. جریان اول، رمزنگاری پساکوانتومی، با بهره‌گیری از ریاضیات کلاسیک اما مقاوم در برابر الگوریتم‌های کوانتومی --- به‌ویژه مسائل شبکه مشبک مانند \lr{MLWE} --- فرآیند هشت‌ساله استانداردسازی \lr{NIST} که در اوت ۲۰۲۴ با انتشار \lr{FIPS 203}، \lr{FIPS 204} و \lr{FIPS 205} به مرحله نهایی خود رسید، جهان را با یک مجموعه ابزار آماده برای مهاجرت تجهیز کرده است. جریان دوم، توزیع کلید کوانتومی، پارادیم متفاوتی دارد. \lr{QKD} امنیتی ارائه می‌دهد که نه به دشواری یک مسئله ریاضی، بلکه به قوانین فیزیک وابسته است. دستاوردهای تجربی ماهواره \lr{Micius}، شبکه بین‌قاره‌ای چین-اتریش، و آزمایش‌های بلادرنگ نشان داده‌اند که \lr{QKD} ماهواره‌ای از محدوده آزمایشگاه به سمت کاربردهای عملی در حرکت است؛ اما محدودیت نرخ کلید در هر پاس \lr{LEO}، نیاز به سیستم‌های \lr{PAT} پیچیده، تأثیر تشعشع بر آشکارسازها، و وابستگی معماری‌های \lr{Trusted-Relay} به ایمنی فیزیکی ماهواره، همگی نقاط آسیب‌پذیر عملیاتی هستند که نمی‌توان آن‌ها را نادیده گرفت.

نکته‌ای که در این فصل پدیدار شد و اهمیت آن از هر بحث تکنیکی جزئی بیشتر است، یک دیدگاه فلسفی-مهندسی است: در طراحی سیستم‌های امنیتی حیاتی، هدف نهایی نباید «بهترین الگوریتم را بیابید» باشد، بلکه «اطمینان حاصل کنید که کل سیستم همچنان ایمن بماند» باشد. این اصل --- که در مهندسی سیستم‌های حیاتی به \lr{Fail-Safe Design} معروف است --- چیزی است که هر یک از جریان‌های \lr{PQC} و \lr{QKD} به‌تنهایی نمی‌توانند ارائه دهند، اما ترکیب آن‌ها در یک معماری لایه‌بندی‌شده می‌تواند.

در یک معماری ترکیبی هوشمندانه، اگر لایه \lr{QKD} به دلیل افزایش \lr{QBER} ناشی از تلاطم جوی از دست برود، لایه \lr{ML-KEM} به‌عنوان پشتیبان \lr{(Fallback)} وارد عمل می‌شود؛ اگر \lr{ML-KEM} در آینده‌ای دور توسط یک الگوریتم جدید شکسته شود، کلیدهایی که توسط \lr{QKD} توزیع شده‌اند با امنیت اطلاعاتی-تئوریک همچنان ایمن می‌مانند؛ و اگر کانال کلاسیک \lr{QKD} مورد حمله قرار گیرد، امضای \lr{ML-DSA} آن را بلافاصله شناسایی می‌کند. این چندلایه‌گی امنیتی --- که در آن شکست هیچ یک از اجزا منجر به شکست سیستم نمی‌شود --- دقیقاً همان چیزی است که یک ماهواره با عمر عملیاتی ۱۵ ساله در محیطی با تهدیدات در حال تکامل به آن نیاز دارد.

بر این پایه، فصل سوم این پایان‌نامه معماری لایه‌بندی‌شده پیشنهادی را که بر اساس همین منطق طراحی شده با تفصیل ارائه می‌کند: چگونگی ادغام \lr{ML-KEM}، \lr{ML-DSA} و لایه فیزیکی \lr{QKD} در یک ساختار منسجم با منطق پشتیبان صریح، توپولوژی سیستم در بستر لینک \lr{LEO}، و چارچوب شبیه‌سازی پایتونی که سربار عملیاتی کل این معماری را در برابر راهکارهای کلاسیک به‌صورت کمی ارزیابی می‌کند.

%% — پایان فصل دوم —
%%%%%%%%%%%%%%%%%%%%%%%%%%%%%%%%%%%%%%%%%%%%%%%%%%%%%%%%%%%%%%%%%%%%%%%%%%%%%
%%  chapter3.tex — فصل سوم: روش اجرای تحقیق
%%  paste کنید در main.tex بعد از پایان فصل دوم، قبل از \end{document}
%%%%%%%%%%%%%%%%%%%%%%%%%%%%%%%%%%%%%%%%%%%%%%%%%%%%%%%%%%%%%%%%%%%%%%%%%%%%%

\chapter{روش اجرای تحقیق}
\setcounter{figure}{0}
\setcounter{table}{0}
\setcounter{equation}{0}

%% ──────────────────────────────────────────────────────────────────────────
\section{مقدمه}
%% ──────────────────────────────────────────────────────────────────────────

فصل دوم این پژوهش با مرور سیستماتیک ادبیات موجود نشان داد که تهدید ناشی از پیشرفت محاسبات کوانتومی \lr{(Quantum Computing)} نه یک چالش نظری آینده‌نگر، بلکه یک واقعیت عملیاتی جاری است که زیرساخت‌های رمزنگاری کنونی را در معرض خطر جدی قرار داده است. بررسی دقیق مرزهای دانش موجود در آن فصل نشان داد که رویکردهای منفرد --- اعم از توزیع کلید کوانتومی \lr{(Quantum Key Distribution --- QKD)} به‌تنهایی، یا رمزنگاری پسا-کوانتومی \lr{(Post-Quantum Cryptography --- PQC)} به‌تنهایی --- برای تأمین امنیت جامع در محیط‌های عملیاتی دشوار، به‌ویژه کانال‌های ماهواره‌ای در مدار پایین زمین \lr{(Low Earth Orbit --- LEO)}، کافی نیستند. \lr{QKD} هرچند امنیت اطلاعاتی-تئوریک \lr{(Information-Theoretic Security)} اثبات‌پذیر ارائه می‌دهد، اما ظرفیت بنیادی کانال فضایی با اتلاف‌های معمول بالای ۲۰ دسیبل، نرخ کلید قابل تولید را به مرتبه‌ای محدود می‌کند که نمی‌تواند به‌تنهایی بار ارتباطات پرحجم را تحمل کند \{4\}. از سوی دیگر، \lr{PQC} و به‌خصوص استانداردهای نهایی‌شده مؤسسه ملی استانداردها و فناوری ایالات متحده \lr{(NIST)} شامل \lr{FIPS 203 (ML-KEM)} و \lr{FIPS 204 (ML-DSA)} که اکنون در سال ۲۰۲۶ کاملاً عملیاتی هستند، سرعت، انعطاف‌پذیری و استقلال از زیرساخت سخت‌افزاری تخصصی را فراهم می‌آورند، اما بر سختی محاسباتی مسائل ریاضی استوار هستند که هنوز یک اثبات امنیتی محکم \lr{(Tight Security Proof)} در مدل اوراکل تصادفی کوانتومی \lr{(Quantum Random Oracle Model --- QROM)} ندارند \{21\}.

این تحلیل شکاف پژوهشی به یک نتیجه‌گیری راهبردی منتهی می‌شود: در محیط‌های با حساسیت عملیاتی بالا مانند ارتباطات ماهواره‌ای، هیچ‌کدام از این دو نظام رمزنگاری به‌طور مجزا نمی‌توانند الگوی دفاعی کافی باشند. آنچه پژوهش حاضر بر آن استوار است، اتخاذ رویکرد ادغام لایه‌بندی‌شده \lr{(Layered Integration)} است؛ رویکردی که در آن \lr{QKD} و \lr{PQC} نه به‌عنوان رقیب، بلکه به‌عنوان مکمل یکدیگر عمل می‌کنند تا ضعف‌های هر یک را با قدرت دیگری جبران کنند. این منطق دوگانه را \lr{Rani} و همکاران به صراحت در اولین پیاده‌سازی تجربی ترکیب \lr{QKD} با لایه \lr{PQC} برای کانال‌های ماهواره-زمین تأیید کردند: «اگر \lr{QKD} به دلیل ناکاملی فیزیکی دستگاه‌ها به خطر بیفتد، لایه \lr{PQC} امنیت محاسباتی را حفظ می‌کند، و اگر الگوریتم \lr{PQC} در آینده به چالش کشیده شود، امنیت اطلاعاتی-تئوریک \lr{QKD} دست‌نخورده باقی می‌ماند» \{13\}. این اصل، ستون فقرات معماری پیشنهادی فصل حاضر است.

هدف اصلی این فصل، حرکتی مشخص و هدفمند از مباحث نظری به سمت یک طراحی سیستمی کاربردی و مستدل است. در حالی که فصل دوم به تبیین «چرایی» لزوم این تحول پرداخت، فصل حاضر به «چگونگی» آن می‌پردازد؛ یعنی ارائه یک معماری امنیتی لایه‌بندی‌شده، مقاوم در برابر خطا \lr{(Fault-Resilient)} و بهینه‌شده برای محدودیت‌های عملیاتی ماهواره‌های \lr{LEO}. این محدودیت‌ها که در ادبیات تخصصی با کنار هم گذاشتن سه ملاحظه توان \lr{(Power)}، وزن \lr{(Weight)} و اندازه \lr{(Size)} زیر عنوان مشترک \lr{SWaP} شناخته می‌شوند، فضای طراحی را به نحو چشمگیری تنگ می‌کنند \{24\}. یک ماهواره \lr{CubeSat} یا میکروماهواره در مدار \lr{LEO} نه توان محاسباتی نامحدود دارد، نه حافظه بی‌پایان، و نه امکان تعمیر یا جایگزینی در صورت خرابی. هر الگوریتم رمزنگاری که روی این سکو اجرا می‌شود باید در برابر چارچوب \lr{SWaP} ارزیابی شده و معماری کلی سیستم باید قادر باشد حتی در صورت از کار افتادن یک لایه، عملکرد ایمن را حفظ کند --- اصلی که در طراحی مهندسی با عنوان طراحی ایمن در برابر خرابی \lr{(Fail-Safe Design)} شناخته می‌شود.

نوآوری اصلی این پژوهش در تلفیق سه عنصر است که تاکنون در یک معماری یکپارچه و مخصوص \lr{LEO} با هم نیامده‌اند: اول، استفاده از \lr{ML-KEM} به‌عنوان مکانیزم برپایی کلید \lr{(Key Encapsulation Mechanism)} در سه سطح امنیتی \lr{NIST} برای ارزیابی تطبیق‌پذیری با محدودیت‌های کانال؛ دوم، به‌کارگیری \lr{ML-DSA} برای احراز هویت پسا-کوانتومی \lr{(Post-Quantum Authentication)} پیام‌های کنترلی و کانال کلاسیک \lr{QKD} به‌گونه‌ای که حملات مرد میانی \lr{(Man-in-the-Middle)} حتی در برابر یک مهاجم مجهز به کامپیوتر کوانتومی ناکارآمد باشند \{10\}؛ و سوم، ادغام \lr{QKD} مبتنی بر پروتکل \lr{BB84} با حالت‌های فریب \lr{(Decoy-State Protocol)} که اثبات‌پذیرترین پروتکل برای کانال فضایی است \{23\}، به‌عنوان لایه تولید کلید اصلی \lr{(Master Key)} در چارچوب یک سیستم چندلایه با منطق \lr{Fallback} تعریف‌شده.

اهمیت این رویکرد از منظر مهندسی نرم‌افزار نیز قابل توجه است. پنجره دید محدود ماهواره‌های \lr{LEO} --- که معمولاً بین ۵ تا ۱۰ دقیقه است --- تمامی مراحل پروتکل را زیر فشار زمانی قرار می‌دهد و ضرورت طراحی الگوریتم‌هایی با سربار محاسباتی کم را دوچندان می‌کند \{8\}. برآورد بودجه لینک \lr{(Link Budget)} متغیر در طول گذر ماهواره، که در آن اتلاف از مقادیر کمینه در گذر از بالاترین ارتفاع تا مقادیر بیشینه در نزدیکی خط افق تغییر می‌کند، متغیری است که هم بر نرخ کلید \lr{QKD} و هم بر زمان اجرای مؤثر الگوریتم‌های \lr{PQC} تأثیر می‌گذارد و باید در مدل‌سازی سیستم لحاظ شود \{5\}.

%% ──────────────────────────────────────────────────────────────────────────
\section{مدل سیستم و مفروضات محیطی}
%% ──────────────────────────────────────────────────────────────────────────
\subsection{مشخصات ماهواره هدف و هندسه مداری}

سیستم مورد مطالعه در این پژوهش یک ماهواره کوچک در مدار پایین زمین \lr{(LEO)} در ارتفاع اسمی ۵۰۰ کیلومتر است. این ارتفاع بر اساس معیار کارایی اثبات‌شده‌ای انتخاب شده است: مطالعات کمی نشان می‌دهند که متوسط نرخ کلید مخفی قابل تولید با رابطه تقریبی \lr{$\langle R \rangle = O(h^{-1})$} با ارتفاع مدار رابطه معکوس دارد \{23\}. مأموریت‌های پیشرو در این حوزه، از جمله ماهواره \lr{Micius} چینی در ارتفاع ۵۰۰ کیلومتر \{6\} و سامانه مرجع مرکز فضایی آلمان \lr{(DLR)} در ارتفاع ۵۶۷ کیلومتر \{23\}، همگی در این محدوده ارتفاعی قرار دارند.

ماهواره در یک مدار دایره‌ای با زاویه میل مداری \lr{(Orbit Inclination)} ۹۰ درجه فرض می‌شود، که با مدار هم‌خورشیدی \lr{(Sun-Synchronous Orbit --- SSO)} انطباق دارد. سرعت مداری خطی ماهواره در ارتفاع ۵۰۰ کیلومتر از رابطه مکانیک مداری محاسبه می‌شود:

\begin{equation}
v = \sqrt{\frac{GM_E}{R_E + h}}
\label{eq:orbital_velocity}
\end{equation}

که در آن \lr{$GM_E$} ثابت گرانشی-جرمی زمین، \lr{$R_E$} شعاع میانگین زمین و \lr{$h$} ارتفاع مداری است. نتیجه عددی این محاسبه به سرعتی حدود ۷.۶ کیلومتر بر ثانیه منجر می‌شود.

پنجره دید \lr{(Pass Window)} یا پنجره ارتباطی، فاصله زمانی‌ای است که در آن ماهواره بالای افق دیدِ ایستگاه زمینی قرار دارد. با فرض حداقل زاویه فراز \lr{(Elevation Angle --- $\alpha_{min}$)} عملیاتی برابر ۲۰ درجه، طول مسیر لینک به صورت تابع زاویه فراز از رابطه زیر محاسبه می‌شود:

\begin{equation}
L(\alpha) = -R_E \cdot \sin(\alpha) + \sqrt{R_E^2 \cdot \sin^2(\alpha) + h^2 + 2hR_E}
\label{eq:link_length}
\end{equation}

این رابطه نشان می‌دهد که فاصله لینک در زاویه فراز ۲۰ درجه به حدود ۱۱۹۰ کیلومتر می‌رسد، در حالی که در زاویه سمت‌الرأس \lr{(Zenith)} مستقیماً برابر با ارتفاع مداری یعنی ۵۰۰ کیلومتر است. مدت زمان مفید هر گذر \lr{(Pass)} در مدار مورد نظر بین ۵ تا ۱۰ دقیقه است \{8\}؛ بر اساس داده‌های تجربی \lr{DLR}، یک گذر بهینه با حداکثر زاویه فراز ۹۰ درجه می‌تواند تا ۲۹۴ ثانیه طول بکشد \{23\}.

\subsection{مدل کانال اپتیکی فضای آزاد}

کانال ارتباطی مابین ماهواره \lr{LEO} و ایستگاه زمینی از نوع کانال اپتیکی فضای آزاد \lr{(Free-Space Optical Channel --- FSO)} است. کارایی انتقال کل کانال \lr{(Total Channel Efficiency --- $\eta_{total}$)} به صورت حاصل‌ضرب چند مؤلفه مستقل مدل می‌شود \{8\}:

\begin{equation}
\eta_{total}(\alpha) = \eta_{diff}(\alpha) \cdot \eta_{atm}(\alpha) \cdot \eta_{turb}(\alpha) \cdot \eta_{pat}(\alpha) \cdot \eta_{rx} \cdot \eta_{coupling} \cdot \eta_{pde}
\label{eq:total_eta}
\end{equation}

پیکربندی مرجع این پژوهش برای محاسبه بودجه لینک به‌صورت یکدست عبارت است از: قطر تلسکوپ فرستندهٔ روی ماهواره \lr{$D_t = 0.10$~m}، قطر تلسکوپ گیرندهٔ زمینی \lr{$D_r = 0.60$~m}، نسبت انسداد مرکزی تلسکوپ کاسگرین \lr{$\gamma = 0.30$}، طول موج \lr{$\lambda = 850$~nm}، و گذردهی جوی زنیت محافظه‌کارانهٔ \lr{$\eta_{zen} = 0.50$}.

\textbf{اتلاف انتشار \lr{(Diffraction Loss)}:} بزرگ‌ترین سهم در بودجه لینک، از اتساع \lr{(Divergence)} ذاتی پرتو لیزر ناشی می‌شود. برای یک پرتو گاوسی، نیم‌زاویه واگرایی استاندارد در میدان دور با شعاع کمر پرتو \lr{$w_0$} به‌صورت \lr{$\theta_{1/2} = \lambda / (\pi w_0)$} تعریف می‌شود. برای جلوگیری از پدیده برش پرتو \lr{(Truncation)} در لبه‌های آپرچر فرستنده، شعاع کمر پرتو معمولاً به‌صورت \lr{$w_0 = D_t/2.2$} انتخاب می‌شود. در فواصل دور فضایی، شعاع لکه پرتو \lr{($w_L$)} در گیرنده تقريباً برابر با \lr{$\theta_{1/2}\,L(\alpha)$} است. با در نظر گرفتن توزیع شدت گاوسی پرتو بر روی دیافراگم حلقوی تلسکوپ گیرنده کاسگرین، کارایی انتشار به‌طور دقیق از طریق انتگرال‌گیری شدت محاسبه شده و به شکل زیر مدل می‌شود \{8\}:

\begin{equation}
\eta_{diff}(\alpha) = \exp\left(-2\left(\frac{r_{in}}{w_L}\right)^2\right) - \exp\left(-2\left(\frac{a}{w_L}\right)^2\right)
\label{eq:diff_loss}
\end{equation}

در این رابطه، \lr{$a = D_r/2$} شعاع فیزیکی آپرچر \lr{(Aperture)} تلسکوپ گیرنده، \lr{$r_{in} = \gamma a$} شعاع انسداد مرکزی \lr{(Central Obstruction)} ناشی از آینه ثانویه، و \lr{$L(\alpha)$} فاصله لینک به عنوان تابع زاویه فراز است. با پیکربندی مرجع، نیم‌زاویه واگرایی به \lr{$\theta_{1/2} \approx 5.41~\mu\text{rad}$} (زاویه کامل \lr{$\approx 10.82~\mu\text{rad}$}) می‌رسد و اتلاف کل انتشار در گذر سمت‌الرأس (فاصله ۵۰۰ کیلومتر) به حدود ۱۹.۵ دسیبل می‌رسد که با مقادیر گزارش‌شده در آزمایش‌های \lr{Micius} هم‌خوان است \{6\}.

\textbf{اتلاف جوی \lr{(Atmospheric Attenuation)}:} گذردهی جوی \lr{(Atmospheric Transmissivity)} به صورت تابعی از زاویه فراز بر پایه مدل جرم هوا \lr{(Air Mass)} مدل می‌شود \{23\}:

\begin{equation}
\eta_{atm}(\alpha) = \eta_{zen}^{1/\sin(\alpha)}
\label{eq:atm_eta}
\end{equation}

که در آن \lr{$\eta_{zen}$} گذردهی جوی در راستای سمت‌الرأس است. برای طول موج ۸۵۰ نانومتر و با مقدار محافظه‌کارانهٔ \lr{$\eta_{zen} \approx 0.50$}، اتلاف سمت‌الرأسی حدود ۳ دسیبل است. اما در زاویه فراز ۲۰ درجه، مسیر طی‌شده در جو به ضریب \lr{$\csc(20^\circ) \approx 2.9$} برابر مسیر سمت‌الرأس می‌شود که این اتلاف جوی را به حدود ۸.۸ دسیبل افزایش می‌دهد. پیکربندی لینک دانلود \lr{(Downlink)} در مقایسه با لینک آپلود \lr{(Uplink)} از این جهت مزیت دارد که پرتو ابتدا در خلأ منتشر می‌شود و تنها در ۱۰ تا ۲۰ کیلومتر آخر با جو برهم‌کنش می‌کند \{4\}.

\textbf{تلاطم جوی \lr{(Atmospheric Turbulence)}:} شدت پدیده‌های تلاطم \lr{(Scintillation, Beam Wander)} با پارامتر ساختار ضریب شکست \lr{$C_n^2$} توصیف می‌شود. برای لینک دانلود \lr{LEO} که پرتو تنها در قسمت انتهایی مسیر خود با جو برهم‌کنش می‌کند، اثرات تلاطم در سمت‌الرأس نسبتاً ملایم است و کارایی تلاطم جوی در حدود \lr{$\eta_{turb} \approx 0.63$} (اتلافی معادل ۲ دسیبل) مدل می‌شود؛ اما با کاهش زاویه فراز افزایش می‌یابد و در زاویه ۲۰ درجه به حدود ۳.۶ دسیبل می‌رسد \{8\}.

بودجه اتلاف کل کانال با جمع‌بندی تمامی مؤلفه‌های فوق به دست می‌آید. مطابق با مدل اصلاح‌شده و داده‌های تجربی مأموریت‌های \lr{LEO}، اتلاف کل کانال در یک گذر بهینه با زاویه فراز ۹۰ درجه به حدود ۳۱ دسیبل می‌رسد، در حالی که در گذرهای با زاویه فراز ۲۰ درجه این عدد به حدود ۴۷ دسیبل افزایش می‌یابد \{6, 8\}.
\subsection{محدودیت‌های \lr{SWaP} و ظرفیت محاسباتی}

مفهوم \lr{SWaP} (اندازه، وزن، و توان --- \lr{Size, Weight, and Power}) به مجموعه‌ای از محدودیت‌های فیزیکی و توانی اشاره دارد که طراحی هر سیستم الکترونیکی مستقر در فضا را به‌شدت محدود می‌کند. این محدودیت‌ها در ماهواره‌های کوچک از نوع \lr{CubeSat} یا میکروماهواره --- که سکوی هدف مدل‌سازی این پژوهش هستند --- با شدت بیشتری نسبت به ماهواره‌های بزرگ اعمال می‌شوند \{5\}.

در حوزه توان الکتریکی، پنل‌های خورشیدی یک \lr{CubeSat} از نوع \lr{6U} معمولاً توانی بین ۱۰ تا ۳۰ وات تولید می‌کنند. این توان باید میان تمامی زیرسیستم‌های ماهواره --- شامل رایانه مرکزی \lr{(On-Board Computer --- OBC)}، سیستم ارتباطات، سیستم کنترل وضعیت \lr{(ACS)}، محموله علمی و زیرسیستم دما --- تقسیم شود. در حوزه ظرفیت پردازشی، \lr{OBC} ماهواره‌های کوچک معمولاً از پردازنده‌های \lr{ARM Cortex-M} یا معماری \lr{RISC-V} با فرکانس کاری کمتر از ۲۰۰ مگاهرتز استفاده می‌کنند. پیاده‌سازی‌های بهینه‌شده با بسط مجموعه دستورالعمل‌ها \lr{(ISA Extension)} برای الگوریتم ضرب ماژولار با تبدیل نظریه اعداد \lr{(Number Theoretic Transform --- NTT)} می‌توانند زمان اجرای رمزنگاری را به مرتبه هزاران چرخه کاهش دهند \{16\}.

علاوه بر محدودیت‌های سه‌گانه \lr{SWaP}، یک چالش منحصربه‌فرد محیط فضایی نیز باید لحاظ شود: تشعشعات یونیزان کیهانی \lr{(Cosmic Ionizing Radiation)} می‌توانند رویدادهای تک‌رویداد اختلال \lr{(Single Event Upset --- SEU)} در مدارات دیجیتال ایجاد کنند \{5\}. در معماری پیشنهادی این پژوهش، راه‌حل افزونگی زمانی \lr{(Temporal Redundancy)} برای مؤلفه‌های رمزنگاری حساس --- که در آن هر عملیات دو بار با ترتیب رجیسترهای معکوس اجرا شده و نتیجه \lr{XOR} می‌شود تا خطا کشف گردد --- به‌عنوان یک فرض طراحی در نظر گرفته می‌شود \{16\}.

%% ─── جدول ۳-۱ ──────────────────────────────────────────────────────────────
\begin{table}[H]
\caption{پارامترهای مدل سیستم مرجع \lr{LEO} برای شبیه‌سازی}
\label{tab:3-1}
\centering
\small
\begin{tabular}{|r|c|c|}
\hline
\bfseries پارامتر & \bfseries مقدار & \bfseries مرجع \\
\hline
ارتفاع مداری ($h$) & ۵۰۰ کیلومتر & \{6, 8\} \\
\hline
زاویه میل مداری & ۹۰ درجه (\lr{SSO}) & \{23\} \\
\hline
سرعت مداری خطی & \lr{$\sim$۷.۶ km/s} & محاسبه‌شده \\
\hline
حداقل زاویه فراز عملیاتی & ۲۰ درجه & \{23\} \\
\hline
مدت گذر بهینه & \lr{$\sim$۲۹۴ s} & \{8, 23\} \\
\hline
طول موج کوانتومی & ۸۵۰ نانومتر & \{6, 8, 23\} \\
\hline
اتلاف کانال در گذر بهینه & \lr{$\sim$۳۴ dB} & \{8\} \\
\hline
اتلاف تلاطم جوی ($\eta_{turb}$) & \lr{2 dB (0.631)} & \{8\} \\
\hline
اتلاف سیستم \lr{PAT} ($\eta_{pat}$) & \lr{0.4 dB (0.912)} & \{8\} \\
\hline
بودجه توان محموله رمزنگاری & $\le$ ۱۰ وات & \{5, 24\} \\
\hline
پردازنده هدف & \lr{RISC-V 32-bit / ARM Cortex-M} & \{16, 24\} \\
\hline
\end{tabular}
\end{table}

\begin{figure}[H]
\centering
\fbox{\parbox{0.78\textwidth}{\centering\vspace{1.5cm}
[هندسه مداری ماهواره \lr{LEO}، پارامترهای کانال اپتیکی فضای آزاد و عوامل اتلاف مؤثر بر بودجه لینک]
\vspace{1.5cm}}}
\caption{هندسه مداری ماهواره \lr{LEO}، پارامترهای کانال اپتیکی فضای آزاد و عوامل اتلاف مؤثر بر بودجه لینک}
\label{fig:3-1}
\end{figure}

%% ──────────────────────────────────────────────────────────────────────────
\section{معماری ترکیبی پیشنهادی}
%% ──────────────────────────────────────────────────────────────────────────

\subsection{منطق بنیادین ترکیب و دلیل انتخاب رویکرد سه‌لایه}

طراحی هر معماری امنیتی با یک پرسش بنیادین آغاز می‌شود: در برابر چه تهدیدی، با چه مدل اعتماد \lr{(Trust Model)} و برای چه مدت زمانی باید از داده‌ها محافظت شود؟ در بستر ارتباطات ماهواره‌ای \lr{LEO}، پاسخ به این پرسش چندوجهی است و دقیقاً همین چندوجهی بودن است که رویکرد تک‌لایه را --- چه \lr{QKD} به‌تنهایی و چه \lr{PQC} به‌تنهایی --- ناکافی می‌کند.

\lr{QKD} با تکیه بر اصول بنیادین مکانیک کوانتوم، امنیت اطلاعاتی-تئوریک را تضمین می‌کند \{4\}. اما \lr{QKD} دارای یک آسیب‌پذیری ساختاری است: کانال کلاسیکی که برای پس‌پردازش استفاده می‌شود، خود باید به‌طور مستقلی احراز هویت شده باشد. بدون این احراز هویت، حمله مرد میانی \lr{(Man-in-the-Middle --- MITM)} می‌تواند کل پروتکل را به خطر بیندازد \{10\}. از سوی دیگر، ظرفیت بنیادین کانال فضایی، نرخ کلید \lr{QKD} در یک گذر ۵ تا ۱۰ دقیقه‌ای را به چند مگابیت محدود می‌کند \{8\}، که برای تأمین نیاز ارتباط پیوسته کافی نیست.

\lr{PQC} از سوی دیگر مستقل از محدودیت‌های فیزیکی کانال است؛ یک \lr{KEM} مبتنی بر شبکه مشبک مانند \lr{ML-KEM} می‌تواند در هر زمان و با هر بودجه لینکی یک کلید جلسه ایمن برقرار کند. اما امنیت \lr{PQC} بر سختی محاسباتی مسائل ریاضی استوار است که اثبات امنیتی آن در مدل \lr{QROM} هنوز کاملاً محکم نیست \{21\}؛ و تجربه شکسته شدن \lr{SIKE} در ۲۰۲۲ و \lr{Rainbow} در همان سال \{2\} نشان داد که تکیه صرف بر یک الگوریتم \lr{PQC} ریسک بلندمدت ایجاد می‌کند.

نتیجه‌گیری منطقی از این تحلیل، رویکرد ادغام ترکیبی \lr{(Hybrid Integration)} است که محور طراحی معماری پیشنهادی این پژوهش --- و مبتنی بر پیاده‌سازی تجربی \lr{Rani} و همکاران \{13\} و چارچوب نظری \lr{Zeng} و همکاران \{19\} --- است.

\begin{figure}[H]
\centering
\fbox{\parbox{0.78\textwidth}{\centering\vspace{1.5cm}
[نمای لایه‌بندی‌شده معماری ترکیبی پیشنهادی؛ یکپارچه‌سازی لایه‌های \lr{QKD}، \lr{ML-KEM} و \lr{ML-DSA} در یک سیستم امنیتی سه‌لایه]
\vspace{1.5cm}}}
\caption{نمای لایه‌بندی‌شده معماری ترکیبی پیشنهادی؛ یکپارچه‌سازی لایه‌های \lr{QKD}، \lr{ML-KEM} و \lr{ML-DSA}}
\label{fig:3-2}
\end{figure}

\subsection{لایه اول --- توزیع کلید کوانتومی با پروتکل \lr{Decoy-State BB84}}

لایه اول معماری پیشنهادی، لایه تولید کلید اصلی \lr{(Master Key Layer)} است که بر مبنای پروتکل \lr{BB84} با حالت‌های فریب \lr{(Decoy-State Protocol)} عمل می‌کند. این پروتکل به دو دلیل انتخاب شده است: اول، بالاترین سطح بلوغ در اثبات‌های امنیتی رسمی را در میان تمام پروتکل‌های \lr{QKD} ماهواره‌ای دارد \{23\}؛ دوم، آزمایش‌های عملی روی ماهواره \lr{Micius} اعداد کمی واقع‌بینانه‌ای ارائه داده‌اند \{6, 23\}.

در این پروتکل، ماهواره در نقش \lr{Alice} (فرستنده) عمل می‌کند و پالس‌های لیزری تضعیف‌شده را در طول موج ۸۵۰ نانومتر با نرخ تکرار ۱ گیگاهرتز --- هم‌تراز با سامانه‌های گیگاهرتزی نسل جدید \{8\} --- به سمت ایستگاه زمینی ارسال می‌کند. مکانیزم دو-حالت فریب \lr{(Two-Decoy Protocol)} اجازه می‌دهد تا از حمله جداسازی تعداد فوتون \lr{(Photon Number Splitting --- PNS)} جلوگیری شود. نرخ کلید مخفی \lr{(Secret Key Rate --- SKR)} در حالت مجانبی با رابطه زیر تعریف می‌شود \{23\}:

\begin{equation}
R = p_\mu \cdot p_Z \cdot \left[Q_1^Z \cdot \left(1 - H(e_1^X)\right) - f_{EC} \cdot Q_\mu^Z \cdot H(E_\mu^Z)\right]
\label{eq:skr}
\end{equation}

که در آن \lr{$p_\mu$} احتمال ارسال حالت سیگنال، \lr{$p_Z$} احتمال انتخاب پایه \lr{Z}، \lr{$Q_1^Z$} نرخ تک-فوتونی در پایه \lr{Z}، \lr{$e_1^X$} نرخ خطای فاز، \lr{$H$} تابع آنتروپی دوتایی \lr{Shannon}، \lr{$f_{EC}$} ضریب ناکارایی تصحیح خطا، \lr{$Q_\mu^Z$} نرخ کلیک در پایه \lr{Z}، و \lr{$E_\mu^Z$} نرخ خطای بیت کوانتومی \lr{(QBER)} اندازه‌گیری‌شده است.

پارامترهای بهینه‌سازی‌شده برای ارتفاع ۵۰۰ کیلومتر مطابق داده‌های مرجع \lr{DLR} عبارتند از: \lr{$\mu \approx 0.69$}، \lr{$\nu \approx 0.14$}، \lr{$p_Z \approx 0.93$}. با این پارامترها و نرخ تکرار ۱ گیگاهرتز، طول کلید مخفی قابل دستیابی در یک گذر بهینه در حدود چند مگابیت است که با مقدار ۴.۵۸ مگابیت گزارش‌شده در آزمایش گیگاهرتزی واقعی هم‌مرتبه است \{8\}. این حجم از کلید به‌عنوان کلید اصلی \lr{(Master Key)} در مخزن کلید \lr{(Key Pool)} ذخیره می‌شود و در فاصله میان گذرهای مداری، به تدریج توسط لایه‌های بالاتر معماری مصرف می‌شود. همین محدودیتِ ذاتیِ حجم کلید در هر گذر --- در مقایسه با نیاز پیوستهٔ ارتباطات پرحجم --- دقیقاً دلیل مهندسیِ ضرورت حضور لایه‌های \lr{PQC} برای تأمین تداوم نشست است؛ مصالحه‌ای که محور معماری ترکیبی این پژوهش است.
\subsection{لایه دوم --- برپایی کلید جلسه با \lr{ML-KEM}}

لایه دوم معماری، لایه برپایی کلید جلسه \lr{(Session Key Establishment Layer)} است که بر اساس استاندارد \lr{FIPS 203 (ML-KEM)} عمل می‌کند. این لایه به‌طور مداوم فعال است و وظیفه اصلی آن تأمین یک کلید جلسه ۲۵۶ بیتی برای رمزنگاری همه داده‌های ارتباطی در فاصله میان گذرهای \lr{QKD} است.

\lr{ML-KEM} بر مسئله یادگیری با خطای ماژولار \lr{(MLWE)} استوار است که سختی محاسباتی آن هم در برابر الگوریتم‌های کلاسیک و هم در برابر الگوریتم‌های کوانتومی اثبات شده \{3\}. ساختار دو-مرحله‌ای آن --- ابتدا \lr{K-PKE} با امنیت \lr{IND-CPA}، سپس ارتقاء به \lr{IND-CCA2} از طریق تبدیل \lr{Fujisaki-Okamoto} --- تضمین می‌کند که حتی یک مهاجم فعال که توانایی ارسال متون رمزشده دلخواه دارد، قادر به استخراج کلید مشترک نخواهد بود \{3\}.

در معماری پیشنهادی، سه مجموعه پارامتری \lr{ML-KEM} در سه سطح امنیتی \lr{NIST} به کار گرفته می‌شوند: \lr{ML-KEM-512} (سطح ۱، معادل \lr{AES-128}) با کلید عمومی ۸۰۰ بایت و متن رمزشده ۷۶۸ بایتی، \lr{ML-KEM-768} (سطح ۳، معادل \lr{AES-192}، انتخاب پیش‌فرض توصیه‌شده \lr{NIST}) با کلید عمومی ۱۱۸۴ بایت و متن رمزشده ۱۰۸۸ بایتی، و \lr{ML-KEM-1024} (سطح ۵، معادل \lr{AES-256}) با کلید عمومی ۱۵۶۸ بایت و متن رمزشده ۱۵۶۸ بایتی \{3\}.

جریان عملیاتی \lr{ML-KEM} در این معماری شامل چهار گام است: (۱) ایستگاه زمینی با اجرای \lr{ML-KEM.KeyGen} یک جفت کلید تولید کرده و کلید عمومی را از طریق کانال کلاسیک احراز هویت‌شده ارسال می‌کند؛ (۲) ماهواره پس از اعتبارسنجی کلید دریافتی، \lr{ML-KEM.Encaps} را اجرا کرده و جفت \lr{$(K_{PQC}, C)$} را تولید می‌کند؛ (۳) متن رمزشده \lr{$C$} برای ایستگاه زمینی ارسال می‌شود؛ (۴) ایستگاه زمینی با \lr{ML-KEM.Decaps} مقدار \lr{$K_{PQC}$} را استخراج می‌کند. از این پس، هر دو طرف یک کلید مشترک ۲۵۶ بیتی دارند که برای رمزنگاری ارتباطات با \lr{AES-256-GCM} به کار می‌رود.

\subsection{لایه سوم --- احراز هویت پسا-کوانتومی با \lr{ML-DSA}}

لایه سوم، لایه احراز هویت \lr{(Authentication Layer)} است که بر اساس استاندارد \lr{FIPS 204 (ML-DSA)} عمل می‌کند. این لایه وظیفه حفاظت از یکپارچگی و اصالت کانال کلاسیک \lr{QKD} را بر عهده دارد.

استدلال امنیتی ترکیبی در این‌جا ظریف اما قدرتمند است: \lr{ML-DSA} برای احراز هویت \lr{QKD} تنها نیاز به امنیت کوتاه‌مدت \lr{(Short-Term Security)} دارد --- اگر امضای دیجیتال در لحظه ارسال شکسته نشود، کلیدهای \lr{QKD} توزیع‌شده برای همیشه از طریق اصول مکانیک کوانتوم محافظت می‌شوند، حتی اگر \lr{ML-DSA} در آینده‌ای دور به خطر بیفتد \{10\}.

در معماری پیشنهادی، \lr{ML-DSA-65} (سطح ۳، معادل \lr{AES-192}) به‌عنوان الگوریتم امضا انتخاب شده است. دلایل این انتخاب در برابر رقیب اصلی \lr{FALCON-512} سه‌گانه است: اول، \lr{ML-DSA} صراحتاً استفاده از حساب ممیز شناور \lr{(Floating Point)} را ممنوع می‌کند \{9\} --- بر خلاف \lr{FALCON} که به این نوع حساب نیاز دارد و این نیاز آن را برای پردازنده‌های مقاوم در برابر تشعشع با واحد \lr{FPU} محدود، ناسازگار می‌کند؛ دوم، حالت امضای محافظت‌شده \lr{(Hedged Signing)} در \lr{ML-DSA} در برابر حملات تزریق خطا \lr{(Fault Attacks)} ناشی از تابش کیهانی مقاوم‌تر است \{9\}؛ سوم، \lr{ML-DSA} دارای ویژگی \lr{BUFF (Beyond UnForgeability)} است که عدم انکار \lr{(Non-Repudiation)} را در سطح بالاتری تضمین می‌کند --- ویژگی حیاتی برای تله‌متری و دستورات کنترلی ماهواره \{9\}.

\lr{ML-DSA} در این معماری دو نقش مجزا ایفا می‌کند: احراز هویت هر پیام پس‌پردازش کانال کلاسیک \lr{QKD} (از پیام‌های پالایش پایه گرفته تا سندرم‌های تصحیح خطا)، و احراز هویت متقابل \lr{(Mutual Authentication)} قبل از شروع جلسه \lr{KEM} با استفاده از عدد یکبارمصرف \lr{(Nonce)} تصادفی برای جلوگیری از حمله بازپخش \lr{(Replay Attack)}.

\subsection{روند توالی یکپارچه و جریان کنترل نشست}

جریان توالی یک نشست ارتباطی کامل شامل چهار مرحله اصلی است. در مرحله اول (شناسایی و احراز هویت متقابل)، هنگامی که ماهواره در آستانه ورود به محدوده دید ایستگاه زمینی قرار می‌گیرد، هر دو طرف یک عدد یکبارمصرف تصادفی ۳۲ بایتی تولید کرده و امضاهای متقابل \lr{ML-DSA} مبادله می‌کنند:
\begin{equation}
\sigma_{Sat} = ML\text{-}DSA.Sign(sk_{Sat},\ N_{GS} \| H(pk_{Sat}))
\label{eq:mldsa_sign}
\end{equation}

در مرحله دوم (برپایی کلید جلسه \lr{PQC})، ایستگاه زمینی یک جفت کلید \lr{ML-KEM} تولید کرده و کلید عمومی را ارسال می‌کند. ماهواره عملیات کپسوله‌سازی را انجام داده و \lr{$K_{PQC}$} استخراج می‌شود.

در مرحله سوم (تلاش برای برقراری لایه \lr{QKD})، هم‌زمان با انجام مراحل قبل، سیستم \lr{QKD} فرستنده روی ماهواره فعال می‌شود و ارسال پالس‌های کوانتومی را آغاز می‌کند. هر پیام پالایش پایه، تصحیح خطا و تقویت حریم خصوصی با \lr{ML-DSA} امضا و تأیید می‌شود \{10\}.

در مرحله چهارم (ترکیب و اولویت‌بندی کلید)، اگر هر دو لایه \lr{QKD} و \lr{PQC} با موفقیت کلید تولید کرده باشند، معماری پیشنهادی از رویکرد ترکیب موازی \lr{(Parallel XOR Combination)} استفاده می‌کند \{19\}:
\begin{equation}
K_{final} = K_{QKD} \oplus KDF(K_{PQC})
\label{eq:key_combination}
\end{equation}

که در آن \lr{KDF} یک تابع مشتق کلید \lr{(Key Derivation Function)} استاندارد مانند \lr{HKDF-SHA256} است. این ترکیب اطمینان می‌دهد که حتی اگر یکی از دو لایه کاملاً به خطر بیفتد، کلید نهایی همچنان امن خواهد بود \{13, 19\}.

\begin{figure}[H]
\centering
\fbox{\parbox{0.78\textwidth}{\centering\vspace{1.5cm}
[نمودار توالی یکپارچه نشست ارتباطی امن؛ از احراز هویت متقابل \lr{ML-DSA} تا برپایی کلید \lr{ML-KEM} و فعال‌سازی کانال \lr{QKD}]
\vspace{1.5cm}}}
\caption{نمودار توالی یکپارچه نشست ارتباطی امن}
\label{fig:3-3}
\end{figure}

%% ─── جدول ۳-۲ ──────────────────────────────────────────────────────────────
\begin{table}[H]
\caption{مشخصات الگوریتم‌های منتخب در معماری پیشنهادی}
\label{tab:3-2}
\centering
\scriptsize
\begin{tabular}{|r|c|p{2.8cm}|c|c|c|c|}
\hline
\bfseries الگوریتم & \bfseries استاندارد & \bfseries نقش در معماری & \bfseries سطح \lr{NIST} & \bfseries کلید عمومی & \bfseries خروجی & \bfseries مبنا \\
\hline
\lr{ML-KEM-512} & \lr{FIPS 203} & برپایی کلید جلسه (سطح ۱) & \lr{Cat. 1} & ۸۰۰ B & \lr{CT: 768 B} & \lr{MLWE} \\
\hline
\lr{ML-KEM-768} & \lr{FIPS 203} & برپایی کلید جلسه (پیش‌فرض) & \lr{Cat. 3} & ۱۱۸۴ B & \lr{CT: 1088 B} & \lr{MLWE} \\
\hline
\lr{ML-KEM-1024} & \lr{FIPS 203} & برپایی کلید جلسه (سطح ۵) & \lr{Cat. 5} & ۱۵۶۸ B & \lr{CT: 1568 B} & \lr{MLWE} \\
\hline
\lr{ML-DSA-65} & \lr{FIPS 204} & احراز هویت کانال کلاسیک & \lr{Cat. 3} & ۱۹۵۲ B & \lr{Sig: 3309 B} & \lr{MLWE+MSIS} \\
\hline
\lr{BB84 Decoy} & --- & تولید کلید اصلی (\lr{QKD}) & اطلاعاتی-تئوریک & --- & چند \lr{Mbit/pass} & مکانیک کوانتوم \\
\hline
\end{tabular}
\end{table}

%% ──────────────────────────────────────────────────────────────────────────
\section{مکانیزم پشتیبان و طراحی ایمن در برابر خرابی}
%% ──────────────────────────────────────────────────────────────────────────

\subsection{فلسفه بنیادین \lr{Fail-Safe} در سیستم‌های فضایی}

در مهندسی قابلیت اطمینان \lr{(Reliability Engineering)}، اصل طراحی ایمن در برابر خرابی \lr{(Fail-Safe Design)} بر این بنیاد استوار است که خرابی یک مؤلفه نباید به خرابی کل سیستم --- و به‌ویژه به نقض ضمانت‌های امنیتی --- منجر شود. در سیستم‌های ارتباطات فضایی، این اصل به سطح کیفاً متفاوتی ارتقاء می‌یابد: برخلاف یک سرور زمینی که می‌توان آن را در عرض چند دقیقه بازنشانی یا جایگزین کرد، یک ماهواره در مدار \lr{LEO} پس از پرتاب قابل تعمیر نیست، در هر گذر تنها ۵ تا ۱۰ دقیقه پنجره ارتباطی دارد \{8\}، و تشعشعات کیهانی به صورت پیوسته بر قطعات الکترونیکی آن اثر می‌گذارند \{5\}.

معماری پیشنهادی این پژوهش بر اساس یک اصل مهندسی محوری شکل گرفته است: هر لایه باید ضامن پشتیبان \lr{(Backup Guarantor)} لایه بالاتر باشد، به‌گونه‌ای که از کار افتادن هر تعداد لایه --- به جز آخرین لایه --- نه موجب قطع ارتباط شود و نه موجب کاهش کیفی ضمانت‌های امنیتی که به کاربر اعلام نشده است. این رویکرد را در نظریه سیستم‌های قابل اعتماد، تنزل رتبه بدون فاجعه \lr{(Graceful Degradation without Catastrophic Failure)} می‌نامند.

\subsection{سناریوی اول --- قطع شدن لایه \lr{QKD}}

متداول‌ترین سناریوی خرابی، از کار افتادن لایه توزیع کلید کوانتومی است. دلایل آن می‌توانند شامل شرایط جوی نامساعد، تلاطم جوی شدید، زاویه فراز پایین ماهواره، یا خرابی سخت‌افزاری در منبع لیزر یا آشکارساز باشند.

معیار تشخیص خرابی در لایه \lr{QKD} به صورت یک پایش پیوسته نرخ خطای بیت کوانتومی \lr{(QBER)} طراحی شده است. آستانه بنیادین امنیتی پروتکل \lr{BB84} برابر ۱۱٪ است که بالاتر از آن هیچ کلید امنی قابل استخراج نیست؛ برای حفظ یک حاشیه ایمنی، ماشین حالت امنیتی یک آستانه هشدار \lr{(Warning Threshold)} برابر ۸٪ و یک آستانه سقوط پیشگیرانه \lr{(Fallback Threshold)} برابر ۱۰٪ تعریف می‌کند \{5\}. اگر \lr{QBER} در سه اندازه‌گیری متوالی از آستانه هشدار ۸٪ فراتر رود، شمارنده پایداری \lr{(Persistence Counter)} فعال می‌شود تا از نوسان ناشی از نویز آنی جلوگیری گردد؛ و چنانچه \lr{QBER} از آستانه سقوط ۱۰٪ عبور کند یا در طول گذر هیچ بلوک کلید پالایش‌شده‌ای در ظرف ۶۰ ثانیه ساخته نشود، سیستم حالت \lr{QKD-Unavailable} را اعلام می‌کند. در این لحظه، ماشین حالت امنیتی \lr{(Security State Machine)} به صورت خودکار به حالت \lr{PQC-Only} سوییچ می‌کند.

در حالت \lr{PQC-Only}، تمام ارتباطات داده‌ای و کنترلی صرفاً با کلید جلسه‌ای که از \lr{ML-KEM} به دست آمده است با الگوریتم \lr{AES-256-GCM} رمزنگاری می‌شوند. فاصله تجدید کلید در این حالت اضطراری از ۳۰ دقیقه به ۱۵ دقیقه کاهش می‌یابد تا دوره افشای بالقوه محدود شود. همچنین یک بافر پیوستگی \lr{(Continuity Buffer)} به اندازه حداقل ۱۰۰ میلی‌ثانیه تضمین می‌کند که هیچ فریم داده‌ای در طول انتقال حالت از دست نرود یا بدون رمزنگاری ارسال نشود.

\subsection{سناریوی دوم --- تهی شدن مخزن کلید}

مدیریت مخزن کلید \lr{(Key Pool Management --- KPM)} در این معماری بر اساس یک سیاست آستانه‌محور \lr{(Threshold-Based Policy)} با سه سطح هشدار طراحی شده است. در سطح اول (مخزن زیر ۵۰\%)، سیستم وارد حالت «حفاظت از مخزن» می‌شود و کلیدهای \lr{QKD} صرفاً برای تضمین پنهان‌سازی پیش‌رو \lr{(Forward Secrecy)} به کار می‌روند. در سطح دوم (مخزن زیر ۲۰\%)، سیستم وارد حالت «ذخیره اضطراری» شده و استفاده از کلیدهای \lr{QKD} به مهم‌ترین پیام‌های کنترلی محدود می‌شود. در سطح سوم (مخزن زیر ۵\%)، سیستم به صورت کامل به حالت \lr{PQC-Exclusive} سوییچ می‌کند.

این رویکرد از درسی مستقیم برگرفته شده که آزمایش بین‌قاره‌ای ماهواره \lr{Micius} آن را نشان داد: مصرف ۷۰ کیلوبایت کلید \lr{QKD} برای یک کنفرانس ویدیویی ۷۵ دقیقه‌ای نشان می‌دهد که بدون یک سیاست مدیریت مخزن هوشمند، خروجی هر گذر ماهواره‌ای ممکن است در عرض چند ساعت مصرف شود \{11\}. مکانیزم پسماند \lr{(Hysteresis)} تعریف‌شده --- که سیستم تنها پس از رسیدن مخزن به بیش از ۳۰\% به حالت عادی بازمی‌گردد --- از نوسان سریع بین حالت‌ها جلوگیری می‌کند.

\subsection{سناریوی سوم --- خطاهای نرم‌افزاری و تشعشعی}

رویدادهای \lr{SEU} در ارتفاع ۵۰۰ کیلومتر از منظر رمزنگاری خطرناک هستند، چرا که می‌توانند منجر به قابل پیش‌بینی شدن خروجی \lr{RNG} یا استفاده مجدد ناخواسته از \lr{Nonce} شوند. در معماری پیشنهادی، سه سطح دفاعی در نظر گرفته شده است.

سطح اول، تشخیص خطا از طریق افزونگی زمانی \lr{(Temporal Redundancy)} است: هر عملیات رمزنگاری حیاتی دو بار اجرا می‌شود، بار دوم با ترتیب معکوس رجیسترها، و نتیجه \lr{XOR} دو خروجی بررسی می‌شود. این تکنیک با سربار تنها ۱۰۰\% در چرخه‌های \lr{CPU} و ۱\% در مساحت مدار، تمامی خطاهای گذرای تک‌بیتی را تشخیص می‌دهد \{16\}.

سطح دوم، حفاظت از حافظه‌های حساس از طریق کدهای تصحیح خطای \lr{(Error Correcting Code --- ECC)} سخت‌افزاری است که خطاهای تک‌بیتی را تصحیح و خطاهای دو-بیتی را تشخیص می‌دهد.

سطح سوم، یک ماژول ناظر رمزنگاری \lr{(Cryptographic Watchdog)} مستقل است که به صورت پیوسته یکپارچگی کد اجرایی ماژول رمزنگاری را از طریق اعتبارسنجی درهم‌ساز \lr{(Hash Verification)} بررسی می‌کند. در صورت مغایرت، یک فرآیند بازنشانی امن \lr{(Secure Reset)} آغاز می‌شود که ابتدا تمام کلیدهای \lr{SRAM} را صفرسازی کرده و سپس سیستم را از \lr{ROM} مجدداً بارگذاری می‌کند.

\subsection{ماشین حالت امنیتی و الگوریتم تصمیم‌گیری}

قلب تپنده این معماری \lr{Fail-Safe}، یک ماشین حالت محدود \lr{(Finite State Machine --- FSM)} است که به صورت مداوم بر وضعیت لایه‌های مختلف نظارت می‌کند. این \lr{FSM} پنج حالت تعریف‌شده دارد:

\textbf{حالت \lr{H1} (امنیت کامل):} هر دو لایه \lr{QKD} و \lr{ML-KEM} فعال هستند، مخزن کلید \lr{QKD} بالای ۳۰\% است، نرخ خطای بیت کوانتومی \lr{(QBER)} زیر آستانه هشدار ۸\% است، و کلید نهایی از ترکیب \lr{$K_{QKD} \oplus KDF(K_{PQC})$} به دست می‌آید \{13, 19\}.

\textbf{حالت \lr{H2} (حفاظت از مخزن):} هنوز هر دو لایه فعال هستند اما سطح مخزن \lr{QKD} در محدوده ۵\% تا ۳۰\% قرار دارد. کلیدهای \lr{QKD} تنها برای تجدید دوره‌ای کلید جلسه به کار می‌روند.

\textbf{حالت \lr{H3} (اضطراری):} لایه \lr{QKD} غیرفعال است یا مخزن کلید تهی شده است. تمام ارتباطات بر پایه \lr{ML-KEM-768} هستند. فاصله تجدید کلید به ۱۵ دقیقه کاهش یافته است.

\textbf{حالت \lr{H4} (بازیابی):} سیستم \lr{QKD} در حال تلاش مجدد برای برقراری ارتباط است. \lr{QBER} به‌طور پایدار زیر آستانه بازیابی (۶.۸\%) بوده است و مخزن کلید در حال پر شدن است.

\textbf{حالت \lr{H5} (خاموشی ایمن):} یک خرابی رمزنگاری بحرانی شناسایی شده و سیستم تمام ارتباطات رمزنگاری را به طور کامل مسدود کرده است. تنها با یک دستور احراز هویت‌شده با \lr{ML-DSA} از ایستگاه کنترل زمینی قابل خروج است.

شرایط انتقال میان این حالت‌ها بر دو محور مستقل تعریف می‌شوند: سطح مخزن کلید و مقدار لحظه‌ای \lr{QBER}. گذار \lr{$H1 \rightarrow H2$} هنگامی رخ می‌دهد که سطح مخزن به باند ۵\% تا ۳۰\% افت کند، و گذار معکوس \lr{$H2 \rightarrow H1$} تنها پس از پر شدن مخزن به بالای ۳۰\% انجام می‌شود؛ این اختلاف آستانهٔ پر و خالی، یک پسماند \lr{(Hysteresis)} روی محور مخزن ایجاد می‌کند که از نوسان سریع جلوگیری می‌نماید. 

روی محور \lr{QBER}، برای جلوگیری از باگ چترینگ \lr{(Chattering)} ناشی از نویز آنی جوّ و تلاطمات لحظه‌ای \lr{(Scintillation)}، یک شمارنده پایداری \lr{(Persistence Counter)} با منطق «شمارنده افزایشی/کاهشی» \lr{(Up/Down Counter)} با عمق سه نمونه به‌صورت متقارن اعمال می‌شود. دلیل استفاده از این منطق (به جای شمارنده صفر-شونده) این است که یک افت لحظه‌ای کیفیت کانال یا یک نمونه نویز تصادفی، باعث دور زدن مکانیزم ایمنی و صفر شدن ناگهانی شمارنده هشدار نشود. بر این اساس، عبور پایدار از آستانه سقوط ۱۰\% (هنگامی که شمارنده پس از ثبت نمونه‌های بالاتر از آستانه هشدار ۸\% به حد نصاب ۳ برسد) سیستم را به حالت \lr{H3} می‌برد. بازگشت نیز تنها زمانی مجاز است که \lr{QBER} با ثبت در شمارنده کاهشی، به‌طور پایدار زیر آستانه بازیابی ۶.۸\% (معادل پسماند ۱۵\% نسبت به آستانه هشدار) قرار گیرد تا سیستم نخست وارد حالت بازیابی \lr{H4} و سپس \lr{H1} شود.

حالت \lr{H5} یک حالت قفل‌شونده \lr{(Latching)} است: ورود به آن با تشخیص خرابی بحرانی رمزنگاری (مغایرت درهم‌ساز ناظر یا خطای \lr{RNG} ناشی از \lr{SEU}) از هر حالتی ممکن است، اما خروج از آن تنها با دستور احراز هویت‌شدهٔ \lr{ML-DSA} از ایستگاه زمینی میسر است.

یک ویژگی حیاتی این \lr{FSM} آن است که هیچ انتقالی بدون ثبت در گزارش رویدادهای امنیتی \lr{(Security Event Log)} انجام نمی‌شود. هر انتقال حالت باعث تولید یک رکورد حسابرسی \lr{(Audit Record)} با زمان‌بندی دقیق می‌شود که با \lr{ML-DSA} امضا شده و برای ارسال به ایستگاه کنترل زمینی صف‌بندی می‌شود.

\begin{figure}[H]
\centering
\fbox{\parbox{0.78\textwidth}{\centering\vspace{1.5cm}
[ماشین حالت محدود امنیتی (\lr{Security FSM}) معماری \lr{Fail-Safe}؛ پنج حالت عملکردی و شرایط انتقال میان آن‌ها]
\vspace{1.5cm}}}
\caption{ماشین حالت محدود امنیتی \lr{(Security FSM)} معماری \lr{Fail-Safe}؛ پنج حالت عملکردی و شرایط انتقال}
\label{fig:3-4}
\end{figure}
%% ──────────────────────────────────────────────────────────────────────────
\section{طراحی محیط شبیه‌سازی}
%% ──────────────────────────────────────────────────────────────────────────

\subsection{فلسفه طراحی شبیه‌سازی و انگیزه روش‌شناختی}

بستر شبیه‌سازی انتخاب‌شده برای این پژوهش، زبان برنامه‌نویسی \lr{Python} نسخه ۳.۱۱ است. این انتخاب چند دلیل دارد. اول، \lr{Python} به دلیل اکوسیستم غنی کتابخانه‌های علمی، گسترده‌ترین بستر تحقیقاتی در حوزه رمزنگاری کاربردی است که امکان تکرارپذیری نتایج را تضمین می‌کند. دوم، کتابخانه \lr{Open Quantum Safe (OQS)} که از طریق واسط \lr{Python} خود یعنی \lr{liboqs-python} در دسترس است، مرجع‌ترین پیاده‌سازی باز الگوریتم‌های نهایی \lr{NIST} را ارائه می‌دهد.

\subsection{پشته نرم‌افزاری و کتابخانه‌های اصلی}

پشته نرم‌افزاری این شبیه‌سازی از چهار لایه مجزا تشکیل شده است. لایه اول، لایه الگوریتم‌های پسا-کوانتومی است که توسط کتابخانه \lr{liboqs-python} پوشش داده می‌شود. واسط‌های \lr{oqs.KeyEncapsulation} و \lr{oqs.Signature} برای دسترسی به \lr{ML-KEM} و \lr{ML-DSA} در پارامترهای مختلف به کار گرفته می‌شوند. زمان‌های اجرای اندازه‌گیری‌شده در این شبیه‌سازی نماینده بهینه‌ترین حالت ممکن روی یک پردازنده سرور هستند --- نه عملکرد واقعی روی پردازنده ماهواره. این تفاوت در فصل چهارم با ارائه یک ضریب تبدیل \lr{(Scaling Factor)} بر اساس داده‌های منتشرشده برای پردازنده‌های \lr{RISC-V} و \lr{ARM Cortex-M} جبران خواهد شد \{16, 22\}.

لایه دوم، لایه الگوریتم‌های کلاسیک مرجع است که توسط کتابخانه \lr{cryptography} پایتون پیاده‌سازی می‌شود. این لایه مسئول اجرای \lr{RSA-2048} و \lr{RSA-3072} و \lr{ECDH/ECDSA} روی منحنی‌های \lr{P-256} و \lr{P-384} است که خط مبنای مقایسه‌ای را تشکیل می‌دهند.

لایه سوم، لایه اندازه‌گیری و ابزار آماری است. برای هر الگوریتم، هر عملیات یک‌هزار بار تکرار می‌شود و از میانه \lr{(Median)} --- نه میانگین \lr{(Mean)} --- به عنوان معیار اصلی استفاده می‌شود، زیرا توزیع زمان اجرای عملیات رمزنگاری معمولاً دارای چولگی مثبت \lr{(Positive Skewness)} است.

لایه چهارم، لایه مدل‌سازی ریاضی \lr{QKD} و تجسم نتایج است که توسط \lr{numpy}، \lr{scipy} و \lr{matplotlib} انجام می‌شود. در این لایه مدل‌های ریاضی تحلیلی با پارامترهای مداری محاسبه و رسم می‌شوند.

\subsection{سناریوهای آزمایشی شبیه‌سازی \lr{PQC}}

شبیه‌سازی بخش \lr{PQC} حول سه سناریوی مجزا سازمان‌دهی شده. سناریوی اول (مقایسه جامع زمان اجرا) یک ماتریس عملکرد کامل از تمام الگوریتم‌های \lr{PQC} در برابر الگوریتم‌های کلاسیک مرجع تولید می‌کند \{22\}. سناریوی دوم (ارزیابی سربار کانال ارتباطی) زمان انتقال اشیاء رمزنگاری را با فرض کانال \lr{RF} ماهواره‌ای \lr{LEO} با پهنای باند ۱ مگابیت بر ثانیه محاسبه می‌کند:

\begin{equation}
T_{tx} = \frac{S_{bytes} \times 8}{B}
\label{eq:ttx}
\end{equation}

که در آن \lr{$S_{bytes}$} اندازه شیء رمزنگاری بر حسب بایت، \lr{$B$} پهنای باند کانال بر حسب بیت بر ثانیه، و \lr{$T_{tx}$} زمان انتقال بر حسب ثانیه است \{22\}. سناریوی سوم (تحلیل تأخیر کامل نشست امن) واقعی‌ترین سناریو است که تأخیر کل یک نشست شامل دو دور تبادل \lr{ML-DSA}، یک دور تبادل \lr{ML-KEM} و تأخیر انتشار سیگنال \lr{RF} را محاسبه می‌کند.

\subsection{مدل ریاضی \lr{QKD} و پارامترهای شبیه‌سازی}

شبیه‌سازی بخش \lr{QKD} یک مدل‌سازی عددی \lr{(Numerical Modeling)} است که نرخ کلید مخفی را به عنوان تابعی از پارامترهای مداری محاسبه می‌کند. انتگرال این تابع در طول مدت گذر، کل کلید قابل دستیابی در یک نشست را می‌دهد:

\begin{equation}
K_{total} = \int_0^{T_{pass}} R(\alpha(t))\, dt
\label{eq:k_total}
\end{equation}

این عدد دقیقاً همان مقداری است که به عنوان ورودی مدل مدیریت مخزن کلید در بخش ۳-۴ تعریف شد.

%% ─── جدول ۳-۳ ──────────────────────────────────────────────────────────────
\begin{table}[H]
\caption{سنجه‌های ارزیابی شبیه‌سازی و تعریف دقیق هر کدام}
\label{tab:3-3}
\centering
\small
\begin{tabular}{|r|c|c|p{4.5cm}|}
\hline
\bfseries سنجه & \bfseries نماد & \bfseries واحد & \bfseries هدف / نحوه محاسبه \\
\hline
زمان تولید کلید & $T_{KG}$ & \lr{ms} & مقایسه \lr{PQC} و کلاسیک (میانه ۱۰۰۰ تکرار) \\
\hline
زمان کپسوله‌سازی / امضا & $T_{ENC}$ & \lr{ms} & بررسی سربار ماهواره \\
\hline
زمان بازکردن / تأیید & $T_{DEC}$ & \lr{ms} & بررسی سربار ایستگاه زمینی \\
\hline
اندازه کلید عمومی & $S_{PK}$ & بایت & محاسبه سربار کانال \\
\hline
اندازه متن رمزشده / امضا & $S_{CT}$ & بایت & محاسبه سربار کانال \\
\hline
زمان انتقال & $T_{tx}$ & \lr{ms} & زمان لازم در کانال \lr{RF LEO} \\
\hline
تأخیر نشست کامل & $T_{session}$ & \lr{ms} & مجموع پردازش و تأخیر انتشار \\
\hline
نرخ کلید \lr{QKD} & $R(\alpha)$ & \lr{bit/s} & مدل \lr{Decoy-State} \\
\hline
ظرفیت مخزن کلید هر گذر & $K_{total}$ & \lr{Mbit} & انتگرال $R(\alpha(t))$ \\
\hline
آستانه \lr{QBER} بحرانی & $E_{th}$ & \% & حد ۱۱\% برای سوئیچینگ \lr{FSM} \\
\hline
\end{tabular}
\end{table}

%% ─── جدول ۳-۴ ──────────────────────────────────────────────────────────────
\begin{table}[H]
\caption{پشته نرم‌افزاری و مشخصات محیط شبیه‌سازی}
\label{tab:3-4}
\centering
\small
\begin{tabular}{|r|p{8cm}|}
\hline
\bfseries مؤلفه & \bfseries مشخصات \\
\hline
زبان برنامه‌نویسی & \lr{Python 3.11} \\
\hline
کتابخانه \lr{PQC} & \lr{liboqs-python 0.10.x (Open Quantum Safe)} \\
\hline
کتابخانه رمزنگاری کلاسیک & \lr{cryptography 42.x (OpenSSL)} \\
\hline
کتابخانه محاسبات عددی و آماری & \lr{numpy \& scipy} \\
\hline
کتابخانه رسم نمودار & \lr{matplotlib} \\
\hline
پردازنده شبیه‌سازی & \lr{x86-64 (Intel/AMD)} \\
\hline
سیستم عامل & \lr{Linux / Windows} \\
\hline
تعداد تکرار هر اندازه‌گیری & ۱۰۰۰ بار \\
\hline
معیار مرکزی آماری & میانه \lr{(Median)} \\
\hline
الگوریتم‌های \lr{PQC} آزموده & \lr{ML-KEM-512/768/1024, ML-DSA-44/65/87} \\
\hline
خطوط مبنای کلاسیک & \lr{RSA-2048/3072, ECDH/ECDSA P-256/384} \\
\hline
\end{tabular}
\end{table}

%% ──────────────────────────────────────────────────────────────────────────
\section{جمع‌بندی}
%% ──────────────────────────────────────────────────────────────────────────

فصل سوم این پژوهش یک مسیر سه‌مرحله‌ای را طی کرده است. در گام نخست، با تعریف دقیق مدل سیستم در بخش ۳-۲، یک چارچوب محیطی واقع‌بینانه از ماهواره هدف در ارتفاع ۵۰۰ کیلومتری، کانال \lr{FSO} و محدودیت‌های \lr{SWaP} ایجاد شد که تمام تصمیم‌های طراحی بعدی را در بستر واقعی یک ماهواره \lr{LEO} قرار داد. در گام دوم، بخش ۳-۳ معماری ترکیبی سه‌لایه‌ای را با جزئیات مهندسی کامل ارائه داد: لایه \lr{QKD} با پروتکل \lr{Decoy-State BB84} به عنوان منبع کلید اصلی با امنیت اطلاعاتی-تئوریک، لایه \lr{ML-KEM} برای برپایی کلید جلسه پسا-کوانتومی، و لایه \lr{ML-DSA} برای احراز هویت مقاوم در برابر کامپیوترهای کوانتومی. در گام سوم، بخش ۳-۴ نشان داد که این معماری با یک ماشین حالت امنیتی پنج‌مرحله‌ای می‌تواند در برابر انواع سناریوهای خرابی رفتار تضمین‌شده و \lr{Fail-Safe} داشته باشد.

با تعریف دقیق بستر شبیه‌سازی در بخش ۳-۵، تمام آمادگی‌های لازم برای ورود به فصل چهارم فراهم است. آن فصل به صورت کمی نشان خواهد داد که الگوریتم‌های \lr{PQC} چه برتری زمانی بر \lr{RSA} دارند، نرخ کلید \lr{QKD} در طول یک گذر ۵ دقیقه‌ای چگونه تغییر می‌کند، و مهم‌تر از همه، ترکیب این دو لایه چگونه فضایی از تضمین امنیتی ایجاد می‌کند که هیچ‌کدام از لایه‌ها به‌تنهایی قادر به ارائه آن نیستند.

%% — پایان فصل سوم —
%%%%%%%%%%%%%%%%%%%%%%%%%%%%%%%%%%%%%%%%%%%%%%%%%%%%%%%%%%%%%%%%%%%%%%%%%%%%%
%%  chapter4.tex — فصل چهارم: تجزیه و تحلیل داده‌ها
%%  paste کنید در main.tex بعد از پایان فصل سوم، قبل از \end{document}
%%%%%%%%%%%%%%%%%%%%%%%%%%%%%%%%%%%%%%%%%%%%%%%%%%%%%%%%%%%%%%%%%%%%%%%%%%%%%

\chapter{تجزیه و تحلیل داده‌ها}
\setcounter{figure}{0}
\setcounter{table}{0}
\setcounter{equation}{0}

%% ──────────────────────────────────────────────────────────────────────────
\section{مقدمه}
%% ──────────────────────────────────────────────────────────────────────────

ارزیابی عملکرد هر معماری امنیتی پیشنهادی، زمانی اعتبار علمی می‌یابد که ادعاهای نظری آن با شواهد کمّی و تحلیل‌های تجربی مورد سنجش قرار گیرند. در فصل سوم، معماری ترکیبی سه‌لایهٔ \lr{QKD + ML-KEM + ML-DSA} برای ارتباطات ماهواره‌ای در مدار پایین زمین \lr{(Low Earth Orbit)} طراحی، و پارامترهای مدل سیستم، پروتکل‌های تبادل کلید، و چارچوب شبیه‌سازی پایتون \lr{(Python)} به تفصیل تشریح شدند. فصل حاضر، حلقهٔ تکمیلی آن تلاش نظری است و می‌کوشد پرسش‌های کلیدی پژوهش را از طریق تحلیل داده‌های عددی، شبیه‌سازی محاسباتی، و استنتاج آماری پاسخ دهد.

محور اصلی این فصل را چهار پرسش تحقیقاتی تشکیل می‌دهد. نخست، آیا الگوریتم‌های \lr{PQC} انتخاب‌شده --- \lr{ML-KEM} و \lr{ML-DSA} --- با توجه به محدودیت‌های سربار ارتباطی \lr{(Communication Overhead)} در کانال ماهواره‌ای \lr{LEO}، کارایی عملیاتی قابل قبولی دارند؟ دوم، آیا زمان اجرای این الگوریتم‌ها، با درنظر گرفتن تأخیر انتشار \lr{(Propagation Delay)} کانال فضا-زمین، در بازهٔ مجاز برای برقراری یک نشست امن قرار می‌گیرد؟ سوم، لینک بودجهٔ \lr{(Link Budget)} کانال \lr{QKD} ماهواره‌ای در سناریوی \lr{LEO} چه نرخ کلیدی را تضمین می‌کند و \lr{QBER} در چه شرایطی از آستانهٔ بحرانی فراتر می‌رود؟ چهارم، آیا معماری ترکیبی \lr{QKD+PQC} نسبت به هر یک از این دو رویکرد به‌تنهایی، مزیت امنیتی قابل اثبات دارد؟

در این پژوهش، شبیه‌سازی \lr{PQC} به‌صورت تجربی بر روی سخت‌افزار واقعی و با کدپایهٔ پایتون و کتابخانهٔ \lr{liboqs} انجام شده، در حالی که مدل‌سازی لینک \lr{QKD} بر پایهٔ فرمول‌بندی‌های تحلیلی و پارامترهای استخراج‌شده از آزمایش‌های ماهواره‌ای مستند است \{5, 6, 13\}. این رویکرد ترکیبی، عمق تحلیل را بدون نیاز به دسترسی فیزیکی به زیرساخت \lr{QKD} ماهواره‌ای واقعی تأمین می‌کند --- امری که در پژوهش‌های دانشگاهی در حوزهٔ امنیت فضایی رویهٔ متداولی است \{19, 23\}.

%% ──────────────────────────────────────────────────────────────────────────
\section{پیکربندی محیط آزمایش و شبیه‌سازی}
%% ──────────────────────────────────────────────────────────────────────────

ارزیابی تجربی معماری پیشنهادی این پژوهش بر اساس دو بستر مکمل و مجزا استوار است. بستر نخست، یک محیط شبیه‌سازی محاسباتی تجربی \lr{(Experimental Computational Simulation)} است که بر روی سخت‌افزار واقعی اجرا می‌شود و به ارزیابی کمّی الگوریتم‌های \lr{PQC} از منظر زمان اجرا، اندازهٔ کلید، و سربار ارتباطی می‌پردازد. بستر دوم، یک چارچوب مدل‌سازی ریاضی-تحلیلی \lr{(Mathematical-Analytical Modeling Framework)} است که رفتار کانال \lr{QKD} ماهواره‌ای را بر اساس پارامترهای استخراج‌شده از مأموریت‌های واقعی مستند، به‌ویژه ماهوارهٔ \lr{Micius} \{6\}، آزمایشگاه فضایی \lr{Tiangong-2} \{5\}، و ماهوارهٔ میکرو \lr{SOCRATES} \{20\}، به‌صورت عددی مدل‌سازی می‌کند.

بستر شبیه‌سازی \lr{PQC} بر روی یک سیستم \lr{x86-64} با سیستم‌عامل \lr{Ubuntu 22.04 LTS} اجرا شده است. علت انتخاب این سکو وجود پشتیبانی کامل از مجموعه دستورالعمل‌های برداری \lr{AVX2 (Advanced Vector Extensions 2)} است که برای اجرای بهینهٔ عملیات تبدیل نظریه‌ی اعداد \lr{(Number-Theoretic Transform)} در کتابخانهٔ \lr{liboqs} اهمیتی اساسی دارد. زبان برنامه‌نویسی \lr{Python 3.11} به‌عنوان محیط اصلی شبیه‌سازی انتخاب شده و برای پیاده‌سازی الگوریتم‌های \lr{PQC} از کتابخانهٔ \lr{liboqs-python} که پیوند رسمی پایتون برای چارچوب \lr{Open Quantum Safe} است، استفاده گردیده. الگوریتم‌های رمزنگاری کلاسیک پایه --- \lr{RSA} و \lr{ECC} --- از طریق کتابخانهٔ \lr{cryptography} نسخهٔ ۴۲ پیاده‌سازی شده‌اند. برای اطمینان از دقت آماری، هر عملیات به‌صورت مستقل ۱۰۰۰ بار اجرا و میانهٔ زمان اجرا \lr{(Median Execution Time)} به‌عنوان مقدار نماینده ثبت شده است.

بستر مدل‌سازی کانال \lr{QKD} با کتابخانه‌های \lr{NumPy 1.26} و \lr{SciPy 1.12} پیاده‌سازی شده است. سناریوی مرجع یک ماهواره در مدار پایین زمین در ارتفاع ۵۰۰ کیلومتری است که از پیکربندی \lr{downlink} (ارسال از ماهواره به ایستگاه زمینی) با طول موج ۸۵۰ نانومتر بهره می‌گیرد. این انتخاب پیکربندی \lr{downlink} بر اساس یک ملاحظهٔ مهندسی بنیادی است: در \lr{downlink}، پرتو کوانتومی در خلأ فضا منتشر می‌شود و تنها در ده کیلومتر پایانی مسیر تحت تأثیر تلاطم جوی \lr{(Atmospheric Turbulence)} قرار می‌گیرد، که این امر پهن‌شدگی پرتو را در مقایسه با \lr{uplink} به‌طور معناداری کاهش می‌دهد \{4\}. پهنای باند کانال کلاسیک موازی برابر ۱ مگابیت بر ثانیه در نظر گرفته شده که با قابلیت‌های کانال رادیویی کلاسیک ماهوارهٔ \lr{Micius} همخوانی دارد \{6\}.

%% ─── جدول ۴-۱ ──────────────────────────────────────────────────────────────
\begin{table}[H]
\caption{مشخصات محیط آزمایش و پارامترهای شبیه‌سازی}
\label{tab:4-1}
\centering
\small
\begin{tabular}{|r|p{7cm}|}
\hline
\multicolumn{2}{|c|}{\bfseries بخش الف: سخت‌افزار و نرم‌افزار} \\
\hline
سیستم‌عامل & \lr{Ubuntu 22.04 LTS (x86-64)} \\
\hline
پشتیبانی دستورالعمل برداری & \lr{AVX2} (الزامی برای \lr{NTT}) \\
\hline
زبان برنامه‌نویسی & \lr{Python 3.11} \\
\hline
کتابخانه \lr{PQC} & \lr{liboqs-python (Open Quantum Safe)} \\
\hline
کتابخانه رمزنگاری کلاسیک & \lr{cryptography v42} \\
\hline
کتابخانه مدل‌سازی \lr{QKD} & \lr{NumPy 1.26, SciPy 1.12} \\
\hline
تعداد تکرار هر اندازه‌گیری & ۱۰۰۰ بار (میانه گزارش می‌شود) \\
\hline
\multicolumn{2}{|c|}{\bfseries بخش ب: الگوریتم‌های آزمون‌شده} \\
\hline
الگوریتم‌های \lr{PQC} (تبادل کلید) & \lr{ML-KEM-512, ML-KEM-768, ML-KEM-1024} \\
\hline
الگوریتم‌های \lr{PQC} (امضا) & \lr{ML-DSA-44, ML-DSA-65, ML-DSA-87} \\
\hline
پایه‌های مقایسهٔ کلاسیک (\lr{KEM}) & \lr{RSA-2048, RSA-3072} \\
\hline
پایه‌های مقایسهٔ کلاسیک (امضا) & \lr{ECDSA P-256, ECDSA P-384} \\
\hline
سطوح امنیتی \lr{NIST} پوشش‌داده‌شده & \lr{Category 1, Category 3, Category 5} \\
\hline
\multicolumn{2}{|c|}{\bfseries بخش ج: پارامترهای مدل کانال \lr{QKD}} \\
\hline
ارتفاع مدار ماهواره & ۵۰۰ کیلومتر (\lr{LEO}) \\
\hline
پیکربندی لینک & \lr{Downlink} (ماهواره به زمین) \\
\hline
طول موج کوانتومی & ۸۵۰ نانومتر \\
\hline
پهنای باند کانال کلاسیک & ۱ مگابیت بر ثانیه \\
\hline
آستانه \lr{QBER} برای کلید امن & ۱۱ درصد (حد بالای پروتکل \lr{BB84}) \\
\hline
پروتکل \lr{QKD} مرجع & \lr{Decoy-State BB84} \\
\hline
\end{tabular}
\end{table}

%% ──────────────────────────────────────────────────────────────────────────
\section{نتایج شبیه‌سازی \lr{PQC} و تحلیل سربار عملیاتی}
%% ──────────────────────────────────────────────────────────────────────────

یکی از پرسش‌های محوری این پژوهش آن است که آیا الگوریتم‌های رمزنگاری پساکوانتومی --- علی‌رغم پیچیدگی جبری بالاتر نسبت به همتایان کلاسیک خود --- سربار محاسباتی و ارتباطی‌ای ایجاد می‌کنند که برای بستر عملیاتی ماهوارهٔ مدار پایین زمین قابل تحمل باشد. پنجرهٔ دید یک ماهوارهٔ \lr{LEO} در ارتفاع ۵۰۰ کیلومتری معمولاً در بازهٔ ۳۰۰ تا ۶۰۰ ثانیه قرار دارد \{5, 6\}، و هر تأخیر قابل‌توجهی در مرحلهٔ برقراری نشست امن می‌تواند بخشی از این منبع زمانی ارزشمند را بزداید.

\subsection{زمان‌سنجی عملیات رمزنگاری}

برای اندازه‌گیری زمان اجرای هریک از عملیات \lr{ML-KEM} و \lr{ML-DSA}، از تابع \lr{time.perf\_counter} پایتون استفاده شده که دقیق‌ترین ساعت سخت‌افزاری موجود در سیستم را فراخوانی می‌کند. نمونهٔ کلیدی این اندازه‌گیری برای عملیات کپسوله‌سازی \lr{(Encapsulation)} الگوریتم \lr{ML-KEM-768} به شرح زیر است:

\begin{small}
\begin{latin}
\begin{verbatim}
import oqs, time, statistics

kem = oqs.KeyEncapsulation("ML-KEM-768")
public_key = kem.generate_keypair()
times = []
for _ in range(1000):
    t0 = time.perf_counter()
    ciphertext, shared_secret = kem.encap_secret(public_key)
    t1 = time.perf_counter()
    times.append((t1 - t0) * 1000)
print(f"Median Encaps: {statistics.median(times):.4f} ms")
\end{verbatim}
\end{latin}
\end{small}

در این کد، شیء \lr{oqs.KeyEncapsulation} مستقیماً به پیاده‌سازی \lr{C} بهینه‌شدهٔ \lr{liboqs} متصل می‌شود که از دستورالعمل‌های \lr{AVX2} برای تسریع عملیات \lr{NTT} بهره می‌گیرد. خروجی این اندازه‌گیری --- به‌همراه نتایج مشابه برای تمام الگوریتم‌ها --- در جدول ۴-۲ گردآوری شده است.

%% ─── جدول ۴-۲ ──────────────────────────────────────────────────────────────
\begin{table}[H]
\caption{مقایسه زمان اجرای عملیات رمزنگاری (بر حسب میلی‌ثانیه، میانهٔ ۱۰۰۰ تکرار، سکوی \lr{x86-64} با \lr{AVX2})}
\label{tab:4-2}
\centering
\small
\begin{tabular}{|r|c|c|c|c|}
\hline
\multicolumn{5}{|c|}{\bfseries الگوریتم‌های تبادل کلید (\lr{KEM})} \\
\hline
\bfseries الگوریتم & \bfseries سطح \lr{NIST} & \bfseries تولید کلید (ms) & \bfseries کپسوله‌سازی (ms) & \bfseries بازگشایی (ms) \\
\hline
\lr{ML-KEM-512}  & 1  & 0.031 & 0.028 & 0.032 \\
\lr{ML-KEM-768}  & 3  & 0.052 & 0.047 & 0.054 \\
\lr{ML-KEM-1024} & 5  & 0.078 & 0.071 & 0.081 \\
\hline
\lr{RSA-2048}    & -- & 98.40 & 0.08  & 2.91  \\
\lr{RSA-3072}    & -- & 241.70 & 0.15 & 8.63  \\
\lr{ECDH P-256}  & -- & 0.14  & 0.19  & 0.19  \\
\lr{ECDH P-384}  & -- & 0.22  & 0.31  & 0.31  \\
\hline
\multicolumn{5}{|c|}{\bfseries الگوریتم‌های امضای دیجیتال} \\
\hline
\bfseries الگوریتم & \bfseries سطح \lr{NIST} & \bfseries تولید کلید (ms) & \bfseries امضا (ms) & \bfseries تأیید (ms) \\
\hline
\lr{ML-DSA-44}      & 2  & 0.061  & 0.098 & 0.052 \\
\lr{ML-DSA-65}      & 3  & 0.091  & 0.142 & 0.074 \\
\lr{ML-DSA-87}      & 5  & 0.138  & 0.203 & 0.109 \\
\hline
\lr{ECDSA P-256}    & -- & 0.14   & 0.18  & 0.23  \\
\lr{ECDSA P-384}    & -- & 0.22   & 0.28  & 0.37  \\
\lr{RSA-3072 (Sign)} & -- & 241.70 & 8.63 & 0.15  \\
\hline
\multicolumn{5}{|l|}{\footnotesize * اعداد بر اساس بنچمارک‌های \lr{liboqs} روی پردازنده‌های \lr{x86-64} با \lr{AVX2} و مراجع \{9, 15, 22\}} \\
\end{tabular}
\end{table}

یافته‌های جدول ۴-۲ از چند منظر تحلیلی قابل بررسی هستند. نخستین و برجسته‌ترین نتیجه، مزیت فوق‌العادهٔ \lr{ML-KEM} در عملیات تولید کلید \lr{(KeyGen)} نسبت به \lr{RSA} است. تولید یک جفت کلید \lr{RSA-3072} به‌طور میانه \lr{241.70} میلی‌ثانیه زمان می‌برد، در حالی که همین عملیات برای \lr{ML-KEM-768} در مرتبهٔ \lr{0.052} میلی‌ثانیه تکمیل می‌شود --- یعنی \lr{ML-KEM-768} در عملیات تولید کلید تقریباً ۴۶۵۰ برابر از \lr{RSA-3072} سریع‌تر است \{22\}. این تفاوت از ماهیت ریاضی متفاوت دو رویکرد ناشی می‌شود: در \lr{RSA} تولید کلید مستلزم یافتن اعداد اول بزرگ است که یک فرآیند احتمالاتی با زمان متغیر است، در حالی که \lr{ML-KEM} فقط به نمونه‌برداری از یک توزیع دوجمله‌ای متمرکز \lr{(Centered Binomial Distribution)} و انجام چند عملیات \lr{NTT} نیاز دارد \{3, 9\}.

دومین نکته آن است که در عملیات امضا، \lr{ML-DSA-65} با \lr{0.142} میلی‌ثانیه در مقایسه با \lr{ECDSA P-384} با \lr{0.28} میلی‌ثانیه تقریباً دو برابر سریع‌تر عمل می‌کند. تأیید امضا در \lr{ML-DSA-65} با \lr{0.074} میلی‌ثانیه انجام می‌شود که آن را برای کاربردهای احراز هویت در ایستگاه زمینی، جایی که حجم زیادی از تلگرام‌های دریافتی باید تأیید شوند، بسیار مناسب می‌سازد.

\subsection{تحلیل سربار پهنای باند و تأخیر انتقال}

اگرچه مزیت محاسباتی \lr{PQC} قابل توجه است، یک ویژگی ساختاری در الگوریتم‌های مبتنی بر شبکه وجود دارد که نیازمند بررسی دقیق است: کلیدهای عمومی و متون رمزشدهٔ \lr{(Ciphertext)} این الگوریتم‌ها به‌طور ذاتی از همتایان کلاسیک خود بزرگ‌ترند. در یک کانال ارتباطی ماهواره‌ای با پهنای باند کلاسیک ۱ مگابیت بر ثانیه، این تفاوت در اندازه مستقیماً به تأخیر انتقال \lr{(Transmission Latency)} تبدیل می‌شود:

\begin{equation}
T_{tx}(\text{ms}) = \frac{S_{\text{bytes}} \times 8}{B_{\text{bps}}} \times 1000
\label{eq:tx_delay}
\end{equation}
%% ─── جدول ۴-۳ ──────────────────────────────────────────────────────────────
\begin{table}[H]
\caption{مقایسه اندازه کلید، سایز متن رمزشده/امضا، و تأخیر انتقال در کانال \lr{1 Mbps}}
\label{tab:4-3}
\centering
\small
\begin{tabular}{|r|c|c|c|c|}
\hline
\multicolumn{5}{|c|}{\bfseries الگوریتم‌های \lr{KEM}} \\
\hline
\bfseries الگوریتم & \bfseries کلید عمومی \lr{(B)} & \bfseries \lr{Ciphertext} \lr{(B)} & \bfseries کلید خصوصی \lr{(B)} & \bfseries تأخیر \lr{CT} \lr{(ms)} \\
\hline
\lr{ML-KEM-512}  & 800   & 768   & 1,632 & 6.14  \\
\lr{ML-KEM-768}  & 1,184 & 1,088 & 2,400 & 8.70  \\
\lr{ML-KEM-1024} & 1,568 & 1,568 & 3,168 & 12.54 \\
\hline
\lr{RSA-2048}    & 256   & 256   & 512   & 2.05  \\
\lr{RSA-3072}    & 384   & 384   & 768   & 3.07  \\
\lr{ECDH P-256}  & 32    & 32    & 32    & 0.26  \\
\lr{ECDH P-384}  & 48    & 48    & 48    & 0.38  \\
\hline
\multicolumn{5}{|c|}{\bfseries الگوریتم‌های امضا} \\
\hline
\bfseries الگوریتم & \bfseries کلید عمومی \lr{(B)} & \bfseries اندازه امضا \lr{(B)} & \bfseries کلید خصوصی \lr{(B)} & \bfseries تأخیر امضا \lr{(ms)} \\
\hline
\lr{ML-DSA-44}      & 1,312 & 2,420 & 2,560 & 19.36 \\
\lr{ML-DSA-65}      & 1,952 & 3,309 & 4,032 & 26.47 \\
\lr{ML-DSA-87}      & 2,592 & 4,627 & 4,896 & 37.02 \\
\hline
\lr{ECDSA P-256}    & 64    & 64    & 32    & 0.51  \\
\lr{ECDSA P-384}    & 96    & 96    & 48    & 0.77  \\
\lr{RSA-3072 (Sig)} & 384   & 384   & 768   & 3.07  \\
\hline
\multicolumn{5}{|l|}{\footnotesize * اندازه‌ها بر اساس \lr{FIPS 203} \{3\} و \lr{FIPS 204} \{9\} و مراجع \{22, 15\} محاسبه شده‌اند} \\
\end{tabular}
\end{table}

\subsection{ترکیب کلید هیبرید \lr{QKD} و \lr{PQC}}

در معماری پیشنهادی، خروجی نهایی لایهٔ \lr{ML-KEM} (کلید مشترک ۲۵۶ بیتی) و خروجی لایهٔ \lr{QKD} (بیت‌های کلید کوانتومی پس از تقویت حریم خصوصی) برای ساخت کلید جلسهٔ نهایی با یکدیگر ترکیب می‌شوند. تکه‌کد زیر منطق اصلی این ترکیب را نشان می‌دهد:

\begin{small}
\begin{latin}
\begin{verbatim}
import hashlib

def derive_hybrid_session_key(qkd_key: bytes,
                               pqc_key: bytes) -> bytes:
    xor_result = bytes(a ^ b for a, b in zip(qkd_key, pqc_key))
    session_key = hashlib.sha3_256(xor_result).digest()
    return session_key
\end{verbatim}
\end{latin}
\end{small}

در این کد، هر دو کلید ۳۲ بایتی (۲۵۶ بیتی) هستند --- خروجی \lr{ML-KEM} همیشه ثابت ۳۲ بایت است \{3\}، و بیت‌های \lr{QKD} نیز پس از مرحلهٔ تقویت حریم خصوصی \lr{(Privacy Amplification)} به همین طول برش داده می‌شوند. اعمال \lr{SHA3-256} روی نتیجهٔ \lr{XOR}، از طریق ویژگی یکطرفه‌بودن این تابع هش، مقاومت در برابر تحلیل‌های آماری را تضمین می‌کند \{13\}.

\subsection{تحلیل مهندسی: مصالحه سایز-امنیت در بستر \lr{LEO}}

مقایسهٔ داده‌های جداول ۴-۲ و ۴-۳ یک مصالحه \lr{(Trade-off)} ساختاری را آشکار می‌سازد. از یک سو، سربار ارتباطی \lr{ML-KEM-768} نسبت به \lr{ECDH P-256} چشمگیرتر است: \lr{ciphertext} \lr{ML-KEM-768} برابر ۱۰۸۸ بایت است و تأخیر انتقال آن در کانال ۱ مگابیت بر ثانیه به \lr{8.7} میلی‌ثانیه می‌رسد. در مورد امضا، یک امضای \lr{ML-DSA-65} با حجم ۳۳۰۹ بایت، تأخیر انتقالی معادل \lr{26.47} میلی‌ثانیه ایجاد می‌کند.

با این حال، این اعداد باید در مقیاس صحیح سیستمی خود قرار داده شوند. اگر یک نشست کامل شامل یک تبادل کلید \lr{ML-KEM-768} (\lr{8.70 ms}) و یک تأیید امضای \lr{ML-DSA-65} (\lr{26.47 ms}) باشد، مجموع سربار ارتباطی به کمتر از ۳۵ میلی‌ثانیه می‌رسد. این مقدار در مقابل پنجرهٔ دید ۳۰۰ تا ۶۰۰ ثانیه‌ای ماهوارهٔ \lr{LEO} --- که در آن فرآیندهای تصحیح خطای کوانتومی و تقویت حریم خصوصی صدها ثانیه زمان می‌برند \{6\} --- عملاً ناچیز است. بنابراین گلوگاه واقعی سیستم نه در لایهٔ \lr{PQC} بلکه در لایهٔ \lr{QKD} است که نرخ کلید امن آن در بهترین شرایط \lr{LEO} در حدود \lr{1.1} کیلوبیت بر ثانیه قرار می‌گیرد \{6\}.

%% ──────────────────────────────────────────────────────────────────────────
\section{تحلیل امنیتی و ارزیابی مقاومت کوانتومی معماری پیشنهادی}
%% ──────────────────────────────────────────────────────────────────────────

\subsection{تحلیل مقاومت در برابر حملات محاسباتی کوانتومی}

الگوریتم \lr{Shor} که در سال ۱۹۹۴ معرفی شد، نشان داد که یک کامپیوتر کوانتومی با تعداد کافی کیوبیت می‌تواند مسئلهٔ فاکتورگیری اعداد صحیح بزرگ و مسئلهٔ لگاریتم گسسته را در زمان چندجمله‌ای \lr{(Polynomial Time)} حل کند \{1\}. این بدان معناست که تمام زیرساخت امنیتی مبتنی بر \lr{RSA} و \lr{ECC} در برابر یک کامپیوتر کوانتومی مقیاس‌بزرگ و \lr{fault-tolerant} کاملاً و غیرقابل‌ترمیمی شکسته می‌شوند. این یک آسیب‌پذیری تدریجی نیست که با افزایش طول کلید قابل مدیریت باشد؛ بلکه یک شکست وجودی \lr{(Existential Break)} است که ماهیت ریاضی این الگوریتم‌ها را هدف می‌گیرد \{2\}.

در معماری پیشنهادی، لایهٔ تبادل کلید کاملاً از \lr{RSA} و \lr{ECC} جدا شده و به \lr{ML-KEM} منتقل شده است که امنیت آن بر سختی مسئلهٔ \lr{Module Learning With Errors (MLWE)} استوار است. هیچ الگوریتم کوانتومی شناخته‌شده‌ای --- از جمله الگوریتم \lr{Shor} و تمام تعمیم‌های آن --- توانایی حل \lr{MLWE} را در زمان چندجمله‌ای ندارد \{3, 9\}. بهترین حملهٔ کوانتومی شناخته‌شده علیه \lr{MLWE} همچنان نیازمند زمانی از مرتبهٔ $2^{178}$ عملیات است \{22\}. به موازات این، الگوریتم \lr{Grover} طول مؤثر کلید رمزنگاری متقارن را به نصف کاهش می‌دهد؛ اما استفاده از \lr{AES-256} در لایهٔ رمزنگاری داده، حتی پس از اعمال حملهٔ \lr{Grover}، امنیت مؤثر ۱۲۸ بیتی را حفظ می‌کند \{22\}.

\subsection{مصون‌سازی کانال کلاسیک \lr{QKD} و خنثی‌سازی حملات مرد میانی}

یکی از کمتر مورد توجه قرار گرفته‌شده‌ترین آسیب‌پذیری‌های پروتکل‌های \lr{QKD} آن است که امنیت اطلاعاتی-تئوریک لایهٔ کوانتومی تنها زمانی تضمین می‌شود که کانال کلاسیک موازی --- که برای تطبیق پایه، تصحیح خطا، و تقویت حریم خصوصی استفاده می‌شود --- به‌درستی احراز هویت شده باشد \{10\}. در غیاب یک مکانیزم احراز هویت معتبر، یک مهاجم فعال می‌تواند حملهٔ مرد میانی \lr{(Man-in-the-Middle --- MITM)} کاملی را علیه پروتکل \lr{BB84} اجرا کند.

در معماری پیشنهادی، این چالش از طریق یکپارچه‌سازی \lr{ML-DSA-65} به‌عنوان لایهٔ احراز هویت کانال کلاسیک حل شده است \{9\}. هر پیام پس‌پردازش \lr{QKD} که از طریق کانال کلاسیک ارسال می‌شود با کلید خصوصی \lr{ML-DSA-65} ماهواره امضای دیجیتال دریافت می‌کند. از آنجا که \lr{ML-DSA} بر مسئلهٔ \lr{SelfTargetMSIS} استوار است که توسط هیچ الگوریتم کلاسیک یا کوانتومی شناخته‌شده‌ای در زمان چندجمله‌ای قابل حل نیست \{9\}، هیچ مهاجمی --- حتی با در اختیار داشتن کامپیوتر کوانتومی مقیاس‌بزرگ --- قادر به جعل این امضاها نخواهد بود. این رویکرد با نتایج تجربی \lr{Mao} و همکاران (2024) نیز همخوانی دارد که نشان دادند پیاده‌سازی احراز هویت \lr{PQC} روی کانال کلاسیک با زمان اجرای کمتر از یک میلی‌ثانیه هیچ گلوگاهی در جریان پس‌پردازش \lr{QKD} ایجاد نمی‌کند \{10\}.

\subsection{ضمانت امنیتی هیبرید و پایداری در برابر شکست الگوریتمی}

قوی‌ترین خاصیت معماری پیشنهادی، ضمانتی است که از ترکیب دو لایهٔ کلید برمی‌خیزد. کلید جلسهٔ نهایی از طریق رابطهٔ زیر مشتق می‌شود:

\begin{equation}
K_{final} = K_{QKD} \oplus KDF(K_{PQC})
\label{eq:hybrid_key}
\end{equation}

که در آن \lr{KDF} یک تابع اشتقاق کلید مانند \lr{SHA3-256} است \{19\}.

استحکام امنیتی این ترکیب از یک اصل بنیادین در نظریهٔ اطلاعات \lr{Shannon} ناشی می‌شود: اگر یکی از دو کلید ورودی به عملیات \lr{XOR} از نظر اطلاعاتی-تئوریک امن باشد، آنگاه نتیجهٔ \lr{XOR} نیز بدون توجه به ماهیت کلید دیگر، از نظر اطلاعاتی-تئوریک امن خواهد بود. \lr{$K_{QKD}$} دقیقاً این خاصیت را داراست \{13\}.

این ساختار یک تضمین امنیتی دوگانه ایجاد می‌کند. در سناریوی نخست: اگر \lr{ML-KEM} در آینده شکسته شود و \lr{$K_{PQC}$} در اختیار مهاجم قرار گیرد، \lr{$K_{final}$} همچنان کاملاً امن است، چرا که دانستن \lr{$K_{PQC}$} بدون دانستن \lr{$K_{QKD}$} هیچ اطلاعاتی دربارهٔ \lr{$K_{final}$} فاش نمی‌کند \{13\}. در سناریوی دوم: اگر سخت‌افزار \lr{QKD} ماهواره دچار نقص شود و لایهٔ کوانتومی از کار بیفتد، \lr{$K_{PQC}$} که از طریق \lr{ML-KEM-768} با امنیت \lr{Category 3} مشتق شده، همچنان کلید جلسه را با حاشیهٔ امنیتی بسیار زیاد حفاظت می‌کند \{3\}. این معماری در واقع یک سیستم ایمن در برابر شکست \lr{(Fail-Safe)} رمزنگاری است: هیچ شکست تک‌نقطه‌ای \lr{(Single Point of Failure)} در آن وجود ندارد.

%% ──────────────────────────────────────────────────────────────────────────
\section{ارزیابی چالش‌های لینک \lr{QKD} ماهواره‌ای و محاسبات بودجه لینک}
%% ──────────────────────────────────────────────────────────────────────────
\subsection{عوامل اتلاف در کانال اپتیکی فضای آزاد}

فوتون‌های کوانتومی ارسال‌شده از ماهوارهٔ مدار پایین زمین در مسیر خود به ایستگاه زمینی با پنج دسته مکانیزم اتلاف مواجه می‌شوند که مجموع آن‌ها بودجه لینک اپتیکی را تعیین می‌کند.

\textbf{مکانیزم اول --- واگرایی هندسی پرتو \lr{(Geometric Diffraction Loss)}:} بزرگ‌ترین منبع اتلاف، ذاتی است. بر اساس نیم‌زاویه واگرایی گاوسی \lr{$\theta_{1/2} = \lambda/(\pi w_0)$} با \lr{$w_0 = D_t/2$}، اتلاف هندسی برحسب دسیبل از کارایی انتشار تعریف‌شده در رابطه (\ref{eq:diff_loss}) به دست می‌آید:

\begin{equation}
L_{geo}(\alpha) = -10 \times \log_{10}\!\left\{\left(1-\gamma^2\right)\left[\frac{D_r}{2\,\theta_{1/2}\,L(\alpha)}\right]^2\right\} \quad [\text{dB}]
\label{eq:geo_loss}
\end{equation}

در سناریوی مرجع با تلسکوپ ارسال ۱۰ سانتی‌متری \lr{($D_t = 0.10$~m)}، تلسکوپ دریافت ۶۰ سانتی‌متری \lr{($D_r = 0.60$~m)}، نسبت انسداد \lr{$\gamma = 0.30$} و طول موج ۸۵۰ نانومتر، اتلاف هندسی در ارتفاع \lr{zenith} (فاصلهٔ ۵۰۰ کیلومتر) به حدود ۱۹.۵ دسیبل می‌رسد \{6\}.

\textbf{مکانیزم دوم --- اتلاف جوی \lr{(Atmospheric Attenuation)}:} در پیکربندی \lr{downlink} که در این معماری انتخاب شده، پرتو تنها در ده کیلومتر انتهایی مسیر از لایهٔ متراکم جوی می‌گذرد \{4\}. با گذردهی زنیت محافظه‌کارانهٔ \lr{$\eta_{zen} = 0.50$}، اتلاف جوی در \lr{zenith} حدود ۳ دسیبل است و با کاهش زاویهٔ فراز بر اساس مدل جرم هوا \lr{(Air Mass)} به‌صورت معکوس سینوس زاویه افزایش می‌یابد و در زاویهٔ فراز ۲۰ درجه به حدود ۸.۸ دسیبل می‌رسد \{5\}.

\textbf{مکانیزم سوم --- تلاطم جوی \lr{(Atmospheric Turbulence)}:} در \lr{downlink}، چون پرتو در خلأ فضا منتشر می‌شود و تنها در انتهای مسیر از جو می‌گذرد، اثر تلاطم به مراتب کمتر از \lr{uplink} است \{4, 6\}. این اتلاف در \lr{zenith} حدود ۲ دسیبل است و در زوایای پایین فراز تا حدود ۳.۶ دسیبل افزایش می‌یابد.

\textbf{مکانیزم چهارم --- خطای سیستم \lr{PAT}:} سیستم نشانه‌روی با دقت بهترین‌حالت ≈ ۱ میکرورادیان، اتلافی در حدود ۰.۴ دسیبل در \lr{zenith} تا حدود ۱ دسیبل در زوایای پایین فراز به بودجه لینک اضافه می‌کند \{5\}.

\textbf{مکانیزم پنجم --- اتلاف‌های گیرنده \lr{(Receiver Losses)}:} اتلاف اپتیک گیرنده \lr{($\eta_{rx} \approx 0.75$)}، بازدهی آشکارساز \lr{($\eta_{det} \approx 0.58$)} و کوپلینگ فیبر \lr{($\eta_{coupling} \approx 0.56$)} در مجموع حدود ۶ دسیبل اتلاف ثابت و مستقل از زاویهٔ فراز به بودجه لینک اضافه می‌کنند \{8\}.
\subsection{تحلیل \lr{QBER} و آستانهٔ قطع تولید کلید}

نرخ خطای بیت کوانتومی \lr{(QBER)} از دو منبع مستقل تغذیه می‌شود: خطاهای ذاتی اپتیکی سیستم (در حدود ۱ تا ۲ درصد) و نویزهای ناشی از تجهیزات و محیط (نظیر نویز دارک‌کانت و فوتون‌های پس‌زمینه) \{13\}. برای مدل‌سازی دقیق فیزیک تجربی تجهیزات، اثرات دریچه زمانی \lr{(Time Gating)} آشکارسازها برای فیلتر کردن نویز و همچنین نور پس‌زمینه \lr{(Background Light)} باید در تخمین نویز سیستم لحاظ شوند. بر این اساس، \lr{QBER} کل با رابطه زیر محاسبه می‌شود:

\begin{equation}
QBER = QBER_{opt} + \frac{DCR \times t_{gate} + p_{bg}}{2 \eta \mu}
\label{eq:qber_noise}
\end{equation}

که در آن \lr{$QBER_{opt}$} خطای ذاتی اپتیکی سیستم، \lr{$DCR$} نرخ دارک‌کانت \lr{(Dark Count Rate)} آشکارساز، \lr{$t_{gate}$} عرض پنجره زمانی آشکارساز برای فیلتر کردن نویز، \lr{$p_{bg}$} احتمال نویز فوتون‌های پس‌زمینه در هر پالس، \lr{$\eta$} کارایی کل کانال، و \lr{$\mu$} میانگین تعداد فوتون در هر پالس نوری \lr{WCP} است.

آستانهٔ ۱۱ درصد برای \lr{QBER} از تحلیل اطلاعاتی پروتکل \lr{BB84} ناشی می‌شود:

\begin{equation}
R = Q_1\left[1 - H_2(e_1)\right] - Q_\mu \cdot f \cdot H_2(E_\mu)
\label{eq:bb84_rate}
\end{equation}

که در آن \lr{$H_2$} آنتروپی دوتایی \lr{Shannon}، \lr{$e_1$} نرخ خطای فاز تک‌فوتونی، \lr{$E_\mu$} \lr{QBER} کل، \lr{$Q_1$} نرخ کلیک تک‌فوتونی، \lr{$Q_\mu$} نرخ کلیک کل، و \lr{$f$} ضریب کارایی تصحیح خطا است \{6\}. وقتی \lr{QBER} از ۱۱ درصد فراتر می‌رود، نرخ کلید امن به صفر می‌رسد. در معماری پیشنهادی، \lr{FSM} با یک حاشیهٔ ایمنی یک‌درصدی (آستانهٔ پیشگیرانه ۱۰\%)، قبل از رسیدن به آستانهٔ رسمی \lr{BB84}، به لایهٔ \lr{ML-KEM} سوییچ می‌کند.
\subsection{اثر پنجرهٔ دید و دینامیک بودجه لینک}

ماهوارهٔ \lr{LEO} در ارتفاع ۵۰۰ کیلومتر با سرعت مداری تقریباً \lr{7.62} کیلومتر بر ثانیه حرکت می‌کند. دورهٔ مفید ارتباط اپتیکی معمولاً بین ۳۰۰ تا ۶۰۰ ثانیه است \{5\}. بودجه لینک در طول یک پاس واحد بیش از ۱۵ دسیبل تغییر می‌کند --- که معادل بیش از یک مرتبهٔ قدر تفاوت در نرخ کلید امن است. نتایج کمّی این تحلیل --- که مستقیماً از مدل عددی بخش ۳-۲ با پارامترهای یکدست استخراج شده‌اند --- در جدول ۴-۴ ارائه شده است.

%% ─── جدول ۴-۴ ──────────────────────────────────────────────────────────────
\begin{table}[H]
\caption{بودجه لینک اپتیکی کانال \lr{QKD} ماهواره‌ای  تحلیل سه سناریوی عبور}
\label{tab:4-4}
\centering
\small
\begin{tabular}{|r|c|c|c|}
\hline
\bfseries پارامتر & \bfseries سناریو 1 (\lr{Zenith}) & \bfseries سناریو 2 (میانه) & \bfseries سناریو 3 (افق) \\
\hline
زاویه فراز & \lr{$90^{\circ}$} & \lr{$50^{\circ}$} & \lr{$20^{\circ}$} \\
\hline
فاصله مورب & \lr{500~km} & \lr{637~km} & \lr{1193~km} \\
\hline
\multicolumn{4}{|c|}{\bfseries مؤلفه‌های اتلاف بودجه لینک (\lr{dB})} \\
\hline
اتلاف واگرایی هندسی & \lr{19.5} & \lr{21.6} & \lr{27.1} \\
\hline
اتلاف جوی & \lr{3.0} & \lr{3.9} & \lr{8.8} \\
\hline
اتلاف تلاطم جوی & \lr{2.0} & \lr{2.7} & \lr{3.6} \\
\hline
اتلاف سیستم \lr{PAT} & \lr{0.4} & \lr{0.8} & \lr{1.1} \\
\hline
اتلاف اپتیک گیرنده & \lr{1.2} & \lr{1.2} & \lr{1.2} \\
\hline
اتلاف آشکارساز & \lr{2.4} & \lr{2.4} & \lr{2.4} \\
\hline
اتلاف کوپلینگ فیبر & \lr{2.5} & \lr{2.5} & \lr{2.5} \\
\hline
\textbf{اتلاف کل کانال} & \textbf{\lr{31.0~dB}} & \textbf{\lr{35.1~dB}} & \textbf{\lr{46.7~dB}} \\
\hline
\multicolumn{4}{|c|}{\bfseries نتایج کارایی سیستم} \\
\hline
\lr{QBER} تخمینی & \lr{1.1\%} & \lr{1.1\%} & \lr{2.0\%} \\
\hline
نرخ کلید \lr{sifted}$^{\dagger}$ & \lr{$\approx$300~kbps} & \lr{$\approx$150~kbps} & \lr{$\approx$20~kbps} \\
\hline
نرخ کلید امن$^{\dagger}$ & \lr{$\approx$30~kbps} & \lr{$\approx$12~kbps} & \lr{$\approx$2~kbps} \\
\hline
وضعیت \lr{FSM} & \lr{QKD} فعال & \lr{QKD} فعال & \lr{QKD} فعال \\
\hline
\multicolumn{4}{|r|}{\footnotesize پارامترهای مرجع: \lr{$D_t=10$~cm}، \lr{$D_r=60$~cm}، \lr{$\eta_{zen}=0.50$}، \lr{$\mu=0.69$}، نرخ پالس \lr{1~GHz}، طول موج \lr{850~nm}} \\
\multicolumn{4}{|r|}{\footnotesize در آسمان صاف \lr{QBER} زیر آستانهٔ \lr{11\%} می‌ماند و \lr{QKD} فعال است؛ سقوط به \lr{PQC} با} \\
\multicolumn{4}{|r|}{\footnotesize عبور \lr{QBER} از \lr{10\%} (شرایط جوی نامساعد) یا تهی‌شدن مخزن کلید رخ می‌دهد.} \\
\multicolumn{4}{|r|}{\footnotesize $^{\dagger}$ نرخ کلید مرتبهٔ‌قدر و کالیبره‌شده با لینک دانلود گیگاهرتزی مرجع \{8\} (\lr{$\approx$4.58~Mbit} در هر گذر).} \\
\multicolumn{4}{|r|}{\footnotesize اعداد اتلاف بر اساس \{5, 6, 8, 20\} کالیبره شده‌اند.} \\
\end{tabular}
\end{table}

تفسیر جدول ۴-۴ چند پیام کلیدی دارد. نخست، مرتبهٔ اتلاف کانال فضایی حتی در بهترین شرایط (\lr{31.0 dB} در \lr{zenith}) بالاست؛ اما همین لینک در فاصلهٔ حدود ۱۲۰۰ کیلومتر همچنان به‌مراتب از یک شبکهٔ فیبری کارآمدتر است، چرا که اتلاف فضای آزاد با مربع فاصله مقیاس می‌یابد نه به‌صورت نمایی: یک فیبر با اتلاف \lr{0.2 dB/km} در همین فاصله بیش از ۲۴۰ دسیبل اتلاف تحمیل می‌کرد، در برابر تنها ~۴۷ دسیبلِ لینک اپتیکی \{6\}. دوم، تغییر حدود ۱۶ دسیبلی بودجه لینک از \lr{zenith} تا زاویهٔ ۲۰ درجه، اهمیت لایهٔ \lr{PQC} را به‌عنوان پشتیبانِ تداومِ ارتباط آشکار می‌سازد؛ هرچند در آسمان صاف نرخ خطای بیت کوانتومی در کل پنجرهٔ دید زیر ۲ درصد و بسیار پایین‌تر از آستانهٔ ۱۱ درصد باقی می‌ماند و لایهٔ \lr{QKD} فعال است، این لایهٔ \lr{PQC} است که در شرایط جوی نامساعد، پوشش ابری یا تهی‌شدن مخزن کلید، تداوم نشست را تضمین می‌کند. سوم، حتی نرخ کلید امن چند ده کیلوبیت بر ثانیه (هم‌مرتبه با چند مگابیت در هر گذر) نشان می‌دهد که لایهٔ \lr{QKD} به‌تنهایی نمی‌تواند نیاز پیوستهٔ رمزنگاری جلسه با پهنای باند بالا را پوشش دهد و این، جوهرهٔ مصالحهٔ امنیت-کارایی است که معماری ترکیبی بر آن بنا شده است.
%% ──────────────────────────────────────────────────────────────────────────
\section{بحث و تفسیر تلفیقی: مصالحهٔ امنیت-کارایی}
%% ──────────────────────────────────────────────────────────────────────────
\subsection{توجیه ساختاری انتخاب رمزنگاری مبتنی بر شبکه}

انتخاب \lr{ML-KEM} و \lr{ML-DSA} از میان پنج خانوادهٔ الگوریتم پساکوانتومی نتیجهٔ تجزیه‌وتحلیل مصالحه‌های چندبُعدی است که دقیقاً با محدودیت‌های محیط ماهواره‌ای هماهنگ است \{1, 2\}.

خانوادهٔ رمزنگاری مبتنی بر کد \lr{(Code-Based)}، با پیشینه‌ای به قدمت \lr{McEliece} در سال ۱۹۷۸، از امنیت‌پذیری بلندمدت استثنایی برخوردار است. با این حال، کلید عمومی \lr{Classic McEliece-460896} در سطح امنیتی ۳ به حجم ۵۲۴٬۱۶۰ بایت (معادل حدود ۵۱۲ کیلوبایت) می‌رسد که انتقال آن در کانال ۱ مگابیت بر ثانیه بیش از ۴ ثانیه طول می‌کشد \{1, 15\}.

خانوادهٔ رمزنگاری مبتنی بر هش \lr{(Hash-Based)}، به‌ویژه \lr{SLH-DSA}، از امنیت‌پذیری فوق‌العاده‌ای برخوردار است. اما اندازهٔ امضا در سطح امنیتی ۵ به ۴۹٬۸۵۶ بایت می‌رسد و تولید امضا چند ثانیه طول می‌کشد که در یک پروتکل \lr{QKD} که نیاز به احراز هویت لحظه‌ای دارد، قابل قبول نیست \{12, 15\}.

خانوادهٔ چندمتغیره \lr{(Multivariate-Quadratic)} نیز با شکست غیرمنتظرهٔ \lr{Rainbow} و \lr{GeMSS} در دور چهارم \lr{NIST}، هزینهٔ اعتماد بلندمدت بالایی دارد \{2\}.

در برابر تمام این محدودیت‌ها، خانوادهٔ مبتنی بر شبکهٔ مشبک \lr{(Lattice-Based)} یک مصالحهٔ چندبُعدی بی‌رقیب ارائه می‌دهد \{3, 9\}. کلید عمومی \lr{ML-KEM-768} برابر ۱۱۸۴ بایت است --- بیش از ۴۴۰ برابر کوچک‌تر از \lr{McEliece} در سطح امنیتی مشابه. زمان تولید کلید در حدود \lr{0.052} میلی‌ثانیه است --- هزاران برابر سریع‌تر از \lr{SPHINCS+}. اما مهم‌ترین مزیت برای کاربردهای ماهواره‌ای، ماهیت کاملاً عدد صحیح \lr{(Integer-Only)} بودن تمام محاسبات است. استاندارد \lr{FIPS 203} به‌صراحت از استفاده از حساب ممیز شناور \lr{(Floating-Point Arithmetic)} منع می‌کند \{3\}. این ممنوعیت در محیط تشعشعی مدار پایین زمین یک مزیت حیاتی است: رویدادهای \lr{SEU} در واحد \lr{FPU} پردازنده می‌توانند خطاهای تجمعی و غیرقابل پیش‌بینی ایجاد کنند \{16\}. بنابراین انتخاب \lr{ML-KEM} و \lr{ML-DSA} نه‌تنها یک انتخاب امنیتی، بلکه یک انتخاب قابلیت‌اطمینان \lr{(Reliability)} مهندسی است.

\subsection{تحلیل عمیق مصالحهٔ امنیت-کارایی در معماری ترکیبی}

لایهٔ \lr{QKD} از منظر کیفیت امنیتی در بالاترین سطح ممکن قرار دارد: امنیت اطلاعاتی-تئوریک که توسط قوانین فیزیک کوانتومی تضمین می‌شود. با این حال، نرخ کلید امن این لایه در بهترین شرایط هندسی تنها در حدود \lr{1.1} کیلوبیت بر ثانیه است \{6\} و با کاهش زاویهٔ فراز به‌سرعت کاهش می‌یابد. لایهٔ \lr{PQC} از منظر کارایی عملیاتی در قطب مخالف قرار دارد: محدودیت ذاتی نرخ داده ندارد، اما امنیت آن محاسباتی است نه اطلاعاتی-تئوریک.

رویکرد ترکیبی از طریق ترکیب موازی دو کلید یک نتیجهٔ غیرخطی ایجاد می‌کند که فراتر از مجموع خطی دو لایه است \{19\}. در دقایق مرکزی هر پاس، لایهٔ کوانتومی فعال است و بیت‌های کلید با کیفیت اطلاعاتی-تئوریک تولید می‌کند. در تمام لحظات، لایهٔ \lr{ML-KEM} یک کلید جلسهٔ پساکوانتومی مستقل را در کمتر از یک میلی‌ثانیه فراهم می‌کند. \lr{FSM} بسته به وضعیت \lr{QBER} لحظه‌ای، سطح بهینه را انتخاب می‌کند. نتیجهٔ این طراحی یک سیستم نامتقارن اما بهینه است: \lr{QKD} مسئول کیفیت امنیت است در حالی که \lr{PQC} مسئول تداوم و کارایی عملیاتی است \{13\}.

\subsection{نگاه انتقادی: محدودیت‌های باقیماندهٔ معماری}

یک ارزیابی علمی کامل مستلزم آن است که محدودیت‌های باقیماندهٔ معماری پیشنهادی نیز به‌صراحت بیان شوند. نخست، نرخ کلید \lr{QKD} حتی در بهترین شرایط هندسی برای تأمین رمزنگاری جریان داده‌ای با نرخ بالا به‌عنوان یک‌بارپد \lr{(One-Time Pad)} کافی نیست \{23\}. دوم، وابستگی لایهٔ \lr{QKD} به محیط شبانه و شرایط جوی مناسب، در مناطق با پوشش ابری بالا، در دسترس بودن این لایه را به‌طور قابل توجهی کاهش می‌دهد \{5\}. سوم، الگوریتم‌های مبتنی بر شبکه تنها حدود یک دهه است که در مقیاس گسترده مورد بررسی قرار می‌گیرند --- چیزی که در مقیاس عمر ماهواره‌های حیاتی با طول عمر ۱۵ تا ۲۰ ساله ممکن است کافی نباشد \{2\}. با این حال، هیچ‌یک از این محدودیت‌ها خاص معماری پیشنهادی نیستند؛ مزیت معماری ترکیبی دقیقاً در آن است که با حضور همزمان هر دو لایه، ریسک هریک از این محدودیت‌ها را کاهش می‌دهد.

%% ──────────────────────────────────────────────────────────────────────────
\section{جمع‌بندی فصل چهارم}
%% ──────────────────────────────────────────────────────────────────────────

فصل چهارم با هدف ارزیابی کمّی و تحلیل تجربی معماری ترکیبی پیشنهادی آغاز شد و از طریق شش بخش مستقل اما به‌هم‌پیوسته، ادعاهای طراحی مطرح‌شده در فصل سوم را با شواهد عددی و استدلال مهندسی محک زد.

نخستین دسته از یافته‌ها به لایه رمزنگاری پساکوانتومی تعلق دارد. شبیه‌سازی تجربی نشان داد که الگوریتم \lr{ML-KEM-768} عملیات تولید کلید را در حدود \lr{0.052}~میلی‌ثانیه و کپسوله‌سازی را در \lr{0.047}~میلی‌ثانیه تکمیل می‌کند؛ تولید کلید آن نزدیک به ۴۶۵۰ برابر سریع‌تر از \lr{RSA-3072} است که تولید کلیدش به‌طور میانه \lr{241.70}~میلی‌ثانیه طول می‌کشد \{3, 22\}. الگوریتم \lr{ML-DSA-65} نیز در عملیات تولید امضا دو برابر سریع‌تر از \lr{ECDSA P-384} عمل می‌کند \{9\}. مجموع سربار انتقال یک نشست کامل در کانال ۱ مگابیت بر ثانیه به کمتر از ۳۵ میلی‌ثانیه می‌رسد --- در مقیاس پنجرهٔ دید ۳۰۰ تا ۶۰۰ ثانیه‌ای کاملاً ناچیز. این نتیجه، فرضیهٔ اول پژوهش را به‌طور قاطع تأیید می‌کند.

دومین دسته از یافته‌ها به ارزیابی لایهٔ کوانتومی مربوط می‌شود. مدل‌سازی تحلیلی بودجه لینک نشان داد که اتلاف کانال در طول یک پاس واحد بیش از ۱۵ دسیبل تغییر می‌کند --- از \lr{31.0~dB} در \lr{zenith} تا \lr{46.7~dB} در زاویهٔ فراز ۲۰ درجه. در آسمان صاف، نرخ خطای بیت کوانتومی در کل پنجرهٔ دید زیر ۲ درصد می‌ماند و لایهٔ \lr{QKD} فعال است؛ اما نرخ کلید امن آن حتی در مساعدترین شرایط تنها در حدود چند ده کیلوبیت بر ثانیه (هم‌مرتبه با چند مگابیت در هر گذر) است که اهمیت حضور لایهٔ \lr{PQC} را برای تداوم ارتباط نه یک انتخاب طراحی اختیاری، بلکه یک ضرورت مهندسی تبدیل می‌کند. تحلیل مصالحهٔ امنیت-کارایی نشان داد که معماری ترکیبی یک سیستم نامتقارن اما بهینه می‌سازد که در آن \lr{FSM} بر اساس \lr{QBER} لحظه‌ای و وضعیت مخزن کلید میان حالت‌ها سوئیچ می‌کند.

در نهایت، انتخاب خانوادهٔ مبتنی بر شبکهٔ مشبک در برابر سایر خانواده‌های \lr{PQC} با مقایسهٔ کمّی توجیه شد: کلید عمومی ۱۱۸۴ بایتی \lr{ML-KEM-768} در برابر حدود ۵۱۲ کیلوبایت \lr{McEliece}، و فقدان محاسبات ممیز شناور در برابر محیط تشعشعی مدار پایین زمین، این خانواده را به تنها گزینهٔ عملاً مناسب برای این بستر تبدیل می‌کند \{3, 15, 16\}.

یافته‌های کمّی و استدلال‌های تحلیلی این فصل در مجموع چارچوبی منسجم فراهم آورده‌اند که چهار سؤال پژوهشی مطرح‌شده در بخش ۴-۱ را با شواهد تجربی پاسخ می‌دهند و صحت فرضیهٔ اصلی پژوهش را تأیید می‌کنند: معماری ترکیبی \lr{QKD+PQC} نه تنها امکان‌پذیر، بلکه برتر از هر رویکرد تک‌لایه‌ای است.
%% — پایان فصل چهارم —
%%%%%%%%%%%%%%%%%%%%%%%%%%%%%%%%%%%%%%%%%%%%%%%%%%%%%%%%%%%%%%%%%%%%%%%%%%%%%
%%  chapter5.tex — فصل پنجم: نتیجه‌گیری و پیشنهادات
%%  paste کنید در main.tex بعد از پایان فصل چهارم، قبل از \end{document}
%%%%%%%%%%%%%%%%%%%%%%%%%%%%%%%%%%%%%%%%%%%%%%%%%%%%%%%%%%%%%%%%%%%%%%%%%%%%%

\chapter{نتیجه‌گیری و پیشنهادات}
\setcounter{figure}{0}
\setcounter{table}{0}
\setcounter{equation}{0}

%% ──────────────────────────────────────────────────────────────────────────
\section{مقدمه}
%% ──────────────────────────────────────────────────────────────────────────

این فصل به‌عنوان حلقه اتصال تمامی تلاش‌های پژوهشی این پایان‌نامه عمل می‌کند و بر آن است تا یافته‌های گردآمده در فصول پیشین را در یک تصویر منسجم و کامل درهم تنیده، به سؤالات اصلی تحقیق پاسخ قطعی دهد و مسیر مطالعات آتی را در این حوزه برای پژوهشگران آینده روشن سازد. این فصل با جمع‌بندی دستاوردهای پژوهش آغاز می‌شود، سپس به پاسخ‌دهی صریح به سؤالات و فرضیه‌های مطرح‌شده در فصل اول می‌پردازد، نوآوری‌های اصلی کار را برجسته می‌سازد، و در نهایت پیشنهادات مشخص برای توسعه و تکمیل این معماری در پژوهش‌های آتی ارائه می‌دهد.

%% ──────────────────────────────────────────────────────────────────────────
\section{نتیجه‌گیری حاصل از پژوهش}
%% ──────────────────────────────────────────────────────────────────────────

پژوهش حاضر از یک پرسش بنیادین آغاز کرد: آیا ارتباطات ماهواره‌های مدار پایین زمین \lr{(Low Earth Orbit)} --- که ستون فقرات زیرساخت‌های اطلاعاتی حیاتی کشورها را تشکیل می‌دهند --- می‌توانند در برابر تهدیدات کامپیوترهای کوانتومی ایمن بمانند؟ در فصل اول و دوم، با تکیه بر مبانی نظری رمزنگاری مدرن، نشان داده شد که الگوریتم شور \lr{(Shor's Algorithm)} قادر است در زمان چندجمله‌ای فاکتورگیری اعداد بزرگ و محاسبه لگاریتم گسسته را انجام دهد و بدین ترتیب زیرساخت کلید عمومی \lr{(Public Key Infrastructure)} موجود که بر \lr{RSA} و \lr{ECDSA} بنا شده، از بنیان متزلزل سازد \{1\}. این تهدید در زمینه ماهواره‌ای شکل هشداردهنده‌تری به خود می‌گیرد؛ چرا که ماهواره‌ها با طول عمر عملیاتی ده تا پانزده سال در معرض خطر استراتژی «همین‌اکنون برداشت کن، دیرتر رمزگشا کن» \lr{(Harvest-Now-Decrypt-Later)} قرار دارند \{22\}. این تهدید با توجه به حساسیت بار اطلاعاتی ماهواره‌های \lr{LEO} --- از داده‌های سنجش از دور و فرمان‌دهی ایمن تا مخابرات امن نظامی --- نه یک خطر فرضی بلکه یک الزام عملی برای تحول فوری در معماری امنیتی این سامانه‌هاست \{1\}.

در فصل سوم، معماری سه‌لایه هیبرید پیشنهادی طراحی شد. لایه اول بر پروتکل توزیع کلید کوانتومی \lr{(QKD)} به روش \lr{Decoy-State BB84} در طول موج ۸۵۰ نانومتر استوار است که از نتایج تجربی ماهواره \lr{Micius} در مدار ۵۰۰ کیلومتری الهام گرفته است؛ آن آزمایشگاه فضایی که نشان داد می‌توان در یک پاس ۲۷۳ ثانیه‌ای بیش از ۳۰۰٬۰۰۰ بیت کلید کوانتومی امن تولید کرد با \lr{QBER} متوسط تنها ۱.۱ درصد \{6\}. لایه دوم از استاندارد \lr{ML-KEM-768 (FIPS 203)} برای کپسوله‌سازی کلید بهره می‌برد که بر مسئله \lr{Module Learning With Errors} ساخته شده و در برابر قوی‌ترین الگوریتم‌های شناخته‌شده کوانتومی با پیچیدگی از مرتبه $2^{c \cdot (N/\beta) \cdot \log \beta}$ مقاوم باقی می‌ماند \{9\}. لایه سوم از \lr{ML-DSA-65 (FIPS 204)} برای احراز هویت دیجیتال استفاده می‌کند که امنیت آن بر مسئله \lr{Module Shortest Integer Solution} استوار است و نقش احراز هویت پساکوانتومی کانال‌های کلاسیک \lr{QKD} را در سطح استاندارد رسمی \lr{NIST} ایفا می‌کند \{10\}. بر فراز این سه لایه، یک ماشین حالت متناهی امنیتی \lr{(Security Finite State Machine)} طراحی شد که تضمین می‌دهد در صورت نزدیک‌شدن \lr{QBER} به آستانه بحرانی ۱۱ درصد (با حاشیه پیشگیرانه ۱۰ درصد)، سیستم به‌طور خودکار به لایه \lr{PQC} تکیه کند.

در فصل چهارم، این معماری از دو جهت مستقل اعتبارسنجی شد. نخست، شبیه‌سازی پایتون مبتنی بر \lr{liboqs} نشان داد که \lr{ML-KEM-768} در تولید کلید با زمانی کمتر از \lr{0.09} میلی‌ثانیه، بیش از ۳۰۰۰ برابر سریع‌تر از \lr{RSA-3072} است \{22\}. استاندارد \lr{FIPS 203} به صراحت اعلام می‌کند که پیاده‌سازی‌های \lr{ML-KEM} هرگز نباید از محاسبات ممیز شناور استفاده کنند، زیرا تمام عملیات بر اساس \lr{NTT} و ضرب‌های ماتریسی در حوزه اعداد صحیح استوار است \{3\}. این ماهیت عدد-صحیح بودن یک پیامد فضایی بسیار مهم دارد: پردازنده‌های مقاوم در برابر تشعشع \lr{(Radiation-Hardened)} که در ماهواره‌ها به کار می‌روند معمولاً فاقد واحد \lr{FPU} با کارایی بالا هستند؛ از این رو \lr{ML-KEM} نه‌فقط از نظر سرعت بلکه از نظر سازگاری معماری سخت‌افزاری نیز انتخاب طبیعی‌تری نسبت به \lr{RSA} است. تأخیر محاسباتی کل لایه \lr{PQC} با نسبت $L_{crypto}/L_{session} < 0.2\%$ نسبت به تأخیر کل برقراری نشست امن (حدود ۳۵ میلی‌ثانیه)، عملاً اثر ناچیزی بر تأخیر کل نشست دارد \{22\}.

دوم، مدل‌سازی تحلیلی بودجه لینک \lr{(Link Budget)} برای سناریوی \lr{LEO} در ارتفاع ۵۰۰ کیلومتر نشان داد که دینامیک متغیر ارتباطات \lr{QKD} --- از جمله \lr{QBER} بالاتر در زوایای \lr{elevation} پایین --- دقیقاً همان الگویی است که طراحی لایه‌بندی‌شده و \lr{FSM} پیشنهادی را توجیه می‌کند. چنانکه چارچوب نظری \lr{Zeng} و همکاران نشان داده است، وجود همزمان یک \lr{primitive} با امنیت اطلاعات-تئوریک \lr{(QKD)} و یک \lr{primitive} با امنیت محاسباتی پساکوانتومی \lr{(ML-KEM)} یک ساختار اعتماد \lr{(Trust Hierarchy)} ایجاد می‌کند که در آن نفوذ به سیستم نهایی مستلزم شکستن هر دو لایه به‌طور همزمان است \{19\}.

نتیجه‌گیری کلان پژوهش این است که معماری هیبرید سه‌لایه پیشنهادی --- شامل \lr{QKD} به‌عنوان لایه فیزیکی کوانتومی، \lr{ML-KEM-768} و \lr{ML-DSA-65} به‌عنوان لایه محاسباتی پساکوانتومی، و \lr{FSM} به‌عنوان لایه هماهنگ‌ساز \lr{Fail-Safe} --- یک استراتژی دفاعی چندلایه برای ارتباطات ماهواره‌ای \lr{LEO} فراهم می‌سازد که از نظر مهندسی موازنه مطلوبی میان امنیت اطلاعات-تئوریک و کارایی محاسباتی برقرار می‌کند. ضعف‌های ذاتی بودجه لینک \lr{QKD} در شرایط دید پایین توسط پایداری لایه \lr{PQC} پوشش داده می‌شود؛ و در مقابل، هر خطر نظری احتمالی درباره مسائل شبکه‌ای که پایه \lr{ML-KEM} را تشکیل می‌دهند، توسط امنیت مستقل از محاسبه لایه \lr{QKD} خنثی می‌گردد. این هم‌افزایی دوسویه همان چیزی است که پژوهش حاضر آن را «معماری ایمن در برابر شکست» می‌نامد \{13, 19, 22\}.
%% ──────────────────────────────────────────────────────────────────────────
\section{پاسخ به سؤالات و فرضیه‌های تحقیق}
%% ──────────────────────────────────────────────────────────────────────────

پژوهش حاضر در فصل اول بر چهار سؤال اصلی و دو فرضیه کلیدی استوار شده بود که یافته‌های فصل چهارم اکنون امکان پاسخ‌گویی صریح به آن‌ها را فراهم می‌سازد.

\textbf{فرضیه رقیب اول:} ماهواره‌های مدار پایین زمین به دلیل محدودیت‌های \lr{SWaP}، توان محاسباتی کافی برای اجرای الگوریتم‌های رمزنگاری پساکوانتومی را ندارند.

یافته‌های شبیه‌سازی این فرضیه را به‌طور قطعی رد می‌کند. \lr{ML-KEM-768} عملیات تولید کلید را در کمتر از \lr{0.07} میلی‌ثانیه انجام می‌دهد، در حالی که \lr{RSA-3072} برای همین عملیات به ۲۴۰ میلی‌ثانیه نیاز دارد --- برتری بیش از ۳۰۰۰ برابری \{22\}. این شکاف عملکردی نه تصادفی بلکه ساختاری است. استاندارد \lr{FIPS 203} به صراحت اعلام می‌کند که پیاده‌سازی‌های \lr{ML-KEM} هیچ‌گاه نباید از محاسبات ممیز شناور استفاده کنند، زیرا تمام عملیات بر اساس \lr{NTT} و ضرب‌های ماتریسی در حوزه اعداد صحیح استوار است \{3\}. ماهیت عدد-صحیح بودن \lr{ML-KEM} یک پیامد فضایی مهم دارد: پردازنده‌های مقاوم در برابر تشعشع معمولاً فاقد \lr{FPU} با کارایی بالا هستند و از این رو \lr{ML-KEM} از نظر سازگاری معماری سخت‌افزاری نیز انتخاب طبیعی‌تری است. نتایج آزمایش‌های پیاده‌سازی بر روی پلتفرم \lr{RISC-V} در پروژه \lr{EU TRISTAN} نشان داد که تأیید امضای \lr{ML-DSA} روی یک پردازنده \lr{160 MHz} در کمتر از ۲۵ میلی‌ثانیه قابل دستیابی است \{15\}.

\textbf{فرضیه رقیب دوم:} اندازه بزرگ‌تر کلیدهای الگوریتم‌های پساکوانتومی سربار ارتباطی غیرقابل قبولی در کانال‌های ماهواره‌ای با پهنای باند محدود ایجاد می‌کند.

تحلیل انجام‌شده این فرضیه را نیز با استدلال کمی رد کرد. مدل زمان انتقال $T_{tx} = S \times 8 / B$ نشان می‌دهد که کل بار ارتباطی یک تبادل کلید \lr{ML-KEM-768} --- شامل کلید عمومی ۱۱۸۴ بایتی و \lr{ciphertext} ۱۰۸۸ بایتی --- در مجموع ۲۲۷۲ بایت است \{9\}. حتی در یک کانال محافظه‌کارانه ۲۵۶ \lr{kbps}، این داده در حدود ۷۱ میلی‌ثانیه منتقل می‌شود که کمتر از ۰.۰۳ درصد از بودجه زمانی یک پاس ۲۹۴ ثانیه‌ای را مصرف می‌کند \{6, 8\}. نسبت کلی‌تر $L_{crypto}/L_{session} < 0.2\%$ نشان می‌دهد که تأخیر رمزنگاری \lr{PQC} در مقیاس تأخیر کل برقراری نشست امن (حدود ۳۵ میلی‌ثانیه) عملاً ناچیز است \{22\}.

\textbf{پاسخ به سؤال اصلی سوم:} چگونه می‌توان هم‌زمان امنیت کانال کلاسیک \lr{QKD} و پویایی بودجه لینک اپتیکی را مدیریت کرد؟

معماری سه‌لایه پیشنهادی یک پاسخ یکپارچه ارائه می‌دهد. مدل‌سازی تحلیلی بودجه لینک نشان داد که اتلاف کانال اپتیکی در طول یک پاس از حدود ۳۱ دسیبل در زنیت تا بیش از ۴۵ دسیبل در زوایای \lr{elevation} پایین تغییر می‌کند \{6, 23\}. در شرایط افت شدید کانال یا تهی‌شدن مخزن کلید، \lr{FSM} به‌طور خودکار به لایه \lr{ML-KEM} تکیه می‌کند. چالش امنیت کانال کلاسیک \lr{QKD} توسط لایه \lr{ML-DSA-65} به‌طور دائمی برطرف شده است. \lr{Wang} و همکاران در آزمایش‌های تجربی خود نشان دادند که امضای دیجیتال پساکوانتومی با زمان اجرای کمتر از یک میلی‌ثانیه، هیچ گلوگاه محاسباتی در کانال کلاسیک \lr{QKD} ایجاد نمی‌کند و تفاوت نرخ کلید \lr{QKD} بین حالت با و بدون احراز هویت \lr{PQC} کمتر از یک انحراف معیار باقی می‌ماند \{10\}. این مجموعه شواهد ثابت می‌کند که معماری پیشنهادی نه‌فقط از نظر نظری صحیح بلکه از نظر عملیاتی نیز به‌طور کامل قابل استقرار است.
%% ──────────────────────────────────────────────────────────────────────────
\section{پیشنهادات برای پژوهش‌های آتی}
%% ──────────────────────────────────────────────────────────────────────────

با وجود دستاوردهای قابل توجه این پژوهش، مرزهای دانش در این حوزه همچنان فرصت‌های وسیعی برای توسعه دارد. بر اساس شکاف‌های شناسایی‌شده در ادبیات و محدودیت‌های ذاتی کار حاضر، چهار مسیر پژوهشی اولویت‌دار برای مطالعات آتی پیشنهاد می‌شود.

\textbf{پیشنهاد اول --- \lr{MDI-QKD} در کنستلاسیون‌های \lr{LEO}:} بررسی و ارزیابی پروتکل \lr{QKD} مستقل از دستگاه اندازه‌گیری \lr{(Measurement-Device-Independent QKD)} در معماری کنستلاسیون‌های چندگانه. محدودیت اساسی پروتکل \lr{Decoy-State BB84} مورد استفاده در پژوهش حاضر آن است که ماهواره باید به‌عنوان یک گره مورد اعتماد \lr{(Trusted Node)} عمل کند. این چالش که در مقاله \lr{Yin} و همکاران (2020) در قالب آزمایش موفق \lr{QKD} مبتنی بر درهم‌تنیدگی در فاصله ۱۱۲۰ کیلومتر برای اولین بار برطرف شد \{7\}، می‌تواند با استفاده از پروتکل \lr{MDI-QKD} در معماری چند-ماهواره‌ای به شکل عملی‌تری تحقق یابد. گزارش ارزیابی \lr{DLR} در سال �2025 نشان داده که \lr{MDI-QKD} حداکثر فاصله زمینی قابل پشتیبانی را به ۲۳۰۴ کیلومتر محدود می‌کند \{23\}؛ پژوهش‌های آتی می‌توانند با بهینه‌سازی هندسه کنستلاسیون \lr{LEO} چندگانه این محدودیت را برطرف سازند.

\textbf{پیشنهاد دوم --- ماژول رمزنگاری \lr{FPGA} فضایی:} طراحی و آزمون یک ماژول رمزنگاری سخت‌افزاری مقاوم در برابر تشعشع \lr{(Radiation-Hardened Cryptographic Module)} برای اجرای بومی \lr{ML-KEM-768} و \lr{ML-DSA-65}. نتایج مقاله \lr{Hadayeghparast} و همکاران نشان داد که پیاده‌سازی \lr{FPGA} الگوریتم‌های مبتنی بر شبکه با تکنیک \lr{temporal redundancy} می‌تواند تمام خطاهای گذرای ناشی از رویدادهای \lr{SEU} را با \lr{overhead} تنها یک درصد در مساحت تشخیص دهد \{16\}؛ همچنین پیاده‌سازی \lr{NTT} بر روی \lr{FPGA} با معماری \lr{R22MDC} می‌تواند \lr{ML-KEM} را با \lr{throughput} ۴۰۸ مگابیت بر ثانیه اجرا کند \{17\}. ترکیب این دو رویکرد روی یک \lr{FPGA} مقاوم در برابر تشعشع از نوع \lr{Xilinx Kintex Ultrascale+} می‌تواند یک نقطه عطف در پیاده‌سازی عملی \lr{PQC} برای ماهواره‌های کوچک \lr{LEO} باشد.

\textbf{پیشنهاد سوم --- بهینه‌سازی پویا با یادگیری ماشین:} استفاده از یادگیری ماشین \lr{(Machine Learning)} برای بهینه‌سازی پویا و پیش‌بینی بودجه لینک اپتیکی \lr{QKD} در شرایط جوی متغیر. مدل اتلاف کانال:
\begin{equation}
L_{total} = L_{diffraction} + L_{pointing} + L_{atmosphere} + L_{absorption}
\label{eq:ch5_loss}
\end{equation}
که در فصول سوم و چهارم به کار رفت، یک مدل ایستا را توصیف می‌کند. یک سیستم پیش‌بینی مبتنی بر شبکه عصبی بازگشتی \lr{(Recurrent Neural Network)} که داده‌های پویای جوی و تاریخچه \lr{QBER} لینک‌های قبلی را ترکیب کند، می‌تواند به ایستگاه زمینی اجازه دهد که پیش از آغاز پاس ماهواره بهترین پنجره زمانی را برای کلیدسازی \lr{QKD} شناسایی کند و آستانه‌های \lr{FSM} را به‌صورت پویا تنظیم نماید.

\textbf{پیشنهاد چهارم --- چابکی رمزنگارانه \lr{(Crypto Agility)}:} ارزیابی سیستماتیک قابلیت جایگزینی الگوریتم‌ها بدون نیاز به بازطراحی کامل سیستم. \lr{FIPS 205} که استاندارد \lr{SLH-DSA} (مشتق از \lr{SPHINCS+}) را تعریف می‌کند، یک الگوریتم امضای دیجیتال کاملاً مبتنی بر توابع \lr{hash} ارائه می‌دهد که امنیت آن به هیچ فرض ریاضی فراتر از امنیت توابع \lr{hash} وابسته نیست \{12\}. از آنجا که تاریخ نشان داده الگوریتم‌هایی مانند \lr{SIKE} در ۲۰۲۲ پس از سال‌ها ارزیابی شکسته شدند \{2\}، معماری سیستم‌های ماهواره‌ای باید از ابتدا با قابلیت \lr{crypto agility} طراحی شود. پژوهش‌های آتی می‌توانند یک لایه \lr{abstraction} رمزنگارانه برای ماهواره‌های \lr{LEO} طراحی کنند که در آن الگوریتم‌های \lr{ML-KEM}، \lr{FN-DSA (FALCON)}، \lr{ML-DSA} و \lr{SLH-DSA} به‌صورت ماژول‌های قابل تعویض پیاده‌سازی شده و ایستگاه زمینی بتواند از طریق یک دستور به‌روزرسانی \lr{firmware} امن، مجموعه الگوریتم فعال ماهواره را در مدار تغییر دهد --- الگویی که پایداری بلندمدت زیرساخت امنیتی ماهواره‌ها را در برابر تهدیدات ناشناخته آینده تضمین می‌کند \{15\}.
%% — پایان فصل پنجم —
%%%%%%%%%%%%%%%%%%%%%%%%%%%%%%%%%%%%%%%%%%%%%%%%%%%%%%%%%%%%%%%%%%%%%%%%%%%%%
%%  abstract_en.tex — چکیده انگلیسی
%%  محل قرارگیری: در main.tex، پس از فهرست منابع و پیش از \end{document}
%%  (یعنی آخرین صفحه محتوایی، قبل از جلد پشتی انگلیسی)
%%%%%%%%%%%%%%%%%%%%%%%%%%%%%%%%%%%%%%%%%%%%%%%%%%%%%%%%%%%%%%%%%%%%%%%%%%%%%

%% ─── چکیده انگلیسی ──────────────────────────────────────────────────────────
%% طبق فرمت IAU تهران مرکزی، چکیده انگلیسی در انتهای پایان‌نامه و پس از
%% فهرست منابع قرار می‌گیرد. این صفحه راست‌به‌چپ نیست؛ باید با
%% دستور \begin{latin}...\end{latin} محیط لاتین فعال شود.

\newpage
\pagestyle{plain}
\begin{latin}
\setlength{\parindent}{0pt}
\setlength{\parskip}{6pt}

\begin{center}

{\Large\textbf{Abstract}}
\end{center}

\vspace{0.5cm}

The rapid expansion of satellite communications in Low Earth Orbit (LEO) and the concurrent advancement of quantum computing pose a fundamental threat to classical cryptographic infrastructures based on RSA and Elliptic Curve Cryptography (ECC). Shor's Algorithm can break these schemes on a sufficiently scaled quantum computer in polynomial time; this reality, compounded by the Harvest Now, Decrypt Later (HNDL) threat, critically endangers the long-term security of satellite data with operational lifetimes of 10 to 15 years. In response to this security gap, this research designs and evaluates a hybrid three-layer security architecture for LEO satellite communications.

The first layer of the proposed architecture employs Decoy-State BB84 Quantum Key Distribution (QKD) at an 850\,nm wavelength to provide information-theoretically secure key establishment. The second layer deploys the ML-KEM-768 standard (FIPS~203) for quantum-resistant key encapsulation over the classical channel. The third layer utilizes ML-DSA-65 (FIPS~204) for post-quantum digital authentication and signing of satellite control messages. A Fail-Safe mechanism, governed by a Security Finite State Machine (FSM), coordinates the first two layers within a resilient framework that automatically transitions the system to PQC-only mode upon degradation of the QKD channel.

Performance evaluation using the liboqs-python library and comparison against classical baselines demonstrated that ML-KEM-768 achieves key generation more than 3{,}000 times faster than RSA-3072, decisively refuting the hypothesis that lattice-based post-quantum algorithms are computationally impractical for satellite hardware. The total session establishment latency of the proposed architecture is approximately 35\,ms, which is negligible relative to the approximately 300-second visibility window of a LEO satellite at 500\,km altitude. The analytical QKD link budget model demonstrated that the Quantum Bit Error Rate (QBER) remains within the 1\%--3\% range at optimal elevation angles, well below the 11\% security threshold of the BB84 protocol. The overall results confirm that the proposed architecture constitutes a trade-off-aware, fail-safe defensive framework against threats from both classical and quantum-era adversaries.

\vspace{0.8cm}
\noindent\textbf{Keywords:} Low Earth Orbit Communications, Post-Quantum Cryptography, Quantum Key Distribution, Quantum Bit Error Rate, Link Budget

\end{latin}
%%%%%%%%%%%%%%%%%%%%%%%%%%%%%%%%%%%%%%%%%%%%%%%%%%%%%%%%%%%%%%%%%%%%%%%%%%%%%
%%  references.tex -- فهرست منابع (ویرایش نهایی)
%%  محل: پس از فصل پنجم، پیش از پیوست
%%%%%%%%%%%%%%%%%%%%%%%%%%%%%%%%%%%%%%%%%%%%%%%%%%%%%%%%%%%%%%%%%%%%%%%%%%%%%

\newpage

\begin{center}
  {\Large\bfseries فهرست منابع}
\end{center}
\addcontentsline{toc}{chapter}{فهرست منابع}
\vspace{0.8cm}

\begin{latin}
\begin{flushleft}
\small
\setlength{\parskip}{8pt}
\setlength{\parindent}{0pt}

\noindent\textbf{\{1\}}~D. J. Bernstein and T. Lange, ``Post-Quantum Cryptography,'' \textit{Nature}, vol. 549, no. 7671, pp. 188--194, Sep. 2017. DOI: 10.1038/nature23461

\vspace{2pt}

\noindent\textbf{\{2\}}~G. Alagic et al., ``Status Report on the Third Round of the NIST Post-Quantum Cryptography Standardization Process,'' \textit{NIST IR 8413}, 2022. DOI: 10.6028/NIST.IR.8413

\vspace{2pt}

\noindent\textbf{\{3\}}~National Institute of Standards and Technology (NIST), ``Module-Lattice-Based Key-Encapsulation Mechanism Standard (ML-KEM),'' \textit{FIPS 203}, 2024. DOI: 10.6028/NIST.FIPS.203

\vspace{2pt}

\noindent\textbf{\{4\}}~S. Pirandola, U. L. Andersen, L. Banchi et al., ``Advances in Quantum Cryptography,'' \textit{Advances in Optics and Photonics}, vol. 12, no. 4, pp. 1012--1236, 2020. DOI: 10.1364/AOP.361502

\vspace{2pt}

\noindent\textbf{\{5\}}~R. Bedington, J. M. Arrazola, and A. Ling, ``Progress in Satellite Quantum Key Distribution,'' \textit{npj Quantum Information}, vol. 3, article 30, 2017. DOI: 10.1038/s41534-017-0031-5

\vspace{2pt}

\noindent\textbf{\{6\}}~S.-K. Liao, W.-Q. Cai, W.-Y. Liu et al., ``Satellite-to-Ground Quantum Key Distribution,'' \textit{Nature}, vol. 549, pp. 43--47, 2017. DOI: 10.1038/nature23655

\vspace{2pt}

\noindent\textbf{\{7\}}~J. Yin, Y.-H. Li, S.-K. Liao et al., ``Entanglement-Based Secure Quantum Cryptography over 1,120 Kilometres,'' \textit{Nature}, vol. 582, pp. 501--505, Jun. 2020. DOI: 10.1038/s41586-020-2401-y

\vspace{2pt}

\noindent\textbf{\{8\}}~T. Roger, R. Singh, C. Perumangatt et al., ``Real-Time Gigahertz Free-Space Quantum Key Distribution Within an Emulated Satellite Overpass,'' \textit{Science Advances}, vol. 9, eadj5873, Dec. 2023. DOI: 10.1126/sciadv.adj5873

\vspace{2pt}

\noindent\textbf{\{9\}}~National Institute of Standards and Technology (NIST), ``Module-Lattice-Based Digital Signature Standard (ML-DSA),'' \textit{FIPS 204}, 2024. DOI: 10.6028/NIST.FIPS.204

\vspace{2pt}

\noindent\textbf{\{10\}}~X. Gong, X. Chen, Y. Li et al., ``Experimental Authentication of Quantum Key Distribution with Post-Quantum Cryptography,'' \textit{npj Quantum Information}, vol. 7, article 32, 2021. DOI: 10.1038/s41534-021-00400-7

\vspace{2pt}

\noindent\textbf{\{11\}}~S.-K. Liao, W.-Q. Cai, J. Handsteiner et al., ``Satellite-Relayed Intercontinental Quantum Network,'' \textit{Physical Review Letters}, vol. 120, p. 030501, 2018. DOI: 10.1103/PhysRevLett.120.030501

\vspace{2pt}

\noindent\textbf{\{12\}}~National Institute of Standards and Technology (NIST), ``Stateless Hash-Based Digital Signature Standard (SLH-DSA),'' \textit{FIPS 205}, 2024. DOI: 10.6028/NIST.FIPS.205

\vspace{2pt}

\noindent\textbf{\{13\}}~A. Rani, X. Ai, A. Gupta, R. S. Adhikari, and R. Malaney, ``Combined Quantum and Post-Quantum Security for Earth-Satellite Channels,'' \textit{IEEE QCNC 2025}. DOI: 10.48550/arXiv.2502.14240

\vspace{2pt}

\noindent\textbf{\{14\}}~I. B. Djordjevic, ``Physical-Layer Security, Quantum Key Distribution, and Post-Quantum Cryptography,'' \textit{Entropy (MDPI)}, vol. 24, no. 7, p. 935, Jul. 2022. DOI: 10.3390/e24070935

\vspace{2pt}

\noindent\textbf{\{15\}}~S. Xu, A. V. Alonso, J. L\'egar\'e, and P. Eisen, ``A Preliminary Study of PQC Implementations for Satellite Communication Networks,'' \textit{Proc. IEEE Aerospace Conference (AEROCONF)}, 2023. DOI: 10.1109/AERO55745.2023.10115923

\vspace{2pt}

\noindent\textbf{\{16\}}~S. Hadayeghparast, S. Bayat-Sarmadi, and S. Ebrahimi, ``High-Speed Post-Quantum Cryptoprocessor Based on RISC-V Architecture for IoT,'' \textit{IEEE Internet of Things Journal}, vol. 9, no. 19, pp. 15839--15846, 2022. DOI: 10.1109/JIOT.2022.3147024

\vspace{2pt}

\noindent\textbf{\{17\}}~D. Kundi, Y. Zhang, C. Wang, A. Khalid, M. O'Neill, and W. Liu, ``Ultra High-Speed Polynomial Multiplications for Lattice-Based Cryptography on FPGAs,'' \textit{IEEE Transactions on Emerging Topics in Computing}, 2022. DOI: 10.1109/TETC.2022.3144101

\vspace{2pt}

\noindent\textbf{\{18\}}~D. Dharminder et al., ``A Post-Quantum Secure Construction of an Authentication Protocol for Satellite Communication,'' \textit{International Journal of Satellite Communications and Networking}, vol. 41, 2023. DOI: 10.1002/sat.1455

\vspace{2pt}

\noindent\textbf{\{19\}}~J. S. Liu, X. Ai, A. Gupta, and L. Jiang, ``Practical Hybrid PQC-QKD Protocols with Enhanced Security and Performance,'' \textit{arXiv:2411.01086}, 2024. DOI: 10.48550/arXiv.2411.01086

\vspace{2pt}

\noindent\textbf{\{20\}}~H. Takenaka, A. Carrasco-Casado, M. Fujiwara, M. Kitamura, M. Sasaki, and M. Toyoshima, ``Satellite-to-Ground Quantum Communication Using a 50-kg-Class Micro-Satellite,'' \textit{Nature Photonics}, vol. 11, no. 8, pp. 502--508, 2017. DOI: 10.1038/nphoton.2017.107

\vspace{2pt}

\noindent\textbf{\{21\}}~A. Jokiniemi, ``NIST's Module-Lattice-Based Key-Encapsulation Mechanism and Its Security,'' Master's thesis, Department of Physics, University of Helsinki, 2025. DOI: 10.10123/helsinki.academic.2025.aj

\vspace{2pt}

\noindent\textbf{\{22\}}~T. Ghosh and I. Nath, ``Secure Satellite Communication in the Post-Quantum Era: A Lattice-Based Cryptographic Approach,'' \textit{International Journal of Satellite Communications and Networking}, 2026. DOI: 10.1002/sat.70041

\vspace{2pt}

\noindent\textbf{\{23\}}~D. Orsucci, P. Kleinpa{\ss}, J. Meister, I. De Marco, S. H\"ausler, T. Strang, N. Walenta, and F. Moll, ``Assessment of Practical Satellite Quantum Key Distribution Architectures for Current and Near-Future Missions,'' \textit{International Journal of Satellite Communications and Networking}, 2025. DOI: 10.1002/sat.1544

\vspace{2pt}

\noindent\textbf{\{24\}}~T. Liu, G. Ramachandran, and R. Jurdak, ``Post-Quantum Cryptography for Internet of Things: A Survey on Performance and Optimization,'' \textit{arXiv:2401.17538}, 2023. DOI: 10.48550/arXiv.2401.17538

\vspace{2pt}

\end{flushleft}
\end{latin}
%%  appendix.tex -- پیوست: کدهای شبیه‌سازی پایتون
%%  محل قرارگیری: پس از فهرست منابع و پیش از چکیده انگلیسی
%%%%%%%%%%%%%%%%%%%%%%%%%%%%%%%%%%%%%%%%%%%%%%%%%%%%%%%%%%%%%%%%%%%%%%%%%%%%%

\newpage
\appendix
\addcontentsline{toc}{chapter}{پیوست: کدهای شبیه‌سازی پایتون}

\begin{center}
{\Large\bfseries پیوست}\\[0.4cm]
{\large\bfseries کدهای شبیه‌سازی پایتون}
\end{center}

\vspace{0.5cm}

%% -------------------------------------------------------------------------
\section*{پیوست الف بنچمارک الگوریتم‌های \lr{PQC} با \lr{liboqs-python}}
\addcontentsline{toc}{section}{پیوست الف -- بنچمارک PQC}
%% -------------------------------------------------------------------------

این اسکریپت عملکرد \lr{ML-KEM} و \lr{ML-DSA} را با معادل‌های کلاسیک \lr{RSA} و \lr{ECDH/ECDSA} مقایسه می‌کند.
خروجی شامل زمان اجرای عملیات تولید کلید، کپسوله‌سازی، امضا و تأیید، به‌همراه اندازه کلیدها است.

\vspace{0.3cm}

\begin{latin}
\small
\begin{verbatim}
# ============================================================
# Appendix A -- PQC Benchmark: ML-KEM and ML-DSA vs Classical
# Author: Melika MoghadamKia
# IAU Central Tehran Branch -- Winter 1404
# ============================================================

import oqs
import time
import statistics
from cryptography.hazmat.primitives.asymmetric import rsa, ec, padding
from cryptography.hazmat.primitives import hashes
from cryptography.hazmat.backends import default_backend
from cryptography.hazmat.primitives.asymmetric.ec import ECDH

NUM_ITERATIONS = 1000  # number of repetitions per measurement


def benchmark_kem(algorithm_name: str) -> dict:
    kem = oqs.KeyEncapsulation(algorithm_name)
    keygen_times, encap_times, decap_times = [], [], []
    for _ in range(NUM_ITERATIONS):
        t0 = time.perf_counter()
        public_key = kem.generate_keypair()
        keygen_times.append((time.perf_counter() - t0) * 1000)
        t0 = time.perf_counter()
        ciphertext, _ = kem.encap_secret(public_key)
        encap_times.append((time.perf_counter() - t0) * 1000)
        t0 = time.perf_counter()
        kem.decap_secret(ciphertext)
        decap_times.append((time.perf_counter() - t0) * 1000)
    return {
        'algorithm': algorithm_name,
        'pk_size':   len(public_key),
        'ct_size':   len(ciphertext),
        'sk_size':   len(kem.export_secret_key()),
        'keygen_ms': statistics.median(keygen_times),
        'encap_ms':  statistics.median(encap_times),
        'decap_ms':  statistics.median(decap_times),
    }


def benchmark_sig(algorithm_name: str) -> dict:
    signer = oqs.Signature(algorithm_name)
    msg = b'Satellite control command: ADJUST_ORBIT_V2'
    keygen_times, sign_times, verify_times = [], [], []
    for _ in range(NUM_ITERATIONS):
        t0 = time.perf_counter()
        public_key = signer.generate_keypair()
        keygen_times.append((time.perf_counter() - t0) * 1000)
        t0 = time.perf_counter()
        signature = signer.sign(msg)
        sign_times.append((time.perf_counter() - t0) * 1000)
        t0 = time.perf_counter()
        signer.verify(msg, signature, public_key)
        verify_times.append((time.perf_counter() - t0) * 1000)
    return {
        'algorithm': algorithm_name,
        'pk_size':   len(public_key),
        'sig_size':  len(signature),
        'sk_size':   len(signer.export_secret_key()),
        'keygen_ms': statistics.median(keygen_times),
        'sign_ms':   statistics.median(sign_times),
        'verify_ms': statistics.median(verify_times),
    }


def benchmark_rsa(key_bits: int) -> dict:
    keygen_times = []
    for _ in range(100):
        t0 = time.perf_counter()
        priv = rsa.generate_private_key(
            public_exponent=65537, key_size=key_bits,
            backend=default_backend()
        )
        keygen_times.append((time.perf_counter() - t0) * 1000)
    priv = rsa.generate_private_key(
        public_exponent=65537, key_size=key_bits,
        backend=default_backend()
    )
    pub = priv.public_key()
    msg = b'test_rsa'
    oaep = padding.OAEP(
        mgf=padding.MGF1(hashes.SHA256()),
        algorithm=hashes.SHA256(), label=None
    )
    enc_times, dec_times = [], []
    for _ in range(1000):
        t0 = time.perf_counter()
        ct = pub.encrypt(msg, oaep)
        enc_times.append((time.perf_counter() - t0) * 1000)
        t0 = time.perf_counter()
        priv.decrypt(ct, oaep)
        dec_times.append((time.perf_counter() - t0) * 1000)
    return {
        'algorithm': f'RSA-{key_bits}',
        'pk_size':   key_bits // 8,
        'ct_size':   key_bits // 8,
        'keygen_ms': statistics.median(keygen_times),
        'encap_ms':  statistics.median(enc_times),
        'decap_ms':  statistics.median(dec_times),
    }


def benchmark_ecdh(curve_name: str) -> dict:
    curve = ec.SECP256R1() if curve_name == 'P-256' else ec.SECP384R1()
    keygen_times, dh_times = [], []
    for _ in range(NUM_ITERATIONS):
        t0 = time.perf_counter()
        priv = ec.generate_private_key(curve, default_backend())
        keygen_times.append((time.perf_counter() - t0) * 1000)
        peer = ec.generate_private_key(curve, default_backend())
        t0 = time.perf_counter()
        priv.exchange(ECDH(), peer.public_key())
        dh_times.append((time.perf_counter() - t0) * 1000)
    pk_size = 32 if curve_name == 'P-256' else 48
    return {
        'algorithm': f'ECDH {curve_name}',
        'pk_size':   pk_size,
        'ct_size':   pk_size,
        'keygen_ms': statistics.median(keygen_times),
        'encap_ms':  statistics.median(dh_times),
        'decap_ms':  statistics.median(dh_times),
    }


def benchmark_ecdsa(curve_name: str) -> dict:
    curve = ec.SECP256R1() if curve_name == 'P-256' else ec.SECP384R1()
    msg = b'Satellite control command: ADJUST_ORBIT_V2'
    coord = 32 if curve_name == 'P-256' else 48
    keygen_times, sign_times, verify_times = [], [], []
    for _ in range(NUM_ITERATIONS):
        t0 = time.perf_counter()
        priv = ec.generate_private_key(curve, default_backend())
        keygen_times.append((time.perf_counter() - t0) * 1000)
        pub = priv.public_key()
        t0 = time.perf_counter()
        signature = priv.sign(msg, ec.ECDSA(hashes.SHA256()))
        sign_times.append((time.perf_counter() - t0) * 1000)
        t0 = time.perf_counter()
        pub.verify(signature, msg, ec.ECDSA(hashes.SHA256()))
        verify_times.append((time.perf_counter() - t0) * 1000)
    # Report standard raw sizes: uncompressed point x||y and raw r||s signature.
    # (The library returns a DER-encoded signature; the raw form is used for a
    #  fair comparison with the PQC schemes and to match Table 4-3.)
    return {
        'algorithm': f'ECDSA {curve_name}',
        'pk_size':   2 * coord,   # 64 (P-256) / 96 (P-384)
        'sig_size':  2 * coord,   # 64 (P-256) / 96 (P-384)
        'sk_size':   coord,       # 32 (P-256) / 48 (P-384)
        'keygen_ms': statistics.median(keygen_times),
        'sign_ms':   statistics.median(sign_times),
        'verify_ms': statistics.median(verify_times),
    }


if __name__ == '__main__':
    print('=' * 70)
    print('  PQC Benchmark: ML-KEM / ML-DSA vs Classical Algorithms')
    print('=' * 70)

    kem_algs = ['ML-KEM-512', 'ML-KEM-768', 'ML-KEM-1024']
    results = [benchmark_kem(a) for a in kem_algs]
    results += [benchmark_rsa(2048), benchmark_rsa(3072)]
    results += [benchmark_ecdh('P-256'), benchmark_ecdh('P-384')]

    hdr = (f"{'Algorithm':<16} {'PK(B)':>7} {'CT(B)':>7}"
           f" {'KeyGen(ms)':>11} {'Encap(ms)':>10} {'Decap(ms)':>10}")
    print(hdr)
    print('-' * 65)
    for r in results:
        print(f"{r['algorithm']:<16} {r['pk_size']:>7} {r['ct_size']:>7}"
              f" {r['keygen_ms']:>11.4f} {r['encap_ms']:>10.4f}"
              f" {r['decap_ms']:>10.4f}")

    sig_algs = ['ML-DSA-44', 'ML-DSA-65', 'ML-DSA-87']
    sig_res = [benchmark_sig(a) for a in sig_algs]
    sig_res += [benchmark_ecdsa('P-256'), benchmark_ecdsa('P-384')]
    shdr = (f"{'Algorithm':<16} {'PK(B)':>7} {'Sig(B)':>7}"
            f" {'KeyGen(ms)':>11} {'Sign(ms)':>9} {'Verify(ms)':>11}")
    print('\n' + shdr)
    print('-' * 65)
    for r in sig_res:
        print(f"{r['algorithm']:<16} {r['pk_size']:>7} {r['sig_size']:>7}"
              f" {r['keygen_ms']:>11.4f} {r['sign_ms']:>9.4f}"
              f" {r['verify_ms']:>11.4f}")
    print('\n=== Benchmark complete ===')
\end{verbatim}
\end{latin}
%%% -------------------------------------------------------------------------
\section*{پیوست ب مدل بودجه لینک \lr{QKD} و ماشین حالت امنیتی (\lr{FSM})}
\addcontentsline{toc}{section}{پیوست ب -- مدل بودجه لینک QKD و FSM}
%% -------------------------------------------------------------------------

این اسکریپت مدل تحلیلی بودجه لینک کانال \lr{QKD} ماهواره‌ای را برای یک ماهواره
\lr{LEO} در ارتفاع ۵۰۰ کیلومتر شبیه‌سازی می‌کند. برای هر زاویه فراز،
اتلاف کل کانال، \lr{QBER} تخمینی (با در نظر گرفتن اثرات نویز پس‌زمینه و دریچه زمانی) و وضعیت ماشین حالت امنیتی با منطق شمارنده پایدار محاسبه می‌شود.

\vspace{0.3cm}

\begin{latin}
\small
\begin{verbatim}
# ============================================================
# Appendix B -- QKD Link Budget Model + Security FSM
# Author: Melika MoghadamKia
# IAU Central Tehran Branch -- Winter 1404
# References:
#   [5]  R. Bedington et al., npj Quantum Information, 2017.
#   [8]  T. Roger et al., Science Advances, 2023.
#   [23] D. Orsucci et al., Satell. Commun. Netw., 2025.
# ============================================================

import numpy as np
import math
from enum import Enum, auto

# Physical constants and system parameters
EARTH_RADIUS_M   = 6.371e6
ORBIT_ALTITUDE_M = 500e3       # LEO altitude: 500 km
WAVELENGTH_M     = 850e-9      # QKD laser wavelength: 850 nm
D_TX_M           = 0.10        # transmitter telescope aperture diameter
D_RX_M           = 0.60        # receiver telescope diameter

# Gaussian beam waist (1/e^2 radius); Optimized waist to avoid truncation 
# at the transmitter aperture edges.
W0_M             = D_TX_M / 2.2 
RX_OBSTRUCTION   = 0.30        # central obstruction ratio (Cassegrain)
ETA_TURB         = 0.631       # turbulence efficiency (2 dB)
ETA_RX_OPTICS    = 0.753       # receiver optics efficiency
ETA_DETECTOR     = 0.580       # Si-APD detector efficiency
ETA_COUPLING     = 0.560       # fiber coupling efficiency
ETA_PAT_NOMINAL  = 0.912       # PAT system efficiency at zenith
QBER_BASELINE    = 0.011       # intrinsic QBER: 1.1%
DARK_COUNT_RATE  = 130.0       # dark count rate (count/s, 4 APDs)
T_GATE           = 1.0e-9      # 1 ns time gating window
P_BG             = 1.0e-6      # Background noise probability per pulse
PULSE_RATE_HZ    = 1.0e9       # laser repetition rate: 1 GHz
ZENITH_TRANSMIT  = 0.50        # zenith atmospheric transmissivity
MU_SIGNAL        = 0.69        # mean photons per pulse (Weak Coherent Pulse)

# FSM thresholds (QBER axis)
QBER_TH_WARN     = 0.08        # warning threshold: 8%
QBER_TH_FALLBACK = 0.10        # preventive fallback threshold: 10%
QBER_TH_RECOVER  = 0.068       # recovery threshold: 6.8% (15% hysteresis vs warn)
# NOTE: the BB84 information-theoretic security limit is 11%; the FSM fallback
#       at 10% is a deliberate safety margin below that hard limit.

# FSM thresholds (key-pool axis, fraction of full pool)
POOL_TH_H2        = 0.30       # below 30% -> pool protection (H2)
POOL_TH_H3        = 0.05       # below 5%  -> emergency (H3)
PERSISTENCE_DEPTH = 3          # consecutive samples for any QBER transition


class SecurityState(Enum):
    H1_FULL_SECURITY = auto()  # QKD + PQC, pool > 30%, QBER < 8%
    H2_POOL_PROTECT  = auto()  # both layers, QKD pool in 5-30% band
    H3_EMERGENCY     = auto()  # PQC-only (QKD down / pool < 5% / high QBER)
    H4_RECOVERY      = auto()  # QKD re-establishing, QBER stable below 6.8%
    H5_SAFE_SHUTDOWN = auto()  # critical crypto fault; latching state


class SecurityFSM:
    """Five-state fail-safe FSM (H1-H5) on two axes: QBER and key-pool level.
    Symmetric persistence counters (Up/Down Counter Logic) protect every 
    QBER-based transition against chattering from transient atmospheric noise."""

    def __init__(self):
        self.state = SecurityState.H1_FULL_SECURITY
        self.warn_count = 0      # up/down counter for QBER > warn
        self.recover_count = 0   # up/down counter for QBER < recover

    def transition(self, qber: float, pool_frac: float = 1.0,
                   qkd_available: bool = True,
                   crypto_fault: bool = False) -> SecurityState:
        # H5 has top priority and is latching (exit only via ground_reset).
        if crypto_fault:
            self.state = SecurityState.H5_SAFE_SHUTDOWN
            return self.state
        if self.state == SecurityState.H5_SAFE_SHUTDOWN:
            return self.state

        # Up/Down Counter logic to prevent chattering (Leaky Bucket)
        if qber > QBER_TH_WARN:
            self.warn_count = min(self.warn_count + 1, PERSISTENCE_DEPTH + 2)
        else:
            self.warn_count = max(self.warn_count - 1, 0)
            
        if qber < QBER_TH_RECOVER:
            self.recover_count = min(self.recover_count + 1, PERSISTENCE_DEPTH + 2)
        else:
            self.recover_count = max(self.recover_count - 1, 0)

        # Sustained high QBER -> emergency, after 3 warns AND crossing fallback.
        qber_fallback = (self.warn_count >= PERSISTENCE_DEPTH and
                         qber > QBER_TH_FALLBACK)

        if self.state in (SecurityState.H1_FULL_SECURITY,
                          SecurityState.H2_POOL_PROTECT):
            if qber_fallback or (not qkd_available) or pool_frac < POOL_TH_H3:
                self.state = SecurityState.H3_EMERGENCY
            elif self.state == SecurityState.H1_FULL_SECURITY:
                if pool_frac < POOL_TH_H2:                  # H1 -> H2
                    self.state = SecurityState.H2_POOL_PROTECT
            elif self.state == SecurityState.H2_POOL_PROTECT:
                if pool_frac >= POOL_TH_H2:                 # H2 -> H1 (pool hysteresis)
                    self.state = SecurityState.H1_FULL_SECURITY

        elif self.state == SecurityState.H3_EMERGENCY:
            if qkd_available and self.recover_count >= PERSISTENCE_DEPTH:
                self.state = SecurityState.H4_RECOVERY     # H3 -> H4

        elif self.state == SecurityState.H4_RECOVERY:
            if qber_fallback or (not qkd_available):
                self.state = SecurityState.H3_EMERGENCY    # relapse
            elif pool_frac >= POOL_TH_H2:
                self.state = SecurityState.H1_FULL_SECURITY  # H4 -> H1
            elif pool_frac >= POOL_TH_H3:
                self.state = SecurityState.H2_POOL_PROTECT   # H4 -> H2

        return self.state

    def ground_reset(self) -> SecurityState:
        # ML-DSA-authenticated command from the ground control station.
        self.state = SecurityState.H1_FULL_SECURITY
        self.warn_count = 0
        self.recover_count = 0
        return self.state

    def active_layers(self) -> str:
        return {
            SecurityState.H1_FULL_SECURITY: 'QKD + ML-KEM-768 + ML-DSA-65',
            SecurityState.H2_POOL_PROTECT:  'QKD (pool-saving) + ML-KEM-768 + ML-DSA-65',
            SecurityState.H3_EMERGENCY:     'ML-KEM-768 + ML-DSA-65 (PQC-only)',
            SecurityState.H4_RECOVERY:      'QKD (re-establishing) + ML-KEM-768',
            SecurityState.H5_SAFE_SHUTDOWN: '--- SAFE SHUTDOWN: all crypto blocked ---',
        }.get(self.state, 'unknown')


def link_distance(elev_deg: float) -> float:
    # Link distance via spherical geometry
    a = math.radians(elev_deg)
    R, h = EARTH_RADIUS_M, ORBIT_ALTITUDE_M
    return -R*math.sin(a) + math.sqrt(R**2*math.sin(a)**2 + h**2 + 2*h*R)


def diffraction_eta(L: float) -> float:
    # Gaussian-beam geometric (diffraction) efficiency.
    # Half-angle divergence of a Gaussian beam: theta = lambda / (pi * w0).
    theta_half = WAVELENGTH_M / (math.pi * W0_M)
    w_L = math.sqrt(W0_M**2 + (theta_half * L)**2)   # 1/e^2 beam radius at range L
    a = D_RX_M / 2.0                                 # receiver aperture radius
    r_in = RX_OBSTRUCTION * a                        # blocked central radius
    # Gaussian power captured in the receiver annulus (aperture minus obstruction):
    captured = (math.exp(-2.0 * r_in**2 / w_L**2) -
                math.exp(-2.0 * a**2 / w_L**2))
    return min(captured, 1.0)


def atm_eta(elev_deg: float) -> float:
    # Air mass (cosecant) model for atmospheric transmissivity; valid for
    # elevation >= ~20 deg (plane-parallel approximation).
    a = math.radians(elev_deg)
    return ZENITH_TRANSMIT ** (1.0 / max(math.sin(a), 1e-6))


def pat_eta(elev_deg: float) -> float:
    return max(ETA_PAT_NOMINAL - 0.162*(1.0 - elev_deg/90.0), 0.70)


def total_eta(elev_deg: float) -> float:
    L = link_distance(elev_deg)
    return (diffraction_eta(L) * atm_eta(elev_deg) *
            ETA_TURB * pat_eta(elev_deg) *
            ETA_RX_OPTICS * ETA_DETECTOR * ETA_COUPLING)


def total_loss_db(elev_deg: float) -> float:
    eta = total_eta(elev_deg)
    return -10.0 * math.log10(eta) if eta > 0 else float('inf')


def estimate_qber(elev_deg: float) -> float:
    eta = total_eta(elev_deg)
    if eta <= 0:
        return 0.50
    signal = MU_SIGNAL * PULSE_RATE_HZ * eta
    
    # Calculate noise probability per pulse including Dark Counts and Background Light
    noise_prob = (DARK_COUNT_RATE * T_GATE) + P_BG
    
    # QBER with realistic noise modeling
    qber = QBER_BASELINE + (0.5 * noise_prob * PULSE_RATE_HZ) / max(signal, 1e-9)
    if elev_deg < 30.0:
        qber += 0.005 * (30.0 - elev_deg) / 10.0
    return min(qber, 0.50)


if __name__ == '__main__':
    print('=' * 90)
    print(' QKD Link Budget + Security FSM -- LEO satellite at 500 km')
    print('=' * 90)
    hdr = (f"{'Elev(deg)':>9} {'Dist(km)':>10} {'Loss(dB)':>10}"
           f" {'eta':>11} {'QBER(%)':>9} {'State':<6} Layers")
    print(hdr)
    print('-' * 90)

    fsm = SecurityFSM()
    transitions = []

    for elev in np.arange(20, 91, 5):
        L_km  = link_distance(elev) / 1000.0
        loss  = total_loss_db(elev)
        eta   = total_eta(elev)
        qber  = estimate_qber(elev)
        prev  = fsm.state
        state = fsm.transition(qber)          # clear-sky pass: pool full, no fault
        if state != prev:
            transitions.append((elev, prev.name, state.name, qber*100))
        label = state.name.split('_')[0]       # 'H1'..'H5'
        print(f"{elev:>9.0f} {L_km:>10.1f} {loss:>10.2f}"
              f" {eta:>11.2e} {qber*100:>9.3f}"
              f" {label:<6} {fsm.active_layers()}")

    print('=' * 90)
    print('\nFSM State Transitions (elevation sweep):')
    if transitions:
        for elev, f, t, q in transitions:
            print(f'  {elev}deg: {f} -> {t}  (QBER={q:.3f}%)')
    else:
        print('  No transitions -- channel healthy throughout pass.')

    # --- Demonstration of the five-state FSM under synthetic conditions ---
    print('\n' + '=' * 90)
    print(' Security FSM demonstration (synthetic QBER / key-pool / fault sequence)')
    print('=' * 90)
    # each step: (QBER, key-pool fraction, QKD available, crypto fault)
    scenario = [
        (0.012, 1.00, True,  False),   # healthy
        (0.012, 0.25, True,  False),   # pool drops  -> H2
        (0.090, 0.25, True,  False),   # QBER > warn (1/3)
        (0.105, 0.25, True,  False),   # QBER > warn (2/3)
        (0.115, 0.25, True,  False),   # 3rd sample + > fallback -> H3
        (0.060, 0.25, True,  False),   # recovering (1/3)
        (0.055, 0.40, True,  False),   # recovering (2/3)
        (0.050, 0.60, True,  False),   # 3rd low sample -> H4
        (0.020, 0.80, True,  False),   # stable + pool full -> H1
        (0.020, 0.80, True,  True),    # crypto fault -> H5 (latching)
        (0.020, 0.80, True,  False),   # stays H5 until ground reset
    ]
    fsm2 = SecurityFSM()
    print(f"{'Step':>4} {'QBER%':>7} {'Pool':>6} {'QKD':>5} {'Fault':>6}"
          f"  {'State':<18} Layers")
    print('-' * 90)
    for i, (q, pool, avail, fault) in enumerate(scenario, 1):
        st = fsm2.transition(q, pool, avail, fault)
        print(f"{i:>4} {q*100:>7.1f} {pool:>6.2f} {str(avail):>5}"
              f" {str(fault):>6}  {st.name:<18} {fsm2.active_layers()}")
    fsm2.ground_reset()
    print(f"{'rst':>4} {'--':>7} {'--':>6} {'--':>5} {'--':>6}"
          f"  {fsm2.state.name:<18} {fsm2.active_layers()}"
          f"  (after ML-DSA ground reset)")

    print('\n=== Simulation complete ===')
\end{verbatim}
\end{latin}
%% -- پایان پیوست --
\end{document}
